\label{chap:83-01}

% RESULTADO_RECUPERADO. Owner integro: sections/hmt/03_app_ortograma.tex
% Destino causal: capitulo 1; la jerarquia se subordina sin alterar el owner.
% Los dos bloques científicos del owner se activan en una única residencia.
\begingroup
\def\HMTAPPFormal{1}
\def\HMTAPPDevelopments{1}
\begin{HMTScientificOwnerSubordinated}
\input{sections/hmt/03_app_ortograma}
\end{HMTScientificOwnerSubordinated}
\endgroup

% Unidad U003. Redacción autónoma; tablas y censos reconstruidos desde 1063/live.

\section{Matrice arithmétique de l’APP et structure de l’anneau local}
\label{p12-83-c1-s1:chap:app-hojas-algebra-local}

La \cref{p12-83-c1-s2:chap:app-completacion-gnomonica} a construit un unique support
\(X_9\). L’APP possède une unique matrice arithmétique principale : la multiplication
nonadique publiée par racine numérique positive. La somme ne constitue pas une
deuxième table APP ; elle est la loi interne d’accumulation qui engendre les lignes et
les colonnes de la matrice multiplicative et fournit une observable auxiliaire de
comparaison. Cette distinction sera maintenue tout au long de l’ouvrage.

\subsection{Matrice multiplicative nonadique et racine numérique}

Pour \(a,b\in\mathcal D_9=\{1,\ldots,9\}\), l’entrée de l’APP s’obtient
directement par
\begin{equation}
 \boxed{\Pi(a,b)=\operatorname{dr}_9(ab)
 =1+((ab-1)\bmod9).}
 \label{p12-83-c1-s1:eq:app-matriz-multiplicativa-directa}
\end{equation}
L’unique matrice APP est donc
\[
\begin{array}{c|ccccccccc}
\Pi&1&2&3&4&5&6&7&8&9\\\hline
1&1&2&3&4&5&6&7&8&9\\
2&2&4&6&8&1&3&5&7&9\\
3&3&6&9&3&6&9&3&6&9\\
4&4&8&3&7&2&6&1&5&9\\
5&5&1&6&2&7&3&8&4&9\\
6&6&3&9&6&3&9&6&3&9\\
7&7&5&3&1&8&6&4&2&9\\
8&8&7&6&5&4&3&2&1&9\\
9&9&9&9&9&9&9&9&9&9
\end{array}
\]
Chaque ligne peut également être reconstruite par addition répétée de son indice,
mais cette récurrence décrit la loi interne de la matrice et non un deuxième
feuillet. Les lignes et colonnes d’indices \(1,2,4,5,7,8\) parcourent les neuf
classes ; celles d’indices \(3,6,9\) sont exactement celles qui contractent le support.

Dans la carte résiduelle, l’ensemble des unités est
\[
 R_9^\times=(\mathbb Z/9\mathbb Z)^\times
   =\{[1],[2],[4],[5],[7],[8]\}.
\]
L’idéal maximal est
\[
 \mathfrak m=(3)=\{[0],[3],[6]\}.
\]

\begin{theorem}[Classification des lignes multiplicatives]
\label{p12-83-c1-s1:thm:app-clasificacion-filas-producto}
Pour \(a\in\mathcal D_9\), l'application
\[
 m_a:\mathcal D_9\longrightarrow\mathcal D_9,
 \qquad b\longmapsto\Pi(a,b)
\]
est bijective si et seulement si \([a]_9\in R_9^\times\). Si
\([a]_9\in\mathfrak m\setminus\{0\}\), son image a trois éléments. Si
\([a]_9=[0]_9\), son image a un seul élément.
\end{theorem}

\begin{proof}
La réduction résiduelle de \(m_a\) est la multiplication par \([a]_9\). Dans un
anneau fini, cette multiplication est bijective exactement lorsque \([a]_9\)
est une unité. Si \([a]_9=[3]\) ou \([6]\), son image est l’idéal
\(\mathfrak m\), de cardinal trois. Si \([a]_9=[0]\), toute entrée est envoyée
sur la classe nulle, publiée comme \(9\).
\end{proof}

\begin{proposition}[Latinité de la loi additive interne]
\label{p12-83-c1-s1:prop:app-latinidad-aditiva}
L’observable auxiliaire
\[
 \Sigma(a,b)=\operatorname{dr}_9(a+b)
\]
agit par translations : pour chaque \(a,c\in\mathcal D_9\), il existe un unique
\(b\in\mathcal D_9\) tel que \(\Sigma(a,b)=c\), et de même pour l’autre
coordonnée.
\end{proposition}

\begin{proof}
Dans la carte résiduelle, \([a+b]_9=[c]_9\) a l’unique solution
\([b]_9=[c-a]_9\). La section positive transporte cette classe vers un unique
élément de \(\mathcal D_9\).
\end{proof}

La différence de géométrie des fibres est ainsi typée sans dupliquer la matrice
APP : la loi additive interne ne fait que translater, tandis que la matrice multiplicative
permute sur les unités et contracte sur le radical.

\begin{proposition}[Structure linéaire du radical]
\label{p12-83-c1-s1:prop:app-radical-lineal}
L'application
\[
 M_3:R_9\longrightarrow R_9,
 \qquad x\longmapsto3x
\]
satisfait
\[
 \ker M_3=\operatorname{im}M_3=\mathfrak m.
\]
L’application induite
\[
 R_9/\mathfrak m\longrightarrow\mathfrak m,
 \qquad [x]\longmapsto3x
\]
est un isomorphisme d’espaces vectoriels sur \(\mathbb F_3\), mais non un
isomorphisme d’anneaux.
\end{proposition}

\begin{proof}
La congruence \(3x\equiv0\pmod9\) équivaut à \(x\equiv0\pmod3\), d’où
le noyau \(\mathfrak m\). L’image est
\(\{[0],[3],[6]\}=\mathfrak m\). Le premier théorème d’isomorphisme donne
l’isomorphisme linéaire. Ce n’est pas un isomorphisme d’anneaux parce que
\(\mathfrak m^2=0\) dans \(R_9\), tandis que le quotient
\(R_9/\mathfrak m\cong\mathbb F_3\) a une unité non nulle.
\end{proof}

La section positive publie le cycle radical non trivial comme
\[
 369\longmapsto693\longmapsto936\longmapsto369,
\]
selon la position à partir de laquelle on lit le triplet. Cette notation conserve
l’ordre des coordonnées ; elle ne transforme pas l’idéal en trois nombres entiers
indépendants.

\subsection{Relèvements et cocycle de retenue}

Soit \(s_0:R_9\to\{0,\ldots,8\}\) la section résiduelle normalisée. La
retenue binaire de la somme est
\[
 \kappa_0(a,b)
 =\frac{s_0(a)+s_0(b)-s_0(a+b)}9.
\]

\begin{lemma}[Identité cocyclique]
\label{p12-83-c1-s1:lem:app-identidad-cociclo-acarreo}
Pour \(a,b,c\in R_9\),
\[
 \kappa_0(a,b)+\kappa_0(a+b,c)
 =\kappa_0(b,c)+\kappa_0(a,b+c).
\]
\end{lemma}

\begin{proof}
Multiplier par neuf et développer l’un quelconque des deux membres produit
\[
 s_0(a)+s_0(b)+s_0(c)-s_0(a+b+c).
\]
Par conséquent, ils sont égaux.
\end{proof}

Nous définissons \(\chi_0(a)=1\) si \(a=[0]_9\) et \(0\) sinon. Comme
\[
 s_+(a)=s_0(a)+9\chi_0(a),
\]
les retenues des deux sections sont reliées par
\[
 \kappa_+(a,b)=
 \kappa_0(a,b)+\chi_0(a)+\chi_0(b)-\chi_0(a+b).
\]
Le changement de représentant modifie le cocycle par un cobord. C’est la
manière exacte de transporter la mémoire lorsque la classe nulle est publiée comme
\(9\) au lieu de \(0\).

\subsection{Recensements intégrés des résidus et des quotients}

En sommant sur toutes les paires \((i,j)\in\mathcal D_9^2\), on obtient
\[
 \sum_{i,j}(i+j)=810,
 \qquad
 \sum_{i,j}ij=2025.
\]
Les résidus positifs satisfont
\[
 \sum_{i,j}R^+_{i,j}=405,
 \qquad
 \sum_{i,j}R^\times_{i,j}=459.
\]
Par les identités relevées,
\[
 \sum_{i,j}Q^+_{i,j}=\frac{810-405}{9}=45,
 \qquad
 \sum_{i,j}Q^\times_{i,j}=\frac{2025-459}{9}=174.
\]

\begin{theorem}[Séparation exacte de l’excès somme--produit]
\label{p12-83-c1-s1:thm:app-separacion-exceso}
La différence intégrée entre produit et somme se décompose de manière unique
en une partie publiée et une partie de quotient :
\[
 \boxed{
  2025-810
  =(459-405)+9(174-45)
  =54+9\cdot129.}
\]
\end{theorem}

\begin{proof}
Nous soustrayons, terme à terme, les identités
\(ij=R^\times_{i,j}+9Q^\times_{i,j}\) et
\(i+j=R^+_{i,j}+9Q^+_{i,j}\), puis sommons sur les quatre-vingt-une paires. Les
recensements précédents donnent l’égalité. L’unicité procède de la division
euclidienne dans chaque paire.
\end{proof}

La valeur \(54\) n’est pas la différence complète entre deux matrices supposées.
C’est le défaut visible entre la matrice multiplicative et l’accumulateur additif
dans la carte positive. Le terme \(9\cdot129\) est la contribution
qui demeure dans les quotients. L’omettre reviendrait à confondre publication et
généalogie.

\subsection{Bloc central de \(2\) positions}

Pour \(k\in\mathbb Z/9\mathbb Z\), nous considérons le bloc symétrique
\[
 B_k=
 \begin{pmatrix}
  k^2&k(k+1)\\
  k(k+1)&(k+1)^2
 \end{pmatrix}\pmod9.
\]
Les diagonales coïncident exactement lorsque
\[
 k^2\equiv(k+1)^2\pmod9
 \quad\Longleftrightarrow\quad
 2k+1\equiv0\pmod9.
\]
Comme \(2^{-1}=5\pmod9\), l’unique solution est \(k=4\). Ainsi apparaît le
bloc
\[
 B_c=
 \begin{pmatrix}7&2\\2&7\end{pmatrix}.
\]
Soit
\[
 R=\begin{pmatrix}0&1\\1&0\end{pmatrix},
 \qquad
 P_\pm=\frac12(I\pm R).
\]
Alors
\[
 B_c=7I+2R=9P_++5P_-.
\]

\begin{theorem}[Unicité et spectre du bloc central]
\label{p12-83-c1-s1:thm:app-bloque-central}
Le bloc \(B_c\) est l’unique bloc adjacent de la famille \(B_k\) ayant
des diagonales égales. Ses valeurs propres dans les sous-espaces symétrique et
antisymétrique sont, respectivement, \(9\) et \(5\).
\end{theorem}

\begin{proof}
La congruence précédente prouve l’unicité. Puisque \(R\) agit par
\(+1\) sur \(\operatorname{im}P_+\) et par \(-1\) sur
\(\operatorname{im}P_-\), la matrice \(7I+2R\) agit par \(9\) et \(5\),
respectivement.
\end{proof}

Les deux lectures ordonnées de ses chiffres donnent
\[
 72+27=99,
 \qquad
 72-27=45=\det B_c,
\]
et les deux diagonales donnent
\[
 77+22=99,
 \qquad
 77-22=55=54+1.
\]
Ces identités sont des contrôles internes du bloc ; elles ne définissent pas à elles seules
le défaut global.

\subsection{Courbure mixte des \(2\) lois}

Dans la carte résiduelle, nous écrivons
\[
 S(x,y)=x+y,
 \qquad
 P(x,y)=xy.
\]
Pour une fonction \(T:X_9\to R_9\), soient
\[
 \Delta_xT(x,y)=T(x+1,y)-T(x,y),
 \qquad
 \Delta_yT(x,y)=T(x,y+1)-T(x,y).
\]
Nous définissons la courbure ordonnée d’une paire de lois par
\[
 \mathcal F(T_x,T_y)
 :=\Delta_x\Delta_yT_y-\Delta_y\Delta_xT_x.
\]

\begin{proposition}[Orientations de la composition mixte]
\label{p12-83-c1-s1:prop:app-curvatura-mixta}
On a
\[
 \mathcal F(S,S)=0,
 \qquad
 \mathcal F(P,P)=0,
\]
et, pour les compositions mixtes,
\[
 \mathcal F(S,P)=+1,
 \qquad
 \mathcal F(P,S)=-1
 \quad\text{dans }R_9.
\]
\end{proposition}

\begin{proof}
Les différences secondes mixtes de la loi linéaire \(S\) s’annulent. Pour
\(P(x,y)=xy\), les deux différences mixtes pures valent un et s’annulent lorsqu’elles sont
comparées dans le même ordre. En combinant une loi linéaire et une loi bilinéaire,
il reste une unique unité dont le signe change lorsque l’ordre des feuillets est inversé.
\end{proof}

La courbure ne crée pas une autre matrice APP. Elle enregistre que l’ordre
\(\Sigma\to\Pi\) et l’ordre \(\Pi\to\Sigma\) sont des orientations opposées,
tandis que les secteurs purs restent plats. Cette triade
\(-1,0,+1\) sera typée dans l’unité TRIT.

\subsection{Obstructions de la matrice arithmétique de l’APP : translations, radical local, bloc central, courbure mixte et conservation de la retenue}

\begin{criterion}[Critères de réfutation de la matrice APP]
L’unité est réfutée si :
\begin{enumerate}
 \item la loi additive interne cesse d’agir par translations ;
 \item une ligne multiplicative non unitaire est déclarée bijective ;
 \item on identifie \(R_9/\mathfrak m\) à \(\mathfrak m\) comme anneau ;
 \item l’excès \(1215\) est remplacé par \(54\) sans conserver les quotients ;
 \item le bloc central est choisi sans résoudre \(2k+1=0\pmod9\) ;
 \item l’inversion de l’ordre mixte ne change pas le signe de la courbure ;
 \item la retenue est éliminée lors du passage entre les sections \(s_0\) et
       \(s_+\).
\end{enumerate}
\end{criterion}

L’APP est maintenant construite comme support, catégorie de trajectoires, matrice
multiplicative nonadique, loi additive interne, radical local, cocycle de
retenue et orientation mixte. L’unité
suivante ne résumera pas ces structures : elle prendra la triade d’orientations
qui vient d’émerger et construira l’état complet du TRIT avec ses trois régimes
quadratiques.


% Unidad U002. Redacción autónoma a partir del generador integral 1063/live.

\section{Complétion gnomonique et support arithmétique de l’APP}
\label{p12-83-c1-s2:chap:app-completacion-gnomonica}

La présente unité construit le support de l’APP à partir de la cellule
nonadique de la \cref{p12-83-c1-s4:chap:segmentos-marcas-intersticios}. Aucun
plan continu, réseau infini ni constante numérique n’est présupposé. L’objet de
départ est fini : deux copies ordonnées de la carte positive de
\(\mathbb Z/9\mathbb Z\). La complétion de leurs interstices intérieurs
produit exactement quatre-vingt-un états, et la loi de translation résiduelle
transforme ce support carré en un tore discret.

Dans toute l’unité, APP désigne une structure mathématique et non une image.
Le sigle APP nomme désormais cette structure ; les propriétés qui
suivent dépendent de ses objets et de ses morphismes, non d’une explicitation verbale du
sigle.

\subsection{Correspondance entre \(2\) cartes d’une même classe résiduelle}

On définit
\[
 \mathcal D_9:=\{1,2,\ldots,9\},
 \qquad
 \mathcal D_8:=\{1,2,\ldots,8\}.
\]
L’application de publication résiduelle est
\[
 \lambda_9:\mathcal D_9\longrightarrow\mathbb Z/9\mathbb Z,
 \qquad
 \lambda_9(a)=[a]_9.
\]
Par conséquent, \(\lambda_9(9)=[0]_9\). Son inverse comme application
d’ensembles est la section positive
\[
 s_+([a]_9)=
 \begin{cases}
  a,&[a]_9\neq[0]_9,\quad 1\leq a\leq8,\\
  9,&[a]_9=[0]_9.
 \end{cases}
\]
Les deux expressions \(9\) et \([0]_9\) ne sont pas deux phases : ce sont deux
représentants du même élément dans des cartes différentes.

Pour deux coordonnées, nous écrirons
\[
 \lambda_9^{\times2}:\mathcal D_9^2
       \longrightarrow X_9:=(\mathbb Z/9\mathbb Z)^2,
 \qquad
 (a,b)\longmapsto([a]_9,[b]_9).
\]
Cette application est une bijection d’ensembles. La carte
\(\mathcal D_9^2\) rend visibles les marques positives ; le groupe \(X_9\)
rend visibles la somme résiduelle et les translations.

\begin{lemma}[Changement de carte sans perte d’état]
\label{p12-83-c1-s2:lem:app-cambio-carta-sin-perdida}
Les applications \(\lambda_9^{\times2}\) et
\(s_+^{\times2}\) sont inverses. En particulier, l’écriture positive et
l’écriture résiduelle contiennent exactement les mêmes quatre-vingt-un états.
\end{lemma}

\begin{proof}
Dans chaque coordonnée, on a
\(\lambda_9\circ s_+=\operatorname{id}_{\mathbb Z/9\mathbb Z}\) et
\(s_+\circ\lambda_9=\operatorname{id}_{\mathcal D_9}\). Le produit
cartésien de ces identités fournit les deux identités requises.
\end{proof}

\subsection{Noyau intérieur et complément gnomonique}

Le noyau intérieur est
\[
 \mathcal N_8:=\mathcal D_8\times\mathcal D_8,
 \qquad |\mathcal N_8|=8^2=64.
\]
Dans le plan des indices, \(\mathcal N_8\) est le carré qui n’utilise la
marque de clôture dans aucune coordonnée. Pour publier également la classe nulle,
il faut compléter la neuvième ligne et la neuvième colonne. Nous définissons le
complément gnomonique disjoint par
\[
 \mathcal G_9:=
  (\{9\}\times\mathcal D_9)
  \sqcup
  (\mathcal D_8\times\{9\}).
\]
Le coin \((9,9)\) appartient à la première pièce ; c’est pourquoi la seconde utilise
\(\mathcal D_8\) et non \(\mathcal D_9\). En conséquence,
\[
 |\mathcal G_9|=9+8=17,
 \qquad
 \mathcal D_9^2=\mathcal N_8\sqcup\mathcal G_9,
 \qquad
 81=64+17.
\]

\begin{proposition}[Minimalité de la complétion]
\label{p12-83-c1-s2:prop:app-minimalidad-completacion}
Soit \(Y\) un sous-ensemble de \(\mathcal D_9^2\) qui contient
\(\mathcal N_8\) et dont la projection sur chaque coordonnée est tout
\(\mathcal D_9\). Si, de plus, \(Y\) contient toutes les paires nécessaires
pour translater une coordonnée tandis que l’autre reste arbitraire,
alors \(\mathcal G_9\subseteq Y\). Par conséquent, l’ajout de dix-sept
états est la complétion minimale compatible avec les deux translations.
\end{proposition}

\begin{proof}
Pour translater la première coordonnée à travers la classe de clôture tandis que
la seconde reste à une valeur arbitraire \(b\in\mathcal D_9\), toutes les paires
\((9,b)\) sont nécessaires. On obtient neuf états. Pour
translater la seconde coordonnée à travers la clôture tandis que la première prend
une valeur arbitraire \(a\in\mathcal D_8\), les huit paires
\((a,9)\) restantes sont nécessaires. Le coin \((9,9)\) a déjà été compté dans la première
famille. Au total, \(9+8=17\) états sont nécessaires, exactement ceux de
\(\mathcal G_9\).
\end{proof}

L’hypothèse de compatibilité avec les deux translations est essentielle. Il suffirait
d’ajouter une seule marque par coordonnée pour que les deux projections soient
surjectives, mais cet ensemble ne contiendrait pas les états de frontière
nécessaires pour déplacer une coordonnée en gardant l’autre fixe. L’APP exige le
produit complet, non deux axes déconnectés.

\subsection{Réalisation exacte de l’intervalle interstitiel par des cellules APP et des coordonnées de frontière}

Soit \(P_9\) le chemin ordonné à neuf marques et huit arêtes intérieures
\[
 E(P_9)=\{e_1,\ldots,e_8\}.
\]
La clôture des extrémités incorpore une arête \(e_9\) et produit le cycle
\(C_9\) :
\[
 E(C_9)=E(P_9)\sqcup\{e_9\}.
\]
Les paires d’interstices intérieurs forment \(E(P_9)^2\). Les paires
d’interstices du cycle complet forment \(E(C_9)^2\). La différence est
\begin{align}
 E(C_9)^2\setminus E(P_9)^2
   &=(\{e_9\}\times E(C_9))
     \sqcup(E(P_9)\times\{e_9\}).
\label{p12-83-c1-s2:eq:app-complemento-intersticial}
\end{align}
La décomposition est disjointe parce que le coin \((e_9,e_9)\) est réservé à
la première pièce.

\begin{theorem}[Équivalence du modèle de marques et du modèle
d’interstices]
\label{p12-83-c1-s2:thm:app-equivalencia-modelos}
L'application
\[
 \beta:E(C_9)^2\longrightarrow\mathcal D_9^2,
 \qquad
 \beta(e_i,e_j)=(i,j),
\]
est bijective. Ses restrictions satisfont
\[
 \beta(E(P_9)^2)=\mathcal N_8,
 \qquad
 \beta(E(C_9)^2\setminus E(P_9)^2)=\mathcal G_9.
\]
Par conséquent, l’identité \(81=64+17\) compte simultanément des paires de
marques publiées et des paires d’interstices fermés.
\end{theorem}

\begin{proof}
Chaque arête de \(C_9\) a un indice unique dans \(\mathcal D_9\).
L’application \((i,j)\mapsto(e_i,e_j)\) est l’inverse explicite de \(\beta\).
Si \(i,j\leq8\), la paire appartient à \(E(P_9)^2\) et son image appartient à
\(\mathcal D_8^2\). Si au moins l’un des indices est neuf, la paire appartient
au complément de \cref{p12-83-c1-s2:eq:app-complemento-intersticial} ; séparer le cas
\(i=9\) du cas \(i\leq8,j=9\) produit exactement les deux pièces de
\(\mathcal G_9\).
\end{proof}

\subsection{Clôture toroïdale et translations}

Soient
\[
 \mathbf e_1=([1]_9,[0]_9),
 \qquad
 \mathbf e_2=([0]_9,[1]_9).
\]
L’ensemble des directions élémentaires est
\[
 \mathscr D_{\rm APP}:=
 \{+\mathbf e_1,-\mathbf e_1,+\mathbf e_2,-\mathbf e_2\}.
\]
Nous définissons le graphe de Cayley orienté
\[
 \Gamma_{\rm APP}:=\operatorname{Cay}(X_9,\mathscr D_{\rm APP}).
\]
Ses sommets sont les quatre-vingt-un éléments de \(X_9\). Pour chaque
\(x\in X_9\) et \(d\in\mathscr D_{\rm APP}\), il existe une arête orientée
\(x\xrightarrow{d}x+d\). L’arête inverse porte l’étiquette \(-d\).

Dans la carte positive, la translation unitaire s’écrit au moyen du successeur
publié et de sa retenue :
\[
 u_+(a)=
 \begin{cases}
  a+1,&a<9,\\
  1,&a=9,
 \end{cases}
 \qquad
 \kappa_+(a)=
 \begin{cases}
  0,&a<9,\\
  1,&a=9.
 \end{cases}
\]
L’identité relevée est
\begin{equation}
 a+1=u_+(a)+9\kappa_+(a).
\label{p12-83-c1-s2:eq:app-sucesor-publicado-acarreo}
\end{equation}
En appliquant \(\lambda_9\), le terme de retenue s’annule et il reste
\[
 \lambda_9(u_+(a))=\lambda_9(a)+[1]_9.
\]

\begin{lemma}[Compatibilité entre clôture, résidu et relèvement]
\label{p12-83-c1-s2:lem:app-compatibilidad-cierre-residuo}
L’ajout de \(e_9\), l’identification \(9\sim[0]_9\) et la translation
\(x\mapsto x+\mathbf e_i\) décrivent trois opérations compatibles dans
des domaines distincts. Le quotient résiduel publie le retour ; la retenue de
\cref{p12-83-c1-s2:eq:app-sucesor-publicado-acarreo} conserve le relèvement que le
quotient omet.
\end{lemma}

\begin{proof}
La clôture topologique incorpore l’arête qui relie les extrémités du chemin. La
carte résiduelle identifie l’extrémité publiée à la classe nulle. Si
\(a<9\), le successeur a une retenue nulle ; si \(a=9\), l’identité relevée
est \(10=1+9\). Les deux cas descendent à la même somme résiduelle, mais la
valeur de \(\kappa_+\) les distingue.
\end{proof}

\begin{corollary}[Support toroïdal]
\label{p12-83-c1-s2:cor:app-soporte-toroidal}
Identifier les côtés opposés de la carte \(\mathcal D_9^2\) produit une
réalisation géométrique du groupe \(X_9\). Chaque translation élémentaire est
définie en tout sommet et possède un inverse.
\end{corollary}

\subsection{Catégorie des trajectoires APP : objets, morphismes de prolongement et composition compatible}

\begin{definition}[Trajectoire orientée]
Une trajectoire de longueur \(n\geq0\) est une paire
\[
 \gamma=(x_0;d_1,\ldots,d_n),
 \qquad
 x_0\in X_9,\qquad d_k\in\mathscr D_{\rm APP}.
\]
Ses sommets sont
\[
 x_k=x_0+\sum_{r=1}^k d_r,
 \qquad 0\leq k\leq n.
\]
La source est \(s(\gamma)=x_0\) et la cible est
\(t(\gamma)=x_n\). Pour \(n=0\), on obtient la trajectoire identité en
\(x_0\).
\end{definition}

Si \(t(\gamma)=s(\delta)\), leur concaténation est
\[
 \delta\circ\gamma
  =(x_0;d_1,\ldots,d_n,e_1,\ldots,e_m),
\]
où \(e_1,\ldots,e_m\) sont les directions de \(\delta\). La
concaténation est associative et les trajectoires de longueur zéro sont des
identités. Nous obtenons ainsi la catégorie libre
\[
 \mathsf{Path}(\Gamma_{\rm APP}).
\]

\begin{proposition}[Finitude du support et dénombrabilité des trajectoires]
\label{p12-83-c1-s2:prop:app-soporte-finito-trayectorias-numerables}
L’ensemble des états de l’APP est fini, de cardinal \(81\). Pour une source
fixée, il existe \(4^n\) mots de directions de longueur \(n\). L’union
de toutes les trajectoires finies est dénombrable.
\end{proposition}

\begin{proof}
La première affirmation procède de \(|X_9|=9^2\). À chacun des \(n\)
pas, il y a quatre choix indépendants, d’où \(4^n\). Enfin,
\(\bigcup_{n\geq0}X_9\times\mathscr D_{\rm APP}^n\) est une union dénombrable
d’ensembles finis.
\end{proof}

Cette proposition sépare deux niveaux qui ne seront pas identifiés : l’APP possède un
support fini, mais admet des trajectoires de toute longueur. Le
prolongement d’une route ne crée pas de nouveaux sommets du support ; il crée une nouvelle
généalogie sur ceux-ci.

\subsection{Matrice multiplicative et loi additive interne sur un unique support}

Pour \(n\in\mathbb Z_{>0}\), la racine numérique positive est
\[
 \operatorname{dr}_9(n):=1+((n-1)\bmod9).
\]
Sur la carte positive, nous définissons la matrice principale et la loi interne
\[
 \Sigma(a,b):=\operatorname{dr}_9(a+b),
 \qquad
 \Pi(a,b):=\operatorname{dr}_9(ab).
\]
Les deux applications ont le même domaine \(\mathcal D_9^2\) et le même
codomaine \(\mathcal D_9\), mais elles n’ont pas le même statut : \(\Pi\) publie
l’unique matrice APP et \(\Sigma\) engendre ses accumulations. Leurs versions résiduelles
sont
\[
 \bar\Sigma([a]_9,[b]_9)=[a+b]_9,
 \qquad
 \bar\Pi([a]_9,[b]_9)=[ab]_9.
\]
Le changement de carte est compatible avec les deux opérations :
\begin{align*}
 \lambda_9\circ\Sigma
   &=\bar\Sigma\circ\lambda_9^{\times2},\\
 \lambda_9\circ\Pi
   &=\bar\Pi\circ\lambda_9^{\times2}.
\end{align*}

Pour ne pas perdre le quotient entier, on introduit les relèvements
\begin{align}
 a+b &= R^+_{a,b}+9Q^+_{a,b},
 &1\leq R^+_{a,b}\leq9,
\label{p12-83-c1-s2:eq:app-lift-suma}\\
 ab &= R^\times_{a,b}+9Q^\times_{a,b},
 &1\leq R^\times_{a,b}\leq9.
\label{p12-83-c1-s2:eq:app-lift-producto}
\end{align}
Ici
\[
 R^+_{a,b}=\Sigma(a,b),
 \qquad
 R^\times_{a,b}=\Pi(a,b),
\]
et les quotients sont déterminés de manière univoque par
\[
 Q^+_{a,b}=\frac{a+b-R^+_{a,b}}9,
 \qquad
 Q^\times_{a,b}=\frac{ab-R^\times_{a,b}}9.
\]

\begin{definition}[APP généalogique]
\label{p12-83-c1-s2:def:app-genealogica-autonoma}
La structure APP à cette étape est le tuple
\[
 \operatorname{APP}:=
 \bigl(
  X_9,
  \Gamma_{\rm APP},
  \mathsf{Path}(\Gamma_{\rm APP}),
  \Sigma,\Pi,
  R^+,Q^+,R^\times,Q^\times
 \bigr).
\]
La généalogie d’une trajectoire conserve, pour chaque sommet visité, la
position, la direction d’entrée, l’entrée multiplicative, l’accumulation
additive et leurs résidus et quotients respectifs.
\end{definition}

\begin{theorem}[Construction causale du support APP]
\label{p12-83-c1-s2:thm:app-cadena-constructiva-autonoma}
La chaîne
\[
 \mathcal D_9
 \longrightarrow
 \mathcal D_8^2
 \longrightarrow
 \mathcal D_8^2\sqcup\mathcal G_9
 \longrightarrow
 \mathcal D_9^2
 \xrightarrow{\lambda_9^{\times2}}
 X_9
 \longrightarrow
 \Gamma_{\rm APP}
 \longrightarrow
 \mathsf{Path}(\Gamma_{\rm APP})
 \longrightarrow
 (\Sigma,\Pi)
\]
construit, sans perte de cardinalité, le support, le mouvement, les trajectoires,
la matrice multiplicative de l’APP et sa loi additive interne.
\end{theorem}

\begin{proof}
La \cref{p12-83-c1-s2:prop:app-minimalidad-completacion} construit la carte complète à
partir du noyau. Le \cref{p12-83-c1-s2:thm:app-equivalencia-modelos} identifie son
modèle interstitiel. Le \cref{p12-83-c1-s2:lem:app-cambio-carta-sin-perdida} transporte
les quatre-vingt-un états vers \(X_9\). Le graphe de Cayley incorpore les quatre
translations réversibles sans changer les sommets. La catégorie libre
incorpore toutes les concaténations finies. Enfin, \(\Sigma\) et
\(\Pi\) sont des applications totales distinctes sur le même domaine ; les
équations \cref{p12-83-c1-s2:eq:app-lift-suma,p12-83-c1-s2:eq:app-lift-producto} conservent la partie
que la réduction résiduelle ne publie pas.
\end{proof}

\subsection{Décomposition de l’orthogramme APP en noyau \(8\times8\), gnomon de \(17\) états et carte \(9\times9\) à côtés opposés identifiés}

La \cref{fig:app-ortograma} doit être lue en trois couches : le carré intérieur
de soixante-quatre états, le gnomon de dix-sept états et
l’identification des côtés opposés. La carte de neuf sur neuf complète ainsi le
noyau de huit sur huit par le gnomon de clôture. La matrice
multiplicative et la loi additive interne sont évaluées en chaque sommet du même
état, conservent des relèvements résidu--quotient distincts et ne constituent pas
deux carrés APP.

\subsection{Obstructions à la fidélité du support APP : cardinalités \(81/17\), translations réversibles et séparation des feuillets somme et produit}

\begin{criterion}[Critères de réfutation du support APP]
La construction est réfutée si l’un quelconque des faits
suivants se produit :
\begin{enumerate}
 \item le support complet a un cardinal différent de \(81\) ;
 \item le complément gnomonique n’a pas pour cardinal \(17\) ;
 \item un état de \(X_9\) est dépourvu de l’une des quatre translations ;
 \item on identifie \(\Sigma\) et \(\Pi\) parce qu’ils coïncident en une valeur ;
 \item on élimine \(Q^+\) ou \(Q^\times\) lors du retour par une frontière ;
 \item les trajectoires de longueurs différentes sont présentées comme de nouveaux
       sommets du support ;
 \item on interprète \(9\) et \([0]_9\) comme des phases indépendantes.
\end{enumerate}
\end{criterion}

L’unité a construit l’APP jusqu’à son évaluation généalogique complète. Elle n’a pas encore
orienté la différence entre la matrice et son accumulateur, elle n’a pas
introduit le TRIT
et elle n’a pas défini le Topological Phase Kernel. L’unité suivante étudiera
la structure algébrique des lignes et des colonnes, le secteur des unités et le
cocycle de retenue qui prépare cette orientation.


% Unidad U005. Redacción autónoma desde 03c/03h/20 de la autoridad 1063/live.

\long\def\HMTDeferredTPKBase{%
\section{Topological Phase Kernel : base observable et état enrichi}
\label{p12-83-c1-s3:chap:tpk-categoria-estado-enriquecido}

L’APP a fourni un support, une matrice multiplicative et une loi additive
interne ; TRIT a
fourni une orientation relevée. Le Topological Phase Kernel, ci-après
TPK, spécifie comment ces données sont transportées conjointement, quelles
coordonnées sont conservées lors de la clôture d’une route et quels lecteurs peuvent être publiés
sans identifier la publication à l’état. Son objet mathématique est une
catégorie sur une base de chemins, munie d’une famille explicite de
fibres et de relations de prolongement.

\subsection{La base double des chemins APP}

Soit \(\mathsf P_{\rm APP}\) la catégorie libre construite dans la
\cref{p12-83-c1-s2:chap:app-completacion-gnomonica}. Nous définissons
\[
 \mathsf P_{\rm APP}^{(2)}
 :=\mathsf P_{\rm APP}\times\mathsf P_{\rm APP}.
\]
Un objet est une paire de positions
\((p^\Sigma,p^\Pi)\in X_9^2\). Un morphisme est une paire de trajectoires ; l’une
d’elles peut être l’identité. La première composante est évaluée par
\(\Sigma\), la seconde par \(\Pi\). Ce ne sont pas deux réseaux : ce sont deux curseurs
typés sur le même support APP.

L’ensemble des directions cardinales est
\[
 \mathscr D_{\rm card}:=\{N,E,S,O\},
\]
avec décalages
\[
 \delta(N)=(-1,0),\quad
 \delta(E)=(0,1),\quad
 \delta(S)=(1,0),\quad
 \delta(O)=(0,-1)
 \quad\text{dans }X_9.
\]
L’involution d’orientation est
\[
 \overline N=S,
 \quad \overline S=N,
 \quad \overline E=O,
 \quad \overline O=E.
\]

\subsection{Phase nonadique, secteur dodécaphasique et calendrier}

L’horloge complète de cette étape est
\[
 \Theta_{108}:=\mathbb Z/108\mathbb Z.
\]
Pour \(\theta\in\Theta_{108}\), soit
\(\widetilde\theta\in\{0,\ldots,107\}\) son représentant normalisé.
Nous définissons trois coordonnées distinctes :
\begin{align*}
 g_0(\theta)&=[\widetilde\theta+1]_9\in C_9,\\
 r(\theta)&=1+(\widetilde\theta\bmod9)\in\mathcal D_9,\\
 \beta(\theta)&=
 \left[\left\lfloor\frac{\widetilde\theta}{9}\right\rfloor\right]_{12}
 \in C_{12}.
\end{align*}
La première est la classe de phase, la deuxième son étiquette positive et la troisième
le secteur dodécaphasique. L’application \(\beta\) est une application par blocs ;
ce n’est pas la réduction canonique \(\theta\bmod12\).

La signature et le feuillet actif sont
\[
 \tau(\theta):=\varepsilon_9(r(\theta))\in\mathsf{TRIT},
\]
et
\[
 \mu(\theta)=
 \begin{cases}
  \Sigma,&\tau(\theta)=+1,\\
  \Pi,&\tau(\theta)=-1,\\
  \varnothing,&\tau(\theta)=0.
 \end{cases}
\]
Le symbole \(\varnothing\) désigne un pas de seuil où aucun curseur
n’est translaté ; il ne désigne pas une absence d’état.

\begin{definition}[Configuration observable]
\label{p12-83-c1-s3:def:tpk-configuracion-observable}
La base observable est
\[
 \mathcal Q_{\rm obs}
 :=X_9^2\times\mathscr D_{\rm card}^2\times\Theta_{108}.
\]
Un élément s’écrit
\[
 q=(p^\Sigma,p^\Pi,d^\Sigma,d^\Pi,\theta).
\]
Il contient deux positions, deux orientations cardinales et un indice de
calendrier.
\end{definition}

\subsection{Sélecteur de feuillet \(M_{\rm ph}:\mathcal D_9\to\mathsf{TRIT}\) et évolution conjointe de la position, de l’orientation et du calendrier}

La fonction
\[
 M_{\rm ph}:\mathcal D_9\longrightarrow\mathsf{TRIT},
 \qquad M_{\rm ph}(r)=\varepsilon_9(r)
\]
est le sélecteur interne de feuillet. Il sélectionne l’évaluation qui agit ; il ne déplace
par lui-même aucune position.

Nous définissons l’indicateur d’inversion de l’orientation par
\[
 \eta_{27}(\theta)=
 \begin{cases}
  -1,&\widetilde\theta+1\in\{27,54,81,108\},\\
  +1,&\text{sinon}.
 \end{cases}
\]
Les nouvelles directions sont
\[
 d^{\nu,+}=
 \begin{cases}
  \overline{d^\nu},&\eta_{27}(\theta)=-1,\\
  d^\nu,&\eta_{27}(\theta)=+1,
 \end{cases}
 \qquad \nu\in\{\Sigma,\Pi\}.
\]
Les positions évoluent par
\[
 (p^{\Sigma,+},p^{\Pi,+})=
 \begin{cases}
 (p^\Sigma+\delta(d^\Sigma),p^\Pi),&\tau(\theta)=+1,\\
 (p^\Sigma,p^\Pi+\delta(d^\Pi)),&\tau(\theta)=-1,\\
 (p^\Sigma,p^\Pi),&\tau(\theta)=0.
 \end{cases}
\]

\begin{definition}[Évolution observable]
\label{p12-83-c1-s3:def:tpk-evolucion-observable}
L'application
\[
 F_{\rm obs}:\mathcal Q_{\rm obs}\longrightarrow\mathcal Q_{\rm obs}
\]
est défini par
\[
 F_{\rm obs}(q)=
 (p^{\Sigma,+},p^{\Pi,+},d^{\Sigma,+},d^{\Pi,+},[\theta+1]_{108}).
\]
La catégorie libre du graphe \(q\to F_{\rm obs}(q)\) est notée
\(\mathsf P_{\rm obs}\).
\end{definition}

Chaque flèche de \(\mathsf P_{\rm obs}\) induit une paire de trajectoires dans
l’APP, de sorte qu’il existe un foncteur
\[
 B:\mathsf P_{\rm obs}\longrightarrow
 \mathsf P_{\rm APP}^{(2)}.
\]
Si \(\tau=+1\), la deuxième composante de \(B\) est l’identité ; si
\(\tau=-1\), c’est la première ; si \(\tau=0\), les deux le sont.

\begin{proposition}[Séparation entre sélection, mouvement et observation]
\label{p12-83-c1-s3:prop:tpk-separacion-seleccion-movimiento}
Le sélecteur \(M_{\rm ph}\), les translations
\(T_d(p)=p+\delta(d)\) et tout échantillonnage temporel possèdent des domaines et des
codomaines distincts. Ce ne sont pas des opérateurs interchangeables.
\end{proposition}

\begin{proof}
\(M_{\rm ph}\) prend une étiquette de phase et renvoie une orientation tritique.
\(T_d\) est un automorphisme de \(X_9\). Un échantillonnage agit sur une
suite ou une catégorie de flèches. Identifier deux d’entre eux violerait le
typage. \(F_{\rm obs}\) les compose dans un ordre déclaré sans effacer cette
distinction.
\end{proof}

\subsection{Décomposition de l’état enrichi en \(6\) familles fibrées}

Sur chaque \(q\in\mathcal Q_{\rm obs}\), nous fixons la fibre totale
\[
 \mathcal F_q=
 \mathsf O_q\times
 \mathsf H_q\times
 \mathsf C_q\times
 \mathsf S_q\times
 \mathsf B_q\times
 \Lambda_q.
\]
Chaque facteur a un type et une fonction propres.

\subsubsection{Détermination conjointe de l’orientation et du feuillet de l’état enrichi}

\[
 \mathsf O_q\subseteq
 \mathscr D_{\rm card}^2\times\{+1,-1\}.
\]
Les deux directions coïncident avec celles stockées dans \(q\) ; la dernière
coordonnée enregistre l’orientation du feuillet lorsque la réalisation l’exige.
L’étiquette de feuillet ne remplace pas les directions cardinales.

\subsubsection{Décomposition euclidienne de l’état en résidu et quotient compatibles}

La première fibre de relèvement est
\[
 \mathsf H^{(1)}=\mathcal D_9\times\mathbb Z,
 \qquad
 \operatorname{rec}_9(r,u)=r+9u.
\]
Pour chaque entier \(n\), il existe d’uniques
\[
 r=\operatorname{dr}_9(n),
 \qquad
 u=\frac{n-r}{9}.
\]
Une profondeur supérieure consiste en une tour déclarée de ces fibres et de leurs
applications de troncature. Le terme « Hensel » ne remplace pas la définition de
cette tour.

\subsubsection{Opérateur de retenue et actualisation conjointe du résidu, du quotient et de la mémoire}

\(\mathsf C_q\) est un groupe abélien. Dans les lecteurs de blocs décimaux,
on prend \(\mathsf C_q=\mathbb Z\). L’accroissement n’est pas une coordonnée
indépendante, mais une cochaîne sur les flèches :
\[
 \kappa(r)\in\mathsf C_{q'}
 \quad\text{pour }r:q\to q'.
\]

\subsubsection{Construction de la signature entière à partir de l’orientation, de la retenue et de la frontière}

La signature est dérivée par une application typée
\[
 \operatorname{Sig}_q:
 \mathsf H_q\times\mathsf C_q\times\Lambda_q
 \longrightarrow\mathsf S_q.
\]
Dans la réalisation de frontière nonadique, on conserve individuellement
\[
 q_\partial,a_\partial,c_\partial\in\mathbb F_3^6,
 \qquad
 \operatorname{colw}\in\mathbb Z_{\geq0}^{6},
 \qquad
 \rho_W,r_\partial\in\mathbb F_3^3.
\]
Ces coordonnées sont la marge, la combinaison axiale, la condition cubique, les poids de
colonne, la classe résiduelle de Witt et le retour ternaire. Elles ne sont pas condensées en un
unique mot si une preuve ultérieure utilise l’une d’elles.

\subsubsection{Cylindre et frontière}

\[
 \mathsf B_q:=\mathsf{Cyl}_q\times\mathsf{Fr}_q.
\]
Le cylindre localise une famille de préfixes ; la frontière détermine quelles
extensions sont compatibles. La survie sera une relation entre
des états munis des deux données, non un nombre ajouté a posteriori.

\subsubsection{Registre généalogique de transport}

\(\Lambda_q\) contient les registres de transport dont l’extrémité est
\(q\). Un prolongement \(r:q\to q'\) agit par concaténation :
\[
 \lambda\longmapsto r\circ\lambda.
\]
Le registre généalogique conserve la route ; son extrémité observable ne la détermine pas.

\subsection{Transport des fibres}

Pour chaque flèche \(r:q\to q'\), on déclare des applications
\[
 \mathsf H(r):\mathsf H_q\to\mathsf H_{q'},
 \quad
 \mathsf C(r):\mathsf C_q\to\mathsf C_{q'},
 \quad
 \Lambda(r):\Lambda_q\to\Lambda_{q'}.
\]
Elles respectent les identités et la composition, et \(\mathsf C(r)\) est un homomorphisme.
La retenue satisfait
\begin{equation}
 \kappa(r_2\circ r_1)
 =\kappa(r_2)+\mathsf C(r_2)\bigl(\kappa(r_1)\bigr).
\label{p12-83-c1-s3:eq:tpk-cocadena-transporte}
\end{equation}
Par conséquent, le transport des coordonnées déterminées est
\begin{align*}
 h'&=\mathsf H(r)(h),\\
 c'&=\mathsf C(r)(c)+\kappa(r),\\
 \lambda'&=\Lambda(r)(\lambda),\\
 \sigma'&=\operatorname{Sig}_{q'}(h',c',\lambda').
\end{align*}

\begin{lemma}[Associativité de la retenue transportée]
\label{p12-83-c1-s3:lem:tpk-asociatividad-acarreo}
L’équation \cref{p12-83-c1-s3:eq:tpk-cocadena-transporte} garantit que transporter un
état par \(r_1\) puis par \(r_2\) produit la même retenue que
le transporter par \(r_2\circ r_1\).
\end{lemma}

\begin{proof}
Le transport successif donne
\[
 \mathsf C(r_2)\bigl(\mathsf C(r_1)c+\kappa(r_1)\bigr)+\kappa(r_2).
\]
Par fonctorialité de \(\mathsf C\), le premier terme est
\(\mathsf C(r_2r_1)c\) ; la somme des deux autres est
\(\kappa(r_2r_1)\) d’après \cref{p12-83-c1-s3:eq:tpk-cocadena-transporte}.
\end{proof}

\subsection{Relations d’extension}

L’orientation et la frontière ne sont pas forcées par une fonction universelle.
Pour chaque flèche \(r:q\to q'\), on déclare une relation
\[
 \operatorname{Ext}_r\subseteq\mathcal F_q\times\mathcal F_{q'}.
\]
On exige les deux conditions
\[
 \Delta_q\subseteq\operatorname{Ext}_{\operatorname{id}_q},
 \qquad
 \operatorname{Ext}_{r_2}\circ\operatorname{Ext}_{r_1}
 \subseteq\operatorname{Ext}_{r_2\circ r_1}.
\]
La première conserve les identités ; la seconde conserve la composition. Une
relation peut admettre plusieurs extensions locales et n’en laisser qu’une après un
horizon plus grand.

\begin{definition}[Topological Phase Kernel]
\label{p12-83-c1-s3:def:tpk-categoria-autonoma}
Un objet du TPK est un tuple
\[
 x=(q,o,h,c,\sigma,C,b,\lambda)
\]
tel que
\[
 q\in\mathcal Q_{\rm obs},\quad
 o\in\mathsf O_q,\quad
 h\in\mathsf H_q,\quad
 c\in\mathsf C_q,\quad
 \lambda\in\Lambda_q,
\]
\[
 \sigma=\operatorname{Sig}_q(h,c,\lambda),
 \qquad
 (C,b)\in\mathsf{Cyl}_q\times\mathsf{Fr}_q.
\]
Nous associons à chaque objet sa composante fibrale
\[
 \mathbf f(x):=(o,h,c,\sigma,(C,b),\lambda)\in\mathcal F_q.
\]
Un morphisme \(x\to x'\) est un triplet \((r,x,x')\) dans lequel
\(r:q\to q'\) appartient à \(\mathsf P_{\rm obs}\), les coordonnées
transportables satisfont les lois précédentes et
\((\mathbf f(x),\mathbf f(x'))\in
\operatorname{Ext}_r\). La composition est
\[
 (r_2,x',x'')\circ(r_1,x,x')=(r_2r_1,x,x''),
\]
et l’identité de \(x\) est \((\operatorname{id}_q,x,x)\).
\end{definition}

\begin{theorem}[Structure catégorielle du TPK]
\label{p12-83-c1-s3:thm:tpk-es-categoria}
La \cref{p12-83-c1-s3:def:tpk-categoria-autonoma} définit une catégorie et l’application
\[
 p:\mathsf{TPK}\longrightarrow\mathsf P_{\rm obs},
 \qquad p(x)=q,\qquad p(r,x,x')=r,
\]
est un foncteur.
\end{theorem}

\begin{proof}
La composition des routes est associative. Le
\cref{p12-83-c1-s3:lem:tpk-asociatividad-acarreo} prouve la compatibilité de la
coordonnée de retenue ; la fonctorialité déclarée de \(\mathsf H\),
\(\mathsf C\) et \(\Lambda\) prouve celle des autres coordonnées
déterminées. La stabilité de \(\operatorname{Ext}\) garantit que le
composé reste admissible, et sa condition diagonale fournit les
identités. L’application \(p\) conserve les identités et la composition par
construction.
\end{proof}

\subsection{Projections et pertes typées}

Les projections élémentaires sont
\[
 \phi(x)=g_0(\theta),
 \qquad
 \tau(x)=\varepsilon_9(r(\theta)),
 \qquad
 \mu(x)=\mu(\theta).
\]
La projection vers l’état TRIT local est
\[
 \pi_{\rm TRIT}(x)
 =\bigl(\tau,\mu,o,h,c,\sigma,\lambda,\phi\bigr).
\]
Elle oublie le cylindre, la frontière et une partie de \(q\) ; elle ne définit pas un autre conteneur.
Deux états peuvent avoir la même image par \(\pi_{\rm TRIT}\) et différer
par l’information oubliée.

\subsection{Critère de survie des routes par conservation de la frontière projective}

Soit \(\mathcal X_{\rm TPK}^{(n)}\) la famille des états à la profondeur
\(n\). Pour \(N\geq n\), nous définissons
\[
 \operatorname{Surv}_n^{(N)}=
 \left\{x_n:
 \begin{array}{l}
 \text{il existe }x_{n+1},\ldots,x_N\text{ avec}\\
 (x_j,x_{j+1})\in\operatorname{Ext}_{r_j}
 \text{ pour }n\leq j<N
 \end{array}\right\}.
\]

\begin{proposition}[Filtration de survie]
\label{p12-83-c1-s3:prop:tpk-filtracion-supervivencia}
Pour \(N\geq n\),
\[
 \operatorname{Surv}_n^{(N+1)}
 \subseteq
 \operatorname{Surv}_n^{(N)}.
\]
Si le schéma est défini à toute profondeur, les états
indéfiniment prolongeables sont
\[
 \operatorname{Surv}_n
 =\bigcap_{N\geq n}\operatorname{Surv}_n^{(N)}.
\]
\end{proposition}

\begin{proof}
Un prolongement jusqu’à \(N+1\) se tronque en un prolongement jusqu’à \(N\),
ce qui prouve l’inclusion. L’intersection exprime exactement la
prolongeabilité à tout horizon fini.
\end{proof}

\subsection{Arbre opératoriel du conteneur}

\begin{figure}[htbp]
 \centering
 \begin{tikzpicture}[
  node distance=7mm and 9mm,
  every node/.style={font=\small,align=center},
  root/.style={draw=black,very thick,rounded corners,inner sep=5pt},
  op/.style={draw=black,rounded corners,inner sep=4pt},
  edge/.style={-{Latex[length=2mm]},semithick}
 ]
  \node[root] (tpk) {TPK\\$\mathcal X_{\rm TPK}^{\rm enr}$};
  \node[op,below left=of tpk] (sel) {sélection\\$M_{\rm ph}$};
  \node[op,below=of tpk] (mov) {transport\\$T_d,F_{\rm obs}$};
  \node[op,below right=of tpk] (fib) {fibres\\$\mathsf H,\mathsf C,\mathsf S,
    \mathsf B,\Lambda$};
  \node[op,below=17mm of mov] (cal) {calendrier\\$C_{12}\hookrightarrow
    C_{108}\twoheadrightarrow C_9$};
  \node[op,below=of cal] (ext) {extension, cylindres et survie};
  \draw[edge] (tpk) -- (sel);
  \draw[edge] (tpk) -- (mov);
  \draw[edge] (tpk) -- (fib);
  \draw[edge] (sel) |- (cal);
  \draw[edge] (mov) -- (cal);
  \draw[edge] (fib) |- (cal);
  \draw[edge] (cal) -- (ext);
 \end{tikzpicture}
 \caption{Inclusion opératorielle du Topological Phase Kernel. Le sélecteur,
 le transport, les lecteurs et les fibres ne sont pas des conteneurs parallèles : ce sont
 des applications ou des coordonnées du même état enrichi.}
 \label{p12-83-c1-s3:fig:tpk-arbol-operatorio-autonomo}
\end{figure}

\subsection{Obstructions de l’état enrichi du TPK : typage du sélecteur, cochaîne de retenue, extension compatible et conservation généalogique}

\begin{criterion}[Critères de réfutation de l’état enrichi]
La définition est réfutée si :
\begin{enumerate}
 \item le sélecteur \(M_{\rm ph}\) est utilisé comme translation ;
 \item on identifie \(\theta\), \(g_0(\theta)\), \(r(\theta)\) et
       \(\beta(\theta)\) ;
 \item on permet de choisir une signature indépendamment de
       \((h,c,\lambda)\) ;
 \item la composition des retenues ne satisfait pas
       \cref{p12-83-c1-s3:eq:tpk-cocadena-transporte} ;
 \item on remplace \(\operatorname{Ext}_r\) par un choix fonctionnel non
       démontré ;
 \item on appelle état complet une projection qui a oublié le cylindre,
       la frontière ou le registre généalogique ;
 \item le retour de phase efface l’une des coordonnées de fibre.
\end{enumerate}
\end{criterion}

Cette unité a défini le conteneur. La suivante énumérera ses lecteurs,
opérateurs, odomètre et calendriers, et démontrera pourquoi les retours de
neuf, vingt-sept et cent huit pas ne sont pas des réinitialisations de l’état.
}%


% Unidad U001. Redacción autónoma; no procede del borrador condensado.

\section{Segments marqués, interstices et cellule nonadique}
\label{p12-83-c1-s4:chap:segmentos-marcas-intersticios}

La construction commence avant toute constante et avant tout choix
d’unité physique. La donnée inaugurale n’est pas un nombre réel déjà constitué, mais un
segment ordonné, une famille de marques et la relation d’adjacence entre
marques consécutives. Cette séparation est nécessaire parce qu’une marque, un
intervalle, la longueur de cet intervalle et le numéral qui publie cette
longueur sont des objets de types différents.

\subsection{Carte positionnelle finie et projection résiduelle décimale}

\begin{definition}[Segment marqué d’ordre dix]
Soit
\[
 \mathsf M_{10}:=\{m_0,m_1,\ldots,m_{10}\}
\]
un ensemble totalement ordonné par
\(m_0<m_1<\cdots<m_{10}\). Son graphe d’adjacence est le chemin
\[
 \mathsf P_{10}:\quad
 m_0\longrightarrow m_1\longrightarrow\cdots\longrightarrow m_{10}.
\]
Nous noterons
\[
 \mathsf I_k:=(m_k,m_{k+1}),\qquad 0\leq k\leq9,
\]
l’interstice abstrait compris entre deux marques consécutives. La
famille des interstices est
\(\mathsf E_{10}:=\{\mathsf I_0,\ldots,\mathsf I_9\}\).
\end{definition}

La notation ouverte \((m_k,m_{k+1})\) ne présuppose pas encore la droite réelle. Elle
indique seulement la paire ordonnée d’extrémités et l’arête qui les relie. Lorsque l’on choisit la
carte archimédienne
\[
 \chi_{10}(m_k)=k,
\]
l’interstice abstrait \(\mathsf I_k\) est publié par l’intervalle
semi-ouvert \([k,k+1)\). Par conséquent, l’indice \(k\), l’arête
\(\mathsf I_k\) et l’intervalle \([k,k+1)\) ne seront pas identifiés dans le texte :
le premier est une étiquette, le deuxième un élément combinatoire et le
troisième sa réalisation dans une carte.

\begin{lemma}[Recensement des marques et des interstices]
\label{p12-83-c1-s4:lem:censo-marcas-intersticios}
Le segment \(\mathsf P_{10}\) contient onze marques et dix interstices. Plus
généralement, un chemin à \(n+1\) marques possède exactement \(n\) arêtes.
\end{lemma}

\begin{proof}
L'application
\[
 \{0,\ldots,n-1\}\longrightarrow E(\mathsf P_n),
 \qquad k\longmapsto(m_k,m_{k+1}),
\]
est bijective. L’affirmation particulière s’obtient avec \(n=10\).
\end{proof}

L’extrémité \(m_0\) fixe l’origine de la carte et \(m_{10}\) fixe son seuil.
Aucune des deux ne constitue par elle-même un interstice supplémentaire. Cette
observation élémentaire empêche trois confusions ultérieures : compter les extrémités
comme des cellules, confondre la clôture avec un nouvel état de phase et traiter le zéro
comme une dixième phase indépendante du cycle nonadique.

\subsection{Extraction de la cellule de \(9\) phases}

La carte décimale distingue l’ancrage \([0,1)\) de la cellule périodique. Les
neuf marques positives
\[
 \mathsf D_9:=\{1,2,3,4,5,6,7,8,9\}
\]
seront utilisées comme représentants publics de
\(\mathbb Z/9\mathbb Z\), en faisant correspondre la marque \(9\) à la classe
résiduelle \(\overline 0\). Nous définissons
\[
 \lambda_9:\mathsf D_9\longrightarrow\mathbb Z/9\mathbb Z,
 \qquad \lambda_9(a)=\overline a.
\]

\begin{proposition}[Carte positive des classes nonadiques]
\label{p12-83-c1-s4:prop:carta-positiva-nonadica}
L’application \(\lambda_9\) est bijective. Son inverse est la section positive
\[
 s_+(\overline a)=
 \begin{cases}
 a,&1\leq a\leq8,\\
 9,&\overline a=\overline0.
 \end{cases}
\]
En particulier, l’écriture publique \(9\) et l’écriture résiduelle \(0\)
représentent la même classe, mais appartiennent à des cartes distinctes.
\end{proposition}

\begin{proof}
Les neuf classes de \(\mathbb Z/9\mathbb Z\) possèdent des représentants uniques dans
\(\{0,1,\ldots,8\}\). Remplacer le représentant \(0\) par \(9\) produit
une autre transversale complète, de sorte que \(\lambda_9\) est bijective et \(s_+\)
est son inverse.
\end{proof}

La cellule de phase ne s’obtient pas en effaçant arbitrairement l’origine. Le tronçon
\([0,1)\) fixe la carte décimale depuis laquelle on observe le cycle ; les neuf
phases sont les neuf classes publiées comme \(1,\ldots,9\). Dans la carte
relevée, le pas qui ferme la cellule est représenté par \([9,10)\) ; dans la
carte résiduelle, ce même pas retourne de \(\overline0\) à \(\overline1\).
La clôture et le changement de représentant ne sont pas la même opération.

\subsection{Position, hauteur et division euclidienne}

\begin{definition}[Coordonnées nonadiques]
Pour chaque entier \(n\geq1\), nous définissons
\[
 \rho_9(n):=1+((n-1)\bmod9),
 \qquad
 q_9(n):=\left\lfloor\frac{n-1}{9}\right\rfloor.
\]
La première coordonnée est appelée position de phase ; la seconde, hauteur de
tour ou quotient de mémoire.
\end{definition}

\begin{theorem}[Décomposition nonadique unique]
\label{p12-83-c1-s4:thm:descomposicion-nonadica-unica}
Tout entier \(n\geq1\) admet une unique écriture
\[
 n=9q+\rho,
 \qquad q\in\mathbb Z_{\geq0},\quad \rho\in\mathsf D_9.
\]
Dans cette écriture, \(q=q_9(n)\) et \(\rho=\rho_9(n)\).
\end{theorem}

\begin{proof}
Nous appliquons la division euclidienne à \(n-1\) : il existe d’uniques
\(q\geq0\) et \(r\in\{0,\ldots,8\}\) tels que
\(n-1=9q+r\). En prenant \(\rho=r+1\), on obtient
\(n=9q+\rho\), avec \(1\leq\rho\leq9\). L’unicité de \(q,r\) implique celle
de \(q,\rho\).
\end{proof}

\begin{corollary}[Retour de phase avec accroissement de mémoire]
\label{p12-83-c1-s4:cor:retorno-fase-incremento-memoria}
Pour tout \(n\geq1\),
\[
 \rho_9(n+9)=\rho_9(n),
 \qquad
 q_9(n+9)=q_9(n)+1.
\]
Par conséquent, un tour complet conserve la phase visible et change
l’état relevé.
\end{corollary}

\begin{proof}
Si \(n=9q+\rho\), alors \(n+9=9(q+1)+\rho\) ; l’unicité du
\cref{p12-83-c1-s4:thm:descomposicion-nonadica-unica} donne les deux identités.
\end{proof}

L’évolution d’un pas s’écrit sans ambiguïté comme
\[
 (q,\rho)\longmapsto
 \begin{cases}
 (q,\rho+1),&1\leq\rho<9,\\
 (q+1,1),&\rho=9.
 \end{cases}
\]
La deuxième ligne n’est pas une réinitialisation. Elle retrouve l’étiquette initiale de phase et
augmente le quotient qui enregistre combien de tours ont été effectués.

\subsection{Périodicité de la projection circulaire et injectivité du relèvement hélicoïdal nonadique}

Soit
\[
 \vartheta_9(\rho)=\frac{2\pi(\rho-1)}9.
\]
La projection de phase sur le cercle unité est
\[
 \mathcal R_9(n)=
 \bigl(\cos\vartheta_9(\rho_9(n)),
       \sin\vartheta_9(\rho_9(n))\bigr).
\]
Pour conserver simultanément la position et la hauteur, nous introduisons
l’immersion hélicoïdale
\[
 \mathcal H_9(n)=
 \left(
  \cos\vartheta_9(\rho_9(n)),
  \sin\vartheta_9(\rho_9(n)),
  q_9(n)+\frac{\rho_9(n)-1}{9}
 \right).
\]

\begin{proposition}[La projection circulaire ne détermine pas l’état]
\label{p12-83-c1-s4:prop:rueda-no-determina-estado}
L’application \(\mathcal R_9\) est périodique de période neuf. L’application
\(\mathcal H_9\) est injective. En particulier,
\[
 \mathcal R_9(n+9)=\mathcal R_9(n),
 \qquad
 \mathcal H_9(n+9)\neq\mathcal H_9(n).
\]
\end{proposition}

\begin{proof}
La première égalité découle du
\cref{p12-83-c1-s4:cor:retorno-fase-incremento-memoria}. La troisième coordonnée de
\(\mathcal H_9(n+9)\) dépasse d’une unité celle de \(\mathcal H_9(n)\),
de sorte que les deux points sont distincts. Enfin, à partir de la troisième coordonnée,
on retrouve \(n\) en multipliant par neuf et en ajoutant un à la partie de phase ;
par conséquent, \(\mathcal H_9\) est injective.
\end{proof}

Cette proposition est le modèle minimal du principe qui traversera tout
l’ouvrage :
\[
 \boxed{\text{retour de phase}\neq\text{retour de l’état}.}
\]
La connexion nonadique du Topological Phase Kernel conservera davantage de champs que
l’hélice élémentaire ---feuillet, orientation, retenue, frontière, terminal et
mémoire de route---, mais elle ne modifiera pas cette distinction inaugurale.

\subsection{Correspondance entre cardinalité discrète et mesure projective}

\begin{definition}[Dénombrement]
Le dénombrement d’un ensemble fini \(X\) est son cardinal \(|X|\). Il n’exige
ni orientation, ni échelle, ni lecteur supplémentaire.
\end{definition}

\begin{definition}[Donnée de mesure]
Une donnée de mesure discrète est un triplet
\[
 (x,\mathcal L,\mathcal C),
\]
où \(x\) est un état, \(\mathcal L\) est un lecteur typé et
\(\mathcal C\) est une carte de publication compatible avec le codomaine de
\(\mathcal L\). Sa sortie visible est
\[
 \mathcal C(\mathcal L(x)).
\]
\end{definition}

Le même ensemble peut avoir un unique cardinal et admettre plusieurs mesures,
parce que la cochaîne locale, l’orientation d’intégration ou la
carte de publication peuvent varier. Inversement, deux états distincts peuvent publier la
même coordonnée. C’est pourquoi l’ouvrage conservera toujours deux niveaux :
\[
 \text{état avec généalogie}
 \quad\longrightarrow\quad
 \text{coordonnée publiée}.
\]

\begin{criterion}[Critères de réfutation de l’unité inaugurale]
La construction est réfutée si l’un quelconque des faits suivants
se produit :
\begin{enumerate}
 \item on compte dix phases indépendantes dans la cellule résiduelle ;
 \item on identifie la marque \(9\) à l’entier \(0\) sans déclarer le
 changement de carte ;
 \item le pas \(9\to1\) force \(q\mapsto0\) ;
 \item deux entiers distincts possèdent la même paire \((q_9,\rho_9)\) ;
 \item une unité physique est attribuée à la cellule avant la définition du
 transducteur dimensionnel.
\end{enumerate}
\end{criterion}

\subsection{Domaine inaugural de la cellule nonadique et exclusion causale de l’APP, du TRIT, du TPK et des constantes analytiques}

L’unité construite jusqu’ici contient uniquement :
\[
 \mathsf P_{10},\quad
 \mathsf D_9,\quad
 \lambda_9,\quad
 (q_9,\rho_9),\quad
 \mathcal R_9,\quad
 \mathcal H_9.
\]
Elle ne contient pas encore l’APP, le TRIT, le Topological Phase Kernel ni une constante
analytique. L’unité suivante prendra deux copies de la cellule nonadique,
construira leur complétion gnomonique et démontrera comment apparaît le support
bidimensionnel sur lequel agissent la matrice multiplicative de l’APP et sa loi
additive interne.


% Unidad U003B. Esqueleto ciclotómico de APP y completación binomial.
% Redacción autónoma a partir de los objetos, no de la arquitectura del PDF abreviado.

\section{Squelette cyclotomique de l’APP et complétion binomiale de l’unité}
\label{p12-83-c1-s5:chap:app-esqueleto-completacion}

Le défaut entre la matrice APP et son accumulateur interne n’est pas un nombre
isolé. Sa loi en
longueur cartésienne, sa décomposition en orbites d’ordre trois et la
complétion de la cellule nonadique produisent trois objets différents qui doivent
rester typés :
\[
 \text{défaut scalaire}\quad D_d,
 \qquad
 \text{module cyclotomique}\quad W_{\rm APP},
 \qquad
 \text{complétion}\quad\overline{\mathcal C}_9.
\]
Le premier compte ; le deuxième conserve l’orientation et le rang ; le troisième ajoute
des strates de frontière. Les égalités entre leurs cardinalités n’autorisent pas à
identifier leurs éléments.

\subsection{Loi du défaut à module fixe et longueur variable}

Pour chaque entier \(d\geq1\), soit
\[
 B_d:=\{1,\ldots,9\}^d.
\]
La racine numérique nonadique est représentée par
\[
 \operatorname{dr}_9(n)=
 \begin{cases}
 n\bmod9,&n\not\equiv0\pmod9,\\
 9,&n\equiv0\pmod9.
 \end{cases}
\]
Nous définissons l’accumulation additive et l’évaluation matricielle multiplicative
sur le même domaine :
\[
 A_d(x_1,\ldots,x_d)
 =\operatorname{dr}_9\!\left(\sum_{j=1}^d x_j\right),
\]
\[
 M_d(x_1,\ldots,x_d)
 =\operatorname{dr}_9\!\left(\prod_{j=1}^d x_j\right).
\]
Leurs masses visibles et leur défaut sont
\[
 S_\Sigma(d):=\sum_{x\in B_d}A_d(x),
 \qquad
 S_\Pi(d):=\sum_{x\in B_d}M_d(x),
\]
\[
 D_d:=S_\Pi(d)-S_\Sigma(d).
\]
Ici, \(d\) compte les positions cartésiennes. Le module \(9\) ne varie pas et
aucune dimension physique n’a encore été introduite.

\begin{lemma}[Masse de l’accumulateur additif]
\label{p12-83-c1-s5:lem:masa-hoja-aditiva-longitud}
Pour tout \(d\geq1\),
\[
 S_\Sigma(d)=45\,9^{d-1}.
\]
\end{lemma}

\begin{proof}
Les \(d-1\) premières coordonnées et un résidu publié
\(r\in\{1,\ldots,9\}\) étant fixés, il existe une unique dernière coordonnée qui donne à la
somme le résidu \(r\). Chaque valeur visible possède ainsi \(9^{d-1}\)
antécédents. En sommant les poids \(1+\cdots+9=45\), on obtient la formule.
\end{proof}

\begin{lemma}[Recensement de la matrice multiplicative]
\label{p12-83-c1-s5:lem:censo-hoja-multiplicativa-longitud}
Pour chacun des six résidus unitaires, il y a \(6^{d-1}\) antécédents
par valeur fixée du produit. Les résidus \(3\) et \(6\) ont
\(d6^{d-1}\) antécédents chacun. Le résidu nul, publié comme \(9\),
a
\[
 9^d-6^d-2d6^{d-1}
\]
antécédents. En conséquence,
\[
 S_\Pi(d)=9^{d+1}-9(d+3)6^{d-1}.
\]
\end{lemma}

\begin{proof}
Le groupe des unités de \(\mathbb Z/9\mathbb Z\) a six éléments. En
fixant \(d-1\) unités, la dernière est déterminée par le produit
unitaire prescrit, ce qui donne \(6^{d-1}\) solutions pour chacune des
six classes.

Pour obtenir un produit de valuation ternaire exactement égale à un, on choisit laquelle
des \(d\) coordonnées appartient à l’une des deux classes \(3,6\), tandis que
les autres sont des unités. Le produit \(3\) ou \(6\) étant fixé, la dernière unité
est déterminée ; d’où \(d6^{d-1}\) pour chaque classe. Tous les tuples
restants ont un produit nul modulo \(9\). Soustraire de \(9^d\) les deux
recensements précédents donne leur multiplicité. La somme pondérée par les
représentants visibles produit la dernière identité.
\end{proof}

\begin{theorem}[Loi dimensionnelle interne de l’APP]
\label{p12-83-c1-s5:thm:ley-dimensional-app-autonoma}
Pour tout \(d\geq1\),
\[
 \boxed{D_d=4\,9^d-9(d+3)6^{d-1}.}
\]
En particulier,
\[
 D_1=0,
 \qquad
 D_2=54,
\]
et
\[
 \sum_{d\geq1}D_dz^d
 =\frac{54z^2(1-3z)}{(1-9z)(1-6z)^2}.
\]
De plus,
\[
 \frac{D_d}{9^d}
 =4-(d+3)\left(\frac23\right)^{d-1}
 \longrightarrow4.
\]
\end{theorem}

\begin{proof}
On soustrait les formules des deux lemmes. Pour la fonction génératrice, on utilise
\[
 \sum_{d\geq1}9^dz^d=\frac{9z}{1-9z},
 \qquad
 \sum_{d\geq1}(d+3)6^{d-1}z^d
 =\frac{z(4-18z)}{(1-6z)^2},
\]
et l’on réduit le numérateur commun. La formule normalisée résulte de la division
par \(9^d\) ; le terme exponentiel multiplié par \(d+3\) converge vers
zéro.
\end{proof}

L’égalité \(D_2=54\) reproduit le recensement \(459-405\), mais la formule
précédente et la loi à module variable
\(\Delta_{3^m}^{\rm pl}=m3^{2m-1}\) sont des énoncés distincts : la première
fait varier \(d\) en maintenant le module \(9\) ; la seconde fait varier le module.

\subsection{Opérateur cyclotomique primaire}

Soit
\[
 G_{A_2}=\begin{pmatrix}2&-1\\-1&2\end{pmatrix},
 \qquad
 C=\begin{pmatrix}0&-1\\1&-1\end{pmatrix}.
\]

\begin{proposition}[Coxeter de type \(A_2\)]
\label{p12-83-c1-s5:prop:coxeter-a2-autonomo}
On vérifie
\[
 C^3=I_2,
 \qquad
 C^{\mathsf T}G_{A_2}C=G_{A_2},
 \qquad
 \det(XI_2-C)=X^2+X+1.
\]
Par conséquent, \(C\) préserve la forme \(A_2\), est d’ordre exactement trois et
ne possède aucun vecteur fixe non nul.
\end{proposition}

\begin{proof}
La multiplication donne
\[
 C^2=\begin{pmatrix}-1&1\\-1&0\end{pmatrix},
 \qquad C^3=I_2.
\]
La substitution directe dans \(C^{\mathsf T}G_{A_2}C\) renvoie
\(G_{A_2}\). Enfin, \(1\) n’est pas racine de \(X^2+X+1\), donc
\(\ker(C-I_2)=0\).
\end{proof}

\subsection{Décomposition de l’APP en \(12\) fibres de type \(A_2\)}

Considérons l’action
\[
 a:\mathbb Z/9\mathbb Z\longrightarrow\mathbb Z/9\mathbb Z,
 \qquad a(r)=4r.
\]
Comme \(4^3\equiv1\pmod9\), sa restriction aux unités a deux orbites
libres :
\[
 \mathcal O_1=(1,4,7),
 \qquad
 \mathcal O_2=(2,8,5).
\]
Pour une orbite \(\mathcal O=(r,4r,7r)\), nous définissons le module d’augmentation
\[
 A(\mathcal O):=
 \ker\!\left(\mathbb Z[\mathcal O]
 \xrightarrow{\,\sum\,}\mathbb Z\right).
\]
Dans la base
\[
 b_1=e_r-e_{4r},
 \qquad b_2=e_{4r}-e_{7r},
\]
la forme euclidienne restreinte a pour matrice de Gram \(G_{A_2}\), et le cycle
\(r\mapsto4r\mapsto7r\mapsto r\) agit par \(C\).

Pour \(x=(x_1,x_2,x_3)\in(\mathbb Z/9\mathbb Z)^3\), les trois paires non
ordonnées sont
\[
 \{1,2\},\quad\{2,3\},\quad\{3,1\}.
\]
Sur chaque paire, on évalue deux lois typées :
\[
 \ell_{ij}^{\Sigma}(x)=x_i+x_j,
 \qquad
 \ell_{ij}^{\Pi}(x)=x_ix_j.
\]
Comme
\[
 \ell_{ij}^{\Sigma}(ax)=a\ell_{ij}^{\Sigma}(x),
 \qquad
 \ell_{ij}^{\Pi}(ax)=a^{-1}\ell_{ij}^{\Pi}(x),
\]
la somme porte l’orientation \(C\) et le produit l’orientation inverse
\(C^{-1}\).

\begin{definition}[Module cyclotomique de l’APP]
\label{p12-83-c1-s5:def:modulo-ciclotomico-app-autonomo}
On définit
\[
 \boxed{
 W_{\rm APP}:=
 \bigoplus_{\{i,j\}}
 \bigoplus_{\mu\in\{\Sigma,\Pi\}}
 \bigoplus_{\mathcal O\in\{\mathcal O_1,\mathcal O_2\}}
 A(\mathcal O).}
\]
Il y a \(3\cdot2\cdot2=12\) termes, chacun de rang deux ; par conséquent
\[
 W_{\rm APP}\cong A_2^{12},
 \qquad \operatorname{rank}W_{\rm APP}=24.
\]
Les étiquettes paire--loi--orbite font partie de l’objet.
\end{definition}

Pour \(u\in\mathbb Z/9\mathbb Z\), nous écrivons
\[
 \beta_{\mathcal O}(u)=
 \begin{cases}
 (1,0),&u=r,\\
 (0,1),&u=4r,\\
 (-1,-1),&u=7r,\\
 (0,0),&u\notin\mathcal O,
 \end{cases}
\]
dans la base \((b_1,b_2)\). L’application d’état est
\[
 \Psi(x)_{ij,\mu,\mathcal O}
 :=\beta_{\mathcal O}(\ell_{ij}^{\mu}(x)).
\]

\begin{theorem}[Équivariance et rang du squelette d’état]
\label{p12-83-c1-s5:thm:esqueleto-estatal-app-autonomo}
Soit \(\rho(a)\) l’action qui utilise \(C\) dans les six fibres additives et
\(C^{-1}\) dans les six fibres multiplicatives. Alors
\[
 \boxed{\Psi(ax)=\rho(a)\Psi(x)}
\]
pour les \(729\) états de \((\mathbb Z/9\mathbb Z)^3\). Le sous-module
rationnel engendré par les images a pour rang \(24\).
\end{theorem}

\begin{proof}
Dans la loi additive, multiplier toutes les coordonnées par \(4\) avance d’une
position dans chaque orbite ; dans la loi multiplicative, cela multiplie la valeur par
\(4^2\equiv7\equiv4^{-1}\pmod9\) et inverse donc le cycle. Cela prouve
l’équivariance composante par composante.

Pour que le rang ne reste pas caché dans une affirmation informatique, nous fixons
l’ordre des coordonnées
\[
 (12,\Sigma,\mathcal O_1),(12,\Sigma,\mathcal O_2),
 (12,\Pi,\mathcal O_1),(12,\Pi,\mathcal O_2),
\]
suivi du même ordre pour les paires \(23\) et \(31\), et, à l’intérieur de chaque
fibre, les deux coordonnées \((b_1,b_2)\). Considérons les états
\[
\begin{gathered}
001,002,004,005,010,011,012,013,014,015,016,020,\\
022,023,040,050,101,102,104,105,110,120,140,150.
\end{gathered}
\]
La matrice entière \(M_\Psi\) dont les lignes sont leurs vingt-quatre images a ses
entrées dans ({-1,0,1}). L’élimination sans fractions de
Bareiss produit comme mineurs principaux successifs
\[
1,1,1,-1,-1,-1,-1,1,-1,1,1,1,1,-2,-1,1,1,1,2,1,4,-6,12,9.
\]
En particulier,
\[
 \det M_\Psi=9\neq0.
\]
Ces vingt-quatre images sont linéairement indépendantes. Comme le
codomaine a déjà pour rang \(24\), elles engendrent \(W_{\rm APP}\otimes\mathbb Q\).
\end{proof}

La preuve présente un certificat fini reproductible ; elle ne déduit pas le rang plein
de la seule surjectivité de chaque projection composante.

\subsection{Localisation de l’agrégat \texorpdfstring{\(54\)}{54}}

Soit
\[
 \mathcal M:=\mathbb Z[\mathbb Z/9\mathbb Z]
\]
et, pour \(\mu\in\{\Sigma,\Pi\}\),
\[
 \nu_\mu:=\sum_{x,y\in\mathbb Z/9\mathbb Z}
 e_{\ell^\mu(x,y)}.
\]
Le recensement des deux tables donne
\[
 \delta_{\Pi\Sigma}:=\nu_\Pi-\nu_\Sigma
 =12e_0+3e_3+3e_6
 -3\sum_{u\in(\mathbb Z/9\mathbb Z)^\times}e_u.
\]
Ses coefficients sont constants sur les orbites de \(a\) ; par conséquent
\[
 (I-a)\delta_{\Pi\Sigma}=0.
\]
Nous définissons la projection d’augmentation
\[
 T:\mathcal M\longrightarrow
 A(\mathcal O_1)\oplus A(\mathcal O_2),
\]
\[
 T(\nu)=
 \left(((I-a)\nu)|_{\mathcal O_1},
       ((I-a)\nu)|_{\mathcal O_2}\right).
\]
Alors
\[
 T(\delta_{\Pi\Sigma})=0.
\]
Cependant, le lecteur visible
\[
 w(e_0)=9,
 \qquad w(e_r)=r\quad(1\leq r\leq8)
\]
satisfait
\[
 w(\delta_{\Pi\Sigma})=54.
\]

\begin{corollary}[Obstruction isotypique]
\label{p12-83-c1-s5:cor:obstruccion-isotipica-app-autonoma}
Le scalaire \(54\) appartient au secteur \(C_3\)-invariant et ne peut
être retrouvé par une fonctionnelle linéaire qui se factorise exclusivement par
les modules d’augmentation. En particulier,
\[
 \sum_x\Psi(x)=0
\]
et, pour toute fonctionnelle linéaire \(L\) sur \(W_{\rm APP}\otimes\mathbb Q\),
\[
 L\!\left(\sum_x\Psi(x)\right)=0\neq54.
\]
\end{corollary}

Cette obstruction empêche de présenter \(54\) comme une somme linéaire rétrospective
des douze fibres. Le couplage APP--Fock ultérieur conserve
l’orientation et utilise un objet supplémentaire ; il ne contredit pas le corollaire.

\subsection{Complétion cubique de la cellule nonadique}

Soit \(E_9\) l’ensemble sous-jacent à \(\mathbb Z/9\mathbb Z\), et soit
\(\star\notin E_9\) un état de frontière. Nous définissons
\[
 \mathcal C_9:=E_9^3,
 \qquad
 \overline{\mathcal C}_9:=(E_9\sqcup\{\star\})^3.
\]
Le symbole \(\star\) n’est pas la classe zéro : il indique qu’une coordonnée a
atteint la strate de clôture.

\begin{theorem}[Stratification binomiale de l’unité cubique]
\label{p12-83-c1-s5:thm:emrh-binomial-autonoma}
La strate ayant exactement \(r\) coordonnées égales à \(\star\) contient
\[
 \binom3r9^{3-r}
\]
éléments. Par conséquent,
\[
 \boxed{1000=729+243+27+1=729+271.}
\]
En retirant seulement le sommet \((\star,\star,\star)\), il reste
\[
 999=729+243+27=729+270.
\]
\end{theorem}

\begin{proof}
On choisit les \(r\) coordonnées de frontière de \(\binom3r\) manières ; chaque
coordonnée restante admet neuf valeurs. Les quatre strates sont disjointes
et épuisent le produit cartésien. La somme est le développement de
\((9+1)^3\).
\end{proof}

\begin{figure}[htbp]
\centering
\resizebox{0.96\linewidth}{!}{\input{figuras/base_autoridad/fig_cubo_emrh_autoridad}}
\caption{Complétion cubique de la cellule nonadique. L’intérieur de
cardinal \(729\) et les strates de frontière de cardinalités \(243\), \(27\) et
\(1\) sont des parties disjointes d’un unique objet.}
\label{p12-83-c1-s5:fig:cubo-emrh-autonomo}
\end{figure}

Cette stratification reçoit dans HMT le nom d’extrême et moyenne raison
holographique. Mathématiquement, c’est la complétion binomiale typée de la cellule
nonadique ; le nom ne remplace pas sa définition.

Pour tout \(d\geq1\), soit
\[
 \overline{\mathcal C}^{(d)}_9=(E_9\sqcup\{\star\})^d.
\]
Attribuons le poids \(1/10\) à chaque symbole d’une coordonnée et prenons la mesure
produit \(\mu_d\).

\begin{theorem}[Compatibilité projective de la mesure]
\label{p12-83-c1-s5:thm:medida-proyectiva-completacion}
La projection qui oublie la dernière coordonnée satisfait
\[
 (\pi_d)_*\mu_d=\mu_{d-1}.
\]
De plus,
\[
 10^d=\sum_{r=0}^{d}\binom dr9^{d-r},
 \qquad
 \mu_d(\overline{\mathcal C}^{(d)}_9)=1.
\]
Le nombre d’états augmente avec \(d\), tandis que la masse totale normalisée
est conservée.
\end{theorem}

\begin{proof}
Pour tout cylindre \(A\times(E_9\sqcup\{\star\})\),
\[
 \mu_d(A\times(E_9\sqcup\{\star\}))
 =\mu_{d-1}(A)\sum_{u\in E_9\sqcup\{\star\}}\frac1{10}
 =\mu_{d-1}(A).
\]
Les cylindres engendrent l’algèbre finie des sous-ensembles, ce qui détermine la
mesure image. L’identité de cardinalités est à nouveau le binôme de Newton.
\end{proof}

\subsection{Morphisme du bloc central de l’APP vers sa réalisation euclidienne}

Le bloc central déjà construit est
\[
 B_c=\begin{pmatrix}7&2\\2&7\end{pmatrix}.
\]
La lecture par lignes et la lecture diagonale donnent
\[
 72+27=77+22=99,
\]
\[
 72-27=45=\det B_c,
 \qquad
 77-22=55=54+1.
\]
La relation avec la complétion ne s’obtient pas en concaténant des symboles, mais
par la division euclidienne
\[
 729=10\cdot72+9,
 \qquad
 271=10\cdot27+1.
\]

\begin{theorem}[Unicité du relèvement décimal]
\label{p12-83-c1-s5:thm:levantamiento-decimal-emrh-autonomo}
Pour \(L_d(n)=10n+d\), avec \(d\in\{0,\ldots,9\}\), il existe une unique paire
\((d,d')\) telle que
\[
 L_d(72)=729,
 \qquad
 L_d(72)+L_{d'}(27)=1000,
 \qquad
 d+d'=10.
\]
C’est la paire
\[
 (d,d')=(9,1).
\]
\end{theorem}

\begin{proof}
La première égalité est \(720+d=729\), donc \(d=9\). La deuxième devient
\(729+270+d'=1000\), d’où \(d'=1\). La troisième est vérifiée et
l’unicité du quotient et du reste exclut une autre solution.
\end{proof}

On obtient ainsi le diagramme causal
\[
 B_c\longrightarrow(72,27)
 \longleftarrow((72,9),(27,1))
 \longleftarrow(729,271).
\]
Le bloc et sa lecture appartiennent à l’APP ; les dividendes et les restes appartiennent
à la complétion. Le diagramme partage des quotients, il n’inverse pas la
causalité.

\subsection{Incidence du cube et séparation d’avec les hexades et les octades}

Soit \(M_d\) la matrice d’incidence des \(2d\) faces de codimension un
d’un hypercube sur ses \(2^d\) sommets. Chaque face contient
\(2^{d-1}\) sommets ; deux faces opposées sont disjointes et deux faces
transversales partagent \(2^{d-2}\) sommets.

\begin{theorem}[Spectre d’incidence hypercubique]
\label{p12-83-c1-s5:thm:espectro-incidencia-cubo-autonomo}
\[
 \operatorname{Spec}(M_dM_d^{\mathsf T})
 =\left\{
 d2^{d-1}\,^{(1)},
 2^{d-1}\,^{(d)},
 0^{(d-1)}
 \right\}.
\]
Para \(d=3\),
\[
 \operatorname{Spec}(M_3M_3^{\mathsf T})
 =\{12,4,4,4,0,0\}.
\]
\end{theorem}

\begin{proof}
L’espace des lignes se décompose en : la droite constante ; les \(d\)
différences de chaque paire de faces opposées ; et les (d-1) combinaisons de
sommes de paires dont la somme totale est nulle. Les intersections décrites précédemment
font que la matrice de Gram agit par \(d2^{d-1}\), \(2^{d-1}\) et \(0\),
respectivement.
\end{proof}

Le cube tridimensionnel possède six faces, huit sommets et douze arêtes. Ces
cardinalités anticipent un motif d’incidence (6/8), mais une face n’est pas
une hexade et l’ensemble des sommets n’est pas une octade. L’identification
exceptionnelle exige un transporteur ultérieur qui conserve les supports et la
jauge.

\subsection{Résonance de l’opérateur d’incidence en rang \(6\)}

Définissons les strates qui conservent exactement \(k\) coordonnées
nonadiques en rang \(d\) :
\[
 \mathcal I_k(d):=\binom dk9^k.
\]
Les défauts de l’APP sont
\[
 D_{\rm vis}=54,
 \qquad
 D_{\rm tot}=1215,
 \qquad
 D_{\rm coc}=129,
 \qquad
 1215=54+9\cdot129.
\]

\begin{theorem}[Unicité de la résonance binomiale]
\label{p12-83-c1-s5:thm:resonancia-binomial-autonoma}
L’unique entier positif \(d\) qui satisfait simultanément
\[
 \mathcal I_1(d)=54,
 \qquad
 \mathcal I_2(d)=1215
\]
est \(d=6\). Dans ce rang,
\[
 \frac{\mathcal I_2(6)-\mathcal I_1(6)}9=129.
\]
\end{theorem}

\begin{proof}
La première égalité est \(9d=54\), qui force \(d=6\). Alors
\[
 \mathcal I_2(6)=81\binom62=81\cdot15=1215,
\]
et
\[
 \frac{1215-54}{9}=129.
\]
Il ne reste aucun autre rang possible parce que la première égalité détermine déjà \(d\).
\end{proof}

\subsection{Obstructions du squelette cyclotomique de l’APP : fibres \(A_2^{12}\), complétion \(729+271\), résonance de rang \(6\) et incidence cubique}

\begin{criterion}[Réfutation typée de l’unité]
La construction est réfutée si l’un quelconque des faits
suivants se produit :
\begin{enumerate}
 \item la loi de \(D_d\) ne reproduit pas \(D_1=0\) et \(D_2=54\) ;
 \item la variation de longueur est fusionnée avec la variation de module ;
 \item \(C\) cesse de préserver \(G_{A_2}\) ou acquiert un vecteur fixe ;
 \item les étiquettes paire--feuillet--orbite ne produisent pas douze fibres ;
 \item le mineur certifié de \(\Psi\) a un déterminant nul ;
 \item une fonctionnelle qui se factorise par le module d’augmentation retrouve \(54\) ;
 \item le symbole de frontière \(\star\) est identifié au résidu zéro ;
 \item la somme des quatre strates n’est pas égale à \(1000\) ;
 \item une face cubique est identifiée à une hexade sans un transporteur
d’incidence.
\end{enumerate}
\end{criterion}

L’unité a construit, sans les reconnaître rétrospectivement, le module
\(A_2^{12}\), la complétion \(729+271\), la mesure projective et la
résonance de rang six. Ces structures alimenteront de manière différente
la génération des constantes et les branches exceptionnelles.


% Unidad U006F. Inventario operatorio completo del Topological Phase Kernel.

\long\def\HMTDeferredTPKDevelopment{%
\section{Algèbre des opérateurs du Topological Phase Kernel et dynamique causale d’actualisation de l’état enrichi}
\label{p12-83-c1-s6:chap:inventario-operatorio-completo-tpk}

Le Topological Phase Kernel n’est ni une réduction modulo trois ni une horloge à
cent huit positions. C’est une catégorie d’états enrichis équipée
d’opérateurs de sélection, de transport, de lecture, de mémoire, de périodicité,
de prolongement et de publication. Pour empêcher une abréviation de remplacer
l’objet, chaque opérateur est enregistré selon son type et l’information qu’il
conserve.

\subsection{Schéma typé des 34 constructions canoniques du TPK : domaine, codomaine, action, invariants et mémoire}

\begin{definition}[Fiche opératoire]
\label{p12-83-c1-s6:def:u006f-ficha-operatoria}
Une réalisation canonique du TPK est un sextuplet
\begin{equation}
 \mathcal O=
 (D_{\mathcal O},C_{\mathcal O},f_{\mathcal O},
 I_{\mathcal O},L_{\mathcal O},M_{\mathcal O}),
 \label{p12-83-c1-s6:eq:u006f-ficha}
\end{equation}
où :
\begin{enumerate}
 \item \(f_{\mathcal O}:D_{\mathcal O}\to C_{\mathcal O}\) est sa règle ;
 \item \(I_{\mathcal O}\) est la famille des relations conservées ;
 \item \(L_{\mathcal O}\) déclare l’information publiée ou perdue ;
 \item \(M_{\mathcal O}\) énumère la mémoire qui demeure dans l’état
 enrichi.
\end{enumerate}
\end{definition}

Deux réalisations ne représentent le même opérateur que s’il existe une
équivalence typée qui entrelace les règles et les invariants. Partager une période,
une cardinalité ou une sortie visible ne suffit pas.

\subsection{Classification de l’état enrichi en \(8\) classes structurelles invariantes}

L’inventaire est partitionné en :
\begin{enumerate}
 \item[\textbf{C1.}] supports et états ;
 \item[\textbf{C2.}] sélection interne ;
 \item[\textbf{C3.}] transporte;
 \item[\textbf{C4.}] lecteurs et quotients ;
 \item[\textbf{C5.}] mémoire et frontière ;
 \item[\textbf{C6.}] périodicité et retour ;
 \item[\textbf{C7.}] prolongement ;
 \item[\textbf{C8.}] salidas.
\end{enumerate}
La classe est déterminée par le domaine, le codomaine et la fonction, non par le nom
d’un chiffre associé.

\Needspace{.92\textheight}
\subsection{Décomposition fonctionnelle en \(14\) regroupements compatibles avec les opérateurs APP–TRIT–TPK}

\begingroup
\small
\renewcommand{\arraystretch}{1.12}
\begin{longtable}{@{}>{\raggedright\arraybackslash}p{.20\textwidth}>{\raggedright\arraybackslash}p{.29\textwidth}>{\raggedright\arraybackslash}p{.43\textwidth}@{}}
\toprule
Regroupement&Opérateurs ou espaces&Fonction et invariants\\
\midrule
\endfirsthead
\toprule
Regroupement&Opérateurs ou espaces&Fonction et invariants\\
\midrule
\endhead
État et mémoire de route&
\(\mathcal X_{\rm TPK}^{\rm enr}\), \(\mathcal F_q\)&
Ils conservent la position, le feuillet, l’orientation, le résidu, le quotient, la retenue, la signature,
la frontière, le cylindre, le préfixe, la provenance et la mémoire.\\

Sélection interne&
\(M_{\rm ph}:\{1,\ldots,9\}\to\{-1,0,+1\}\)&
Active l’évaluation additive, multiplicative ou le pas de seuil ; ne déplace pas
l’état.\\

Transporte espacial&
\(T_N,T_E,T_S,T_O,F_{\rm loc},F_{\rm obs}\)&
Réalise des translations cardinales, une rétroaction de feuillet et une chronologie à deux
curseurs.\\

Lectores modulares&
\(\rho_9,q_9,\rho_3^{\rm or},c_3,q_{1000},r_{1000}\)&
Publient la phase, l’orientation et le bloc décimal en conservant les quotients de
relèvement.\\

Bloc et mémoire&
\(\mathcal K_{27}=F_{\rm obs}^{27}\), \(\mathcal N_{27}\)&
Le premier compose vingt-sept transports ; le second filtre des événements de
mémoire. Un indice identique n’implique pas un opérateur identique.\\

État dodécaphasique&
\(\mathcal R_{12},C_{12},\Theta_{108},\Gamma_{108}\)&
Distingue le registre enrichi, le secteur, la loi cyclique et le support ordonné.\\

Périodicité et profondeurs&
\(C_{12}\hookrightarrow C_{108}\twoheadrightarrow C_9\),
\(w_{108},w_{1080}\)&
Conserve l’extension non scindée, les mots de longueurs différentes et
la mémoire verticale.\\

Canaux \(90/120\)&
\(\mathcal F_6,\mathcal F_8,\mathcal I_{6\to8},
S_{90},S_{120}\)&
Sépare les incidences finies, les queues analytiques et leurs normalisations.\\

Bidirectionnalité&
\(P_\pm,\operatorname{rev},\operatorname{Ext}_\pm,
N_\pm,\beta,C_M\)&
Réalise l’échange, l’inversion de route, le prolongement futur et passé,
la valence, le bilan et la conjugaison.\\

Émission et coinduction&
\(L_0,L_1,R_{36},G_9,\varprojlim\operatorname{Cyl}_m\)&
Relève le germe, ouvre la frontière et organise des prolongements finis et
compatibles.\\

Lecteurs numériques&
\(\mathscr C_\pi,\mathscr P_e,F_{\rm av}\)&
Construisent des encadrements rationnels, une propagation factorielle et une auto-échelle de
Fibonacci.\\

Projection exceptionnelle&
\(\Gamma_{A_W}\), poids, support projectif&
Transporte la même émission vers une incidence ternaire sans sélectionner de chiffres
archimédiens.\\

Atlas et réalisation&
signature de route, atlas \(81/56/13/324\), \(C_M\)&
Classe les histoires et les caractères avant d’appliquer une carte dimensionnelle.\\

Validation finie&
recensements, preuves de nécessité structurelle, identités d’intégrité et
contrôles d’indépendance&
Vérifie la couverture à la profondeur exécutée et sépare la construction,
la vérification indépendante et les données de confrontation.\\
\bottomrule
\end{longtable}
\endgroup

Les huit classes décrivent la nature mathématique ; les quatorze regroupements,
la fonction d’exécution. Ce sont deux gradations du même inventaire.

\Needspace{.92\textheight}
\subsection{Domaine et typage des \(34\) entrées de l’opérateur d’actualisation}

\begin{equation}
 \mathfrak O_{\rm TPK}:=(\mathfrak o_1,\ldots,\mathfrak o_{34}).
 \label{p12-83-c1-s6:eq:u006f-registro-operadores}
\end{equation}

\begingroup
\scriptsize
\renewcommand{\arraystretch}{1.12}
\begin{longtable}{@{}r p{.19\textwidth}p{.28\textwidth}p{.42\textwidth}@{}}
\toprule
\#&Symbole&Domaine et codomaine&Règle ou fonction\\
\midrule
\endfirsthead
\toprule
\#&Symbole&Domaine et codomaine&Règle ou fonction\\
\midrule
\endhead
1&\(\mathcal A_9\)&\((\mathbb Z/9)^2\), matrice et loi interne&
Cellule arithmétique additive--multiplicative.\\
2&\((\rho_9,q_9)\)&\(\mathbb Z_{>0}\to
\{1,\ldots,9\}\times\mathbb Z_{\geq0}\)&
Division nonadique avec représentant positif et quotient.\\
3&\(\rho_3^{\rm or},c_3\)&\(\mathbb Z\to
\{-1,0,+1\}\times\mathbb Z\)&
Relèvement ternaire orienté \(n=\rho_3^{\rm or}(n)+3c_3(n)\).\\
4&\(\iota_{1000}:=(q_{1000},r_{1000})\)&\(\mathbb Z\to
\mathbb Z\times\{0,\ldots,999\}\)&
Division euclidienne \(N=1000q_{1000}+r_{1000}\).\\
5&\(T_d\)&\(X_9\to X_9\), \(d\in\{N,E,S,O\}\)&
Translation cardinale orientée.\\
6&\(F_{\rm loc}\)&\(\mathcal Q_{\rm loc}\times
\mathscr D_{\rm card}\to\mathcal Q_{\rm loc}\)&
Rétroaction locale de position et de feuillet.\\
7&\(T_{\min}\)&\(\Omega_{\rm seed}\to\mathcal U_6\)&
Émetteur minimal des deux curseurs, y compris la signature coordonnée.\\
8&\(G\)&\(\operatorname{im}s_\Sigma\times
\operatorname{im}s_\Pi\to\{0,\ldots,999\}^6\)&
Couplage décimal coordonnée par coordonnée.\\
9&\(\Psi_6\)&\(\mathcal U_6\to\mathbb F_3^6\)&
Réduction ternaire du catalogue.\\
10&\(\mathcal R_3\)&\(\mathcal Q_{\rm loc}\times
\mathscr D_{\rm card}^3\to\{-1,0,+1\}\)&
Facteur observable local de trois pas.\\
11&\(\mathcal K_{27}\)&\(\mathcal Q_{\rm obs}\to\mathcal Q_{\rm obs}\)&
Composition \(F_{\rm obs}^{27}\).\\
12&\(\mathcal X_{\rm TPK}^{\rm enr}\)&catégorie des états enrichis&
Conserve toutes les coordonnées typées du noyau.\\
13&\(\mathcal H\)&relation sur les états enrichis&
Transition de Hensel qui retient le résidu et le quotient.\\
14&\(\mathcal C_{3,1000}\)&états croisés
\((N,L,D,K,A,B)\)&
Actualisation exacte des cylindres ternaires--décimaux.\\
15&\(\Sigma_{\rm B}\)&mémoires de frontière \(\to\) signature conjointe&
Publie l’orientation de bloc en conservant son registre généalogique.\\
16&\({\rm EF}\text{--}G_9\)&relation d’extension&
Filtre les extensions compatibles et organise la généalogie nonadique.\\
17&\(\Gamma_9\)&histoires enrichies \(\longrightarrow\) histoires&
Holonomie d’un tour de phase avec mémoire de frontière.\\
18&\(X_\chi\)&\(\varprojlim_m\operatorname{Surv}^{(\chi)}_m\)&
Section globale lorsque la non-vacuité, la compatibilité et l’unicité sont
démontrées à toute profondeur.\\
19&\(\mathcal P_{\rm APP}^{(2)}\)&catégorie bivariée de chemins&
Transporte simultanément l’évaluation multiplicative et l’accumulation
additive.\\
20&\(F_{\rm obs}\)&\(\mathcal Q_{\rm obs}\to\mathcal Q_{\rm obs}\)&
Chronologie observable déterministe à deux curseurs.\\
21&\(\mathcal A_{42}\)&alphabet de quarante-deux triades&
Support décimal élargi de l’émission finie.\\
22&\(\operatorname{Ext}_r\)&
\(\operatorname{Ext}_r\subseteq\mathcal F_q\times\mathcal F_{q'}\)&
Relation typée de prolongement des fibres.\\
23&\(\mathcal K_{\rm ph}\)&\(\mathcal U_6\times C_9\) sur lui-même&
Transfert de continuations compatibles avec phase explicite.\\
24&\((P,H,Q)\)&\(\mathcal X_\times\to\mathcal X_\times\)&
Avancée, demi-tour et quart de tour des germes.\\
25&\((c_{\rm mode},c_{\rm or})\)&caminos \(\to\mathbb Z^2\)&
Cocycles de changement de mode et d’orientation.\\
26&\(E_{20}\)&états finis \(\to\) vingt blocs&
Prolongement triadique fini avec livre de provenance.\\
27&\(\mathcal R_{12}\)&histoire enrichie \(\to\) registre de douze
secteurs&
Système temporel orienté avec route, signature, retenue et registre généalogique.\\
28&\(R_{\rm sgn}\)&état terminal \(\to U_{12}^{\rm sgn}
\subseteq\mathbb Z^{12}\)&
Lecteur de l’observable harmonique signée.\\
29&\(\mathcal H_{12}\)&\(\mathbb Z^{12}\to\mathbb Z^{12}\)&
Transformation par trois orbites de pas trois ; inverse
\(\mathcal H_{12}/4\) sur l’image pertinente.\\
30&\(\mathcal E_{\rm dod}=(D_3,D_4,Q)\)&
\(\mathbb Z^{12}\to\mathbb Z^{25}\)&
Extracteur transversal conjoint d’arités trois et quatre.\\
31&\(\mathcal C_\alpha\)&blocs et retenues \(\to\) blocs&
Récurrence dodécaphasique de clôture en base mille.\\
32&\((L_{\rm tick},\chi_{\rm mass})\)&histoire enrichie \(\to\)
caractère temporel&
Interface dimensionnelle ultérieure ; elle ne participe pas au constructeur numérique de
cet ouvrage.\\
33&\(\mathcal W_\pm\)&réseau bidirectionnel d’indice douze&
Lecture dans les deux sens en conservant l’orientation et la provenance.\\
34&\(\mathcal A_{\rm species}\)&familles d’extension \(\to\) atlas&
Interface de réalisation ultérieure, non un quatrième primitif.\\
\bottomrule
\end{longtable}
\endgroup

\subsection{Composition causale de l’actualisation de l’état enrichi}

\begin{definition}[État canonique complet]
\label{p12-83-c1-s6:def:u006f-esquema-estado}
Soit
\[
 \mathcal E:=
 \mathbb Z_{\rm quo}\times\mathbb Z_{\rm car}\times
 \mathsf{Mem}\times\mathsf{Term}\times\mathsf{Cyl}\times
 \mathsf{Fr}\times\mathsf{Pref}\times\Lambda .
\]
L’état enrichi canonique à l’instant \(t\) est
\begin{equation}
 x_t=(p_t,m_t,\sigma_t,b_t,r_t;
 q_t,c_t,\eta_t,s_t,C_t,\partial_t,N_t,\lambda_t)
 \in
 X_9\times\{\Sigma,\Pi\}\times\mathsf{TRIT}\times
 C_{12}\times C_9\times\mathcal E.
 \label{p12-83-c1-s6:eq:u006f-esquema-estado}
\end{equation}
Les abréviations d’état utilisées dans d’autres unités sont des projections de
ce tuple. En particulier, ni le quotient \(q_t\), ni la retenue
\(c_t\), ni le terminal \(s_t\), ni le cylindre \(C_t\), ni la frontière
\(\partial_t\), ni le préfixe \(N_t\), ni le livre \(\lambda_t\) ne sont omis.
\end{definition}

Soit \(\rho_t\) la flèche observable déterminée par \(x_t\). Les lois de
transport de U005 et la transition \(\mathcal C_{3,1000}\) déterminent un
endomorphisme de mémoire
\begin{align}
 \mathbf E_t:\mathcal E&\longrightarrow\mathcal E,\notag\\
 (q,c,\eta,s,C,\partial,N,\lambda)&\longmapsto
 \left(
 \begin{aligned}
 &\mathsf H(\rho_t)q,
  \mathsf C(\rho_t)c+\kappa(\rho_t),
  \mathsf M(\rho_t)\eta,\\
 &\mathsf T(\rho_t)s,
  \mathsf{Cyl}(\rho_t)C,
  \mathsf{Fr}(\rho_t)\partial,\\
 &\mathsf{Pref}(\rho_t)N,
  \Lambda(\rho_t)\lambda
 \end{aligned}
 \right).
 \label{p12-83-c1-s6:eq:u006f-actualizacion-memoria}
\end{align}
Chaque composante publiée dans \eqref{p12-83-c1-s6:eq:u006f-actualizacion-memoria} possède le
domaine indiqué par sa fibre ; \(\kappa(\rho_t)\) est la cochaîne de retenue.

Nous définissons trois endomorphismes du produit canonique. Si
\(\sigma'=M_{\rm ph}(g(t))\) et \(m'=\mu(\sigma',m)\), alors
\begin{align}
 \mathsf{Sel}_t(p,m,\sigma,b,r;e)
  &:=(p,m',\sigma',b,r;e),\label{p12-83-c1-s6:eq:u006f-selector-tipado}\\
 \mathsf{Tra}_t(p,m,\sigma,b,r;e)
  &:=(T_{d(t)}p,m,\sigma,b,r;e),\label{p12-83-c1-s6:eq:u006f-transporte-tipado}\\
 \mathsf{Upd}_t(p,m,\sigma,b,r;e)
  &:=(p,m,\sigma,b',r';\mathbf E_t(e)),
 \label{p12-83-c1-s6:eq:u006f-memoria-tipada}
\end{align}
où \((b',r')\) est la destination de \(\rho_t:(b,r)\to(b',r')\). L’actualisation
causale est la composition d’applications, non la composition d’une application avec une valeur
tritique :
\begin{equation}
 \boxed{\mathcal U_t:=
 \mathsf{Upd}_t\circ\mathsf{Tra}_t\circ\mathsf{Sel}_t.}
 \label{p12-83-c1-s6:eq:u006f-tuberia}
\end{equation}

\begin{theorem}[Clôture de l’actualisation causale]
\label{p12-83-c1-s6:thm:u006f-cierre-tuberia}
L’application \(\mathcal U_t\) est un endomorphisme de l’état canonique
complet. Les lecteurs, prolongements et projections sont appliqués après
\(\mathcal U_t\) uniquement lorsque leur domaine a été produit.
\end{theorem}

\begin{proof}
\(\mathsf{Sel}_t\) ne change que le mode et le TRIT ;
\(\mathsf{Tra}_t\) ne change que la position APP ; et
\(\mathsf{Upd}_t\) applique le produit typé
\eqref{p12-83-c1-s6:eq:u006f-actualizacion-memoria} et la flèche de phase. Les trois
codomaines sont le même produit de la
\cref{p12-83-c1-s6:def:u006f-esquema-estado} ; leur composition est donc définie et
conserve toutes les coordonnées de l’état.
\end{proof}

\subsection{Séparation typée entre opérateurs de transport, de mémoire, de calendrier, d’extraction et de retour}

L’inventaire impose
\begin{equation}
 \mathcal S_4\neq D_4\neq\operatorname{sd}_4,
 \qquad
 \mathcal R_{12}\neq C_{12}\neq\Theta_{108}\neq\Gamma_{108},
 \qquad
 F_{\rm obs}\neq\mathcal K_{\rm ph},
 \label{p12-83-c1-s6:eq:u006f-no-identificacion-uno}
\end{equation}
et
\begin{equation}
 U_{12}^{\rm cat}\neq U_{12}^{\rm sgn},
 \qquad
 \mathcal E_{\rm dod}\neq D_{\rm raw}^{90,120}
 \neq\mathcal K_{6\to8}.
 \label{p12-83-c1-s6:eq:u006f-no-identificacion-dos}
\end{equation}
Ce sont des incompatibilités de type, non de simples inégalités de valeur.

\begin{criterion}[Audit de complétude opératorielle]
\label{p12-83-c1-s6:crit:u006f-completitud}
Une exécution n’est opératoriellement complète que si :
\begin{enumerate}
 \item elle déclare les entrées de \(\mathfrak O_{\rm TPK}\) employées ;
 \item elle enregistre le domaine et le codomaine de chaque composition ;
 \item elle conserve l’ordre causal de \eqref{p12-83-c1-s6:eq:u006f-tuberia} ;
 \item elle distingue dynamique, transfert, lecture, prolongement et sortie ;
 \item elle permet de reconstruire à partir de \(\lambda_t\) chaque changement de bloc et
 chaque retenue ;
 \item elle ne promeut pas une vérification finie en une section infinie sans
 prouver la non-vacuité, la compatibilité et l’unicité à toute profondeur.
\end{enumerate}
Tout symbole utilisé sans type ou toute fusion interdite par
\eqref{p12-83-c1-s6:eq:u006f-no-identificacion-uno} et
\eqref{p12-83-c1-s6:eq:u006f-no-identificacion-dos} constitue un échec vérifiable.
\end{criterion}


% Unidad U006E. Separación tipada de los canales 90/120 y de los símbolos
% epsilon. La realización de incidencia se construye después de W12 y W24.

\section{Décomposition opératorielle des canaux d’incidence \texorpdfstring{\(90/120\)}{90/120} : sélecteur temporel, générateur nilpotent et composante angulaire}
\label{p12-83-c1-s7:chap:canales-90-120-epsilon}

Les entiers \(90\) et \(120\) apparaissent dans plusieurs strates de la construction. Leur
coïncidence numérique n’établit pas une identité d’objets. De même,
la lettre epsilon désigne trois constructions aux domaines incompatibles. La
finalité de cette unité est de fixer leurs types avant de les employer et d’indiquer
avec précision quelles réalisations exigent des structures qui n’ont pas encore été
construites.

\subsection{Opérateurs temporel, nilpotent et angulaire : domaines et codomaines}

\subsubsection{Action du sélecteur temporel tritique sur \(\{-1,0,+1\}\) : ouverture, seuil et retour avec mémoire}

Le sélecteur interne du calendrier est
\begin{equation}
 \varepsilon_9:\{1,\ldots,9\}
 \longrightarrow\{-1,0,+1\},
 \qquad
 \varepsilon_9(r)=
 \begin{cases}
  +1,&r\in\{1,4,7\},\\
  -1,&r\in\{2,5,8\},\\
  0,&r\in\{3,6,9\}.
 \end{cases}
 \label{p12-83-c1-s7:eq:u006e-epsilon-temporal}
\end{equation}
Son codomaine est le TRIT orienté. Il décide quel feuillet agit et ne déplace
par lui-même aucune position.

\subsubsection{Générateur parabolique nilpotent \(J_0^2=0\) et évolution au seuil temporel}

La classe
\begin{equation}
 \varepsilon_{\rm nil}
 \in
 \mathcal Q_{\rm n}:=
 \mathbb F_3[\varepsilon_{\rm nil}]
 /(\varepsilon_{\rm nil}^{\,2})
 \label{p12-83-c1-s7:eq:u006e-epsilon-nilpotente}
\end{equation}
satisfait
\begin{equation}
 \varepsilon_{\rm nil}\neq0,
 \qquad
 \varepsilon_{\rm nil}^{\,2}=0.
 \label{p12-83-c1-s7:eq:u006e-epsilon-nilpotente-relacion}
\end{equation}
C’est un élément d’une algèbre quadratique dégénérée ; ce n’est pas un sélecteur
temporel.

\subsubsection{Résidu angulaire induit après la sélection temporelle et conservation de l’orientation}

Le symbole
\begin{equation}
 \varepsilon_{\rm res}^{\circ}\in\mathbb R
 \label{p12-83-c1-s7:eq:u006e-epsilon-residual-tipo}
\end{equation}
désigne, lorsque ses paramètres géométriques seront introduits, le résidu d’une
identité angulaire. À cette étape, seul son type est réservé. Il n’est obtenu
ni à partir de \(\varepsilon_9\) ni à partir de
\(\varepsilon_{\rm nil}\), et n’intervient pas dans la construction de
\(\pi,e,\varphi\) ou de \(\alpha\).

\begin{proposition}[Non-identification des symboles epsilon]
\label{p12-83-c1-s7:prop:u006e-no-identificacion-epsilon}
Il n’existe pas d’égalité bien typée entre
\[
 \varepsilon_9,\qquad
 \varepsilon_{\rm nil},\qquad
 \varepsilon_{\rm res}^{\circ}.
\]
\end{proposition}

\begin{proof}
Le premier objet est une fonction entre ensembles finis ; le deuxième, un
élément nilpotent d’une algèbre sur \(\mathbb F_3\) ; le troisième, un
scalaire réel avec unité angulaire. Leurs domaines, codomaines et lois de
composition sont distincts.
\end{proof}

\subsection{Distribution des indices \(90/120\) entre les \(6\) objets du système d’incidence}

\begin{definition}[Extracteur dodécaphasique]
\label{p12-83-c1-s7:def:u006e-extractor-dodecafasico}
L’extracteur entier
\begin{equation}
 \mathcal E_{\rm dod}:=(D_3,D_4,Q):
 \mathbb Z^{12}\longrightarrow\mathbb Z^{25}
 \label{p12-83-c1-s7:eq:u006e-extractor-dodecafasico}
\end{equation}
lit des différences d’arités trois et quatre et une coordonnée de somme. Les
indices ne sont pas des longueurs angulaires.
\end{definition}

\begin{definition}[Différence brute de séries]
\label{p12-83-c1-s7:def:u006e-diferencia-cruda}
Pour \(|q|<1\), on définit
\begin{equation}
 S_m^{\rm raw}(q):=\frac{q^m}{1-q^{3m}},
 \qquad
 D_{\rm raw}^{90,120}(q)
 :=S_{90}^{\rm raw}(q)-S_{120}^{\rm raw}(q).
 \label{p12-83-c1-s7:eq:u006e-series-crudas}
\end{equation}
Cet objet est une fonction analytique de \(q\) ; ce n’est pas un recensement fini.
\end{definition}

\begin{definition}[Familles marquées d’incidence]
\label{p12-83-c1-s7:def:u006e-familias-incidencia}
Une fois construites une hexade \(H\), son image
\(\ell(H)\) dans une dodécade binaire \(D\), et l’octade relevée
\(O_H\), on définira
\begin{align}
 \mathcal F_6(H)
 &:=\ell(H)\times\binom{\ell(H)}2,
 \label{p12-83-c1-s7:eq:u006e-familia-seis}\\
 \mathcal F_8(H)
 &:=O_H\times\binom{\ell(H)}2.
 \label{p12-83-c1-s7:eq:u006e-familia-ocho}
\end{align}
Par construction,
\begin{equation}
 |\mathcal F_6(H)|
 =6\binom62=90,
 \qquad
 |\mathcal F_8(H)|
 =8\binom62=120.
 \label{p12-83-c1-s7:eq:u006e-cardinales-incidencia}
\end{equation}
L’inclusion \(\ell(H)\hookrightarrow O_H\) induira
\begin{equation}
 \mathcal I_{6\to8}:
 \mathcal F_6(H)\hookrightarrow\mathcal F_8(H).
 \label{p12-83-c1-s7:eq:u006e-inclusion-incidencia}
\end{equation}
\end{definition}

Les \cref{p12-83-c1-s7:eq:u006e-familia-seis,p12-83-c1-s7:eq:u006e-familia-ocho} fixent dès maintenant le
type. Leur existence effective sera démontrée après la construction de
\(C_W\), de \(W_{12}\), de la donnée binaire \(G_{24},D,\ell\) et de \(W_{24}\).
Cette branche n’est pas anticipée ici.

\begin{definition}[Combinaison normalisée de queues]
\label{p12-83-c1-s7:def:u006e-combinacion-colas}
Pour des paramètres \(q_-,q_+\in(0,1)\), soient
\begin{equation}
 S_m(q):=\frac{q^m}{1-q^{3m}},
 \qquad
 \Delta S_m:=S_m(q_-)-S_m(q_+).
 \label{p12-83-c1-s7:eq:u006e-colas-normalizadas}
\end{equation}
La fonctionnelle
\begin{equation}
 \mathcal K_{6\to8}
 :=12\,\Delta S_{90}-\Delta S_{120}
 \label{p12-83-c1-s7:eq:u006e-funcional-seis-ocho}
\end{equation}
est une combinaison analytique de queues. Son coefficient douze appartient à sa
normalisation et n’en fait pas le registre
\(\mathcal R_{12}\).
\end{definition}

\begin{definition}[Semi-groupe des degrés et des moments]
\label{p12-83-c1-s7:def:u006e-semigrupo-momentos}
Le semi-groupe additif engendré par les deux degrés est
\begin{equation}
 \mathfrak S_{90,120}
 :=\langle90,120\rangle
 =\{90a+120b:a,b\in\mathbb N_0\}.
 \label{p12-83-c1-s7:eq:u006e-semigrupo}
\end{equation}
Pour \(n\in\mathfrak S_{90,120}\), le moment orienté est
\begin{equation}
 s_n:=q_-^n-q_+^n.
 \label{p12-83-c1-s7:eq:u006e-momento}
\end{equation}
\end{definition}

Les moments sont des monômes différentiels indexés par un semi-groupe ; les
séries de \eqref{p12-83-c1-s7:eq:u006e-colas-normalizadas} sont des sommes infinies de ces
moments sur des progressions spécifiques.

\subsection{Défaut normalisé de l’incidence}

Lorsque l’inclusion \(\mathcal I_{6\to8}\) aura été construite, sa
différence de cardinalités sera
\[
 |\mathcal F_6(H)|-|\mathcal F_8(H)|=90-120=-30.
\]
Si l’on déclare la normalisation
\[
 24\binom62=360,
\]
on obtient
\begin{equation}
 \frac{90-120}{24\binom62}
 =\frac{6-8}{24}
 =-\frac1{12}.
 \label{p12-83-c1-s7:eq:u006e-defecto-normalizado}
\end{equation}
L’inclusion détermine \(-30\) ; le rationnel \(-1/12\) dépend en outre de
la normalisation déclarée. Il ne doit pas être présenté comme si la cardinalité
fixait à elle seule le dénominateur.

\subsection{Invariants de typage}

\begin{theorem}[Séparation des six objets]
\label{p12-83-c1-s7:thm:u006e-separacion-seis-objetos}
Les objets
\begin{equation}
 \mathcal E_{\rm dod},\quad
 D_{\rm raw}^{90,120},\quad
 \mathcal I_{6\to8},\quad
 \mathcal K_{6\to8},\quad
 \mathfrak S_{90,120},\quad
 \varepsilon_{\rm res}^{\circ}
 \label{p12-83-c1-s7:eq:u006e-seis-objetos}
\end{equation}
sont mutuellement distincts par leur type : ce sont, respectivement, un lecteur entier,
une fonction analytique, un morphisme fini, une combinaison de séries, un
semi-groupe d’indices et un scalaire réel angulaire.
\end{theorem}

\begin{proof}
La conclusion s’obtient en comparant leurs domaines et codomaines :
\[
 \mathbb Z^{12}\to\mathbb Z^{25},\quad
 \mathbb D\to\mathbb C,\quad
 \mathcal F_6(H)\to\mathcal F_8(H),\quad
 (0,1)^2\to\mathbb R,\quad
 \mathbb N_0^2\to\mathbb N_0,\quad
 \{\ast\}\to\mathbb R.
\]
Il n’y a pas deux signatures égales.
\end{proof}

\begin{criterion}[Critères de réfutation]
\label{p12-83-c1-s7:crit:u006e-refutacion}
La classification est réfutée si l’un de ses domaines ou
codomaines est altéré, si le recensement de
\eqref{p12-83-c1-s7:eq:u006e-cardinales-incidencia} ne reproduit pas \(90\) et \(120\), si
l’on attribue \(-1/12\) à l’inclusion sans déclarer la normalisation, ou si
l’on présente le résidu angulaire comme une sortie du générateur de constantes de
cette monographie.
\end{criterion}


% Unidad U006A. Factor local, observacion estroboscopica y transicion exacta
% de cilindros.

\section{Facteur local de transport et transition ternaire--décimale}
\label{p12-83-c1-s8:chap:tpk-factor-local-transicion-cilindros}

Cette unité sépare trois opérations qui interviennent dans le transport du
Topological Phase Kernel et qui ne partagent pas de domaine : la réalisation cardinale
d’un trit, l’échantillonnage stroboscopique d’une suite et l’actualisation
exacte d’une paire de cylindres ternaires--décimaux. La première est une dynamique
finie ; la deuxième est un lecteur temporel ; la troisième conserve les entiers
qui permettent de décider de la compatibilité sans arithmétique à virgule flottante.

\subsection{État local et rétroaction de feuillet}

Soit
\[
 X_9=(\mathbb Z/9\mathbb Z)^2,
 \qquad
 \mathscr D_{\rm card}=\{N,E,S,O\},
\]
avec
\[
 \delta(N)=(-1,0),\quad \delta(E)=(0,1),\quad
 \delta(S)=(1,0),\quad \delta(O)=(0,-1).
\]
Les évaluations visibles de l’APP sont
\[
 L_\Sigma(x,y)=\rho_9((x+1)+(y+1)),
 \qquad
 L_\Pi(x,y)=\rho_9((x+1)(y+1)),
\]
où \(\rho_9(n)=1+((n-1)\bmod9)\). Nous définissons en outre
\[
 \rho_3^{\rm or}([0])=0,
 \qquad
 \rho_3^{\rm or}([1])=+1,
 \qquad
 \rho_3^{\rm or}([2])=-1.
\]

\begin{definition}[Facteur local de l’APP--TRIT]
\label{p12-83-c1-s8:def:u006a-factor-local}
L’espace des états locaux est
\[
 \mathcal Q_{\rm loc}=X_9\times\{\Sigma,\Pi\},
 \qquad |\mathcal Q_{\rm loc}|=162.
\]
Pour \(q=(p,m)\), \(d\in\mathscr D_{\rm card}\), nous posons
\[
 p_d=p+\delta(d),
 \qquad
 a_d(q)=\rho_3^{\rm or}([L_m(p_d)]_3),
\]
et
\[
 \mu(a,m)=
 \begin{cases}
  \Sigma,&a=+1,\\
  \Pi,&a=-1,\\
  m,&a=0.
 \end{cases}
\]
La transition locale est
\[
 F_d^{\rm loc}(p,m)=(p_d,\mu(a_d(p,m),m)).
\]
\end{definition}

La position est modifiée avant l’évaluation du feuillet. Le trit nul conserve le
dernier mode actif ; il n’élimine pas le feuillet et ne réinitialise pas l’histoire.

Pour un mot cardinal \(u=d_1d_2d_3\), soit
\[
 q_j=F_{d_j}^{\rm loc}(q_{j-1}),\qquad q_0=q,
\]
et notons
\[
 R_3(q,u)\in\{-1,0,+1\}
\]
le trit produit au troisième déplacement.

\begin{definition}[Observation de pas trois]
\label{p12-83-c1-s8:def:u006a-obs3}
Sur une suite orientée \(a=(a_t)_{t\geq1}\), nous définissons
\[
 \operatorname{Obs}_3(a)=(a_3,a_6,a_9,\ldots).
\]
L’application \(R_3\) prend un état et un mot cardinal de longueur trois ;
\(\operatorname{Obs}_3\) prend une suite et publie une sous-suite. Ce sont
des opérateurs de types distincts.
\end{definition}

\begin{theorem}[Surjectivité exacte du facteur échantillonné]
\label{p12-83-c1-s8:thm:u006a-sobreyectividad-r3}
Pour tout \(q\in\mathcal Q_{\rm loc}\) et tout
\(b\in\{-1,0,+1\}\), la fibre
\[
 \mathcal R_3(q,b)
 =\{u\in\mathscr D_{\rm card}^3:R_3(q,u)=b\}
\]
est non vide et satisfait
\[
 \boxed{8\leq |\mathcal R_3(q,b)|\leq40.}
\]
En particulier, les \(162\cdot3=486\) fibres état--sortie sont non vides.
\end{theorem}

\begin{proof}
À partir de chacun des \(162\) états, on énumère les \(4^3=64\) mots
cardinaux et l’on applique littéralement la récurrence de la
\cref{p12-83-c1-s8:def:u006a-factor-local}. Le recensement exhaustif des \(486\) paires
\((q,b)\) a pour minimum huit, pour maximum quarante et aucune valeur nulle.
L’énoncé est donc une identité finie sur le système défini, non
une hypothèse de mélange.
\end{proof}

\begin{corollary}[Réalisation d’une suite orientée arbitraire]
\label{p12-83-c1-s8:cor:u006a-realizacion-sucesion}
Pour toute suite
\((b_k)_{k\geq1}\in\{-1,0,+1\}^{\mathbb N}\), il existe une histoire cardinale
dont l’observation satisfait
\[
 (a_3,a_6,a_9,\ldots)=(b_1,b_2,b_3,\ldots).
\]
\end{corollary}

\begin{proof}
Un ordre total de \(\mathscr D_{\rm card}^3\) étant fixé, à l’étape \(k\), on
prend le premier élément de \(\mathcal R_3(q_{k-1},b_k)\). L’état terminal
du bloc appartient de nouveau à \(\mathcal Q_{\rm loc}\), de sorte que le
choix peut être itéré indéfiniment.
\end{proof}

Le corollaire démontre la capacité de réalisation du facteur échantillonné. Il ne
sélectionne pas les sections de \(\pi,e,\varphi\), ne fixe pas les deux trits
intermédiaires et ne remplace pas la compatibilité de l’état enrichi.

\subsection{Cylindres des \(2\) lecteurs}

Un mot ternaire de longueur \(L\), interprété comme l’entier \(N\),
détermine
\[
 I_3(N,L)=\left[\frac{N}{3^L},\frac{N+1}{3^L}\right).
\]
Un préfixe de \(K\) triades décimales, interprété comme l’entier \(D\),
détermine
\[
 I_{1000}(D,K)=
 \left[\frac{D}{1000^K},\frac{D+1}{1000^K}\right).
\]
Nous introduisons les résidus croisés
\[
 A=N1000^K-D3^L,
 \qquad
 B=(N+1)1000^K-D3^L=A+1000^K.
\]

\begin{lemma}[Critère entier de compatibilité]
\label{p12-83-c1-s8:lem:u006a-criterio-entero}
Les cylindres \(I_3(N,L)\) et \(I_{1000}(D,K)\) se rencontrent si et seulement si
\[
 \boxed{A<3^L,\qquad B>0.}
\]
De manière équivalente,
\[
 N1000^K<(D+1)3^L,
 \qquad
 (N+1)1000^K>D3^L.
\]
\end{lemma}

\begin{proof}
Deux intervalles semi-ouverts \([a,b)\) et \([c,d)\) se rencontrent
exactement lorsque \(a<d\) et \(c<b\). La multiplication par
\(3^L1000^K>0\) produit les deux inégalités. Soustraire \(D3^L\) et
\(N1000^K\), respectivement, donne la forme en \(A,B\).
\end{proof}

\subsection{Opérateur de complétion \(C_{3,1000}\) : extension ternaire et conservation de l’unité en \(10^3\) positions}

Un nouveau bloc de six trits est représenté par
\[
 u\in\{0,\ldots,3^6-1\}=\{0,\ldots,728\}.
\]
L’actualisation ternaire est
\[
 N'=729N+u,
 \qquad L'=L+6.
\]
La publication décimale présente deux cas. Si le cylindre raffiné ne détermine pas
encore une nouvelle triade, alors \(K'=K\), \(D'=D\). S’il détermine
\(\delta\in\{0,\ldots,999\}\), alors
\[
 K'=K+1,
 \qquad D'=1000D+\delta.
\]

\begin{definition}[État croisé et transition exacte]
\label{p12-83-c1-s8:def:u006a-c31000}
L’état des cylindres est le sextuplet
\[
 \mathfrak c=(N,L,D,K,A,B)
\]
soumise à
\[
 A=N1000^K-D3^L,
 \qquad B=A+1000^K.
\]
La transition \(\mathcal C_{3,1000}\) adjoint \(u\) et, le cas échéant,
\(\delta\), et actualise
\[
 (N,L,D,K,A,B)\longmapsto
 (N',L',D',K',A',B')
\]
par
\[
 A'=
 \begin{cases}
  729A+u1000^K,&K'=K,\\[1mm]
  729000A+u1000^{K+1}-\delta\,729\,3^L,&K'=K+1,
 \end{cases}
 \qquad
 B'=A'+1000^{K'}.
\]
Le prolongement est admissible exactement lorsque
\[
 A'<3^{L'},\qquad B'>0.
\]
\end{definition}

\begin{theorem}[Exactitude et suffisance de l’état croisé]
\label{p12-83-c1-s8:thm:u006a-exactitud-c31000}
La transition de la \cref{p12-83-c1-s8:def:u006a-c31000} conserve exactement les
identités
\[
 A'=N'1000^{K'}-D'3^{L'},
 \qquad
 B'=(N'+1)1000^{K'}-D'3^{L'}.
\]
Par conséquent, \((A',B',L',K')\) suffit à décider de la compatibilité de la
nouvelle paire de cylindres, tandis que \((N',D')\) conserve sa provenance
reconstructive.
\end{theorem}

\begin{proof}
Si aucune triade décimale n’apparaît,
\[
 N'1000^K-D3^{L+6}
 =(729N+u)1000^K-729D3^L
 =729A+u1000^K.
\]
Si aparece \(\delta\),
\begin{align*}
 N'1000^{K+1}-D'3^{L+6}
 &=(729N+u)1000^{K+1}-(1000D+\delta)7293^L\\
 &=729000A+u1000^{K+1}-\delta\,729\,3^L.
\end{align*}
Dans les deux cas, augmenter \(N'\) d’une unité ajoute \(1000^{K'}\), ce qui
donne \(B'\). Le critère de la \cref{p12-83-c1-s8:lem:u006a-criterio-entero} conclut.
\end{proof}

\begin{proposition}[Compatibilité avec la concaténation]
\label{p12-83-c1-s8:prop:u006a-concatenacion-c31000}
Appliquer successivement \(\mathcal C_{3,1000}\) à deux blocs produit la
même paire \((N'',D'')\) et les mêmes résidus croisés qu’adjoindre d’un seul
coup les deux mots ternaires et les triades décimales publiées dans le
même ordre.
\end{proposition}

\begin{proof}
Les récurrences \(N\mapsto729N+u\) et
\(D\mapsto1000D+\delta\) sont les lois de concaténation positionnelle en
bases \(729\) et \(1000\). Le
\cref{p12-83-c1-s8:thm:u006a-exactitud-c31000} identifie \(A,B\) à des fonctions de ces
quatre entiers ; leur actualisation par étapes coïncide donc avec
l’évaluation postérieure à la concaténation.
\end{proof}

\subsection{Décomposition du Topological Phase Kernel en opérateurs d’avancée, de rotation, de retenue et de mémoire}

La chaîne correcte est
\[
 (q,u)\xmapsto{\ R_3\ }b,
 \qquad
 (a_t)_{t\geq1}
 \xmapsto{\ \operatorname{Obs}_3\ }
 (a_{3k})_{k\geq1},
 \qquad
 (\mathfrak c,u,\delta)
 \xmapsto{\ \mathcal C_{3,1000}\ }
 \mathfrak c'.
\]
\(R_3\) prouve la réalisabilité locale ; \(\operatorname{Obs}_3\) publie une
sous-suite ; \(\mathcal C_{3,1000}\) transporte la compatibilité entre deux
bases. Aucun des trois ne peut remplacer la relation de survie
du TPK complet.

\begin{criterion}[Critères de réfutation]
\label{p12-83-c1-s8:crit:u006a-refutacion}
La construction est réfutée dans son domaine si l’un quelconque des
faits suivants se produit :
\begin{enumerate}
 \item l’une des \(486\) fibres \(\mathcal R_3(q,b)\) est vide ou se trouve
       en dehors de l’intervalle \([8,40]\) ;
 \item le trit nul change le mode actif ;
 \item on identifie \(R_3\) à \(\operatorname{Obs}_3\) ;
 \item la récurrence de \(A'\) ne reproduit pas sa définition directe ;
 \item un prolongement déclaré compatible viole \(A'<3^{L'}\) ou
       \(B'>0\) ;
 \item la décision dépend d’une tolérance numérique absente des
       inégalités entières.
\end{enumerate}
\end{criterion}


% Unidad U006D. Registro dodecafásico integral y descomposición 1+5+6.

\section{Registre dodécaphasique entier, extension cyclique et décomposition \texorpdfstring{\(1+5+6\)}{1+5+6}}
\label{p12-83-c1-s9:chap:registro-dodecafasico-integral}

La période observable de cent huit positions contient douze secteurs de
neuf positions. Cette égalité de cardinalités ne suffit pas à décrire
l’état temporel : il faut distinguer le support ordonné, la loi cyclique,
l’indice sectoriel et le registre enrichi des transports, orientations,
retenues et incidences. Cette unité construit les quatre objets et démontre
la reconstruction entière des douze coordonnées de mémoire.

\subsection{Support ordonné et coordonnée cyclique}

Pour \(m\in\{1,\ldots,12\}\), nous définissons
\begin{equation}
 E_m:=\{9(m-1)+1,\ldots,9m\}.
 \label{p12-83-c1-s9:eq:u006d-bloques-nonadicos}
\end{equation}
Le support ordonné et le groupe cyclique sont
\begin{equation}
 \Gamma_{108}:=
 E_1\sqcup\cdots\sqcup E_{12},
 \qquad
 \Theta_{108}:=\mathbb Z/108\mathbb Z.
 \label{p12-83-c1-s9:eq:u006d-gamma-theta}
\end{equation}
\(\Gamma_{108}\) conserve la position, l’ordre et l’appartenance à un bloc ;
\(\Theta_{108}\) conserve la loi de translation. Ce ne sont pas le même objet.

Pour \(\theta\in\Theta_{108}\), soit
\(\widetilde\theta\in\{0,\ldots,107\}\) son représentant. Nous définissons
\begin{equation}
 \beta(\theta):=
 \left[
 \left\lfloor\frac{\widetilde\theta}{9}\right\rfloor
 \right]_{12},
 \qquad
 r(\theta):=1+(\widetilde\theta\bmod9).
 \label{p12-83-c1-s9:eq:u006d-coordenadas-bloque-residuo}
\end{equation}
La première coordonnée localise le secteur dodécaphasique ; la seconde publie
la position positive \(1,\ldots,9\) à l’intérieur du secteur.

\subsection{Extension cyclique non scindée}

Soient \(C_n=\mathbb Z/n\mathbb Z\). Les applications
\begin{equation}
 \iota:C_{12}\longrightarrow C_{108},
 \quad \iota([b]_{12})=[9b]_{108},
 \qquad
 p:C_{108}\longrightarrow C_9,
 \quad p([\theta]_{108})=[\theta]_9
 \label{p12-83-c1-s9:eq:u006d-mapas-extension}
\end{equation}
forment la suite
\begin{equation}
 0\longrightarrow C_{12}
 \xrightarrow{\;\iota\;}C_{108}
 \xrightarrow{\;p\;}C_9
 \longrightarrow0.
 \label{p12-83-c1-s9:eq:u006d-extension}
\end{equation}
Pour la section ensembliste
\[
 s([r]_9):=[\widetilde r]_{108},
 \qquad 0\leq\widetilde r<9,
\]
le défaut d’additivité est
\begin{equation}
 c(r,r'):=
 \left[
 \left\lfloor
 \frac{\widetilde r+\widetilde r'}9
 \right\rfloor
 \right]_{12},
 \qquad
 s(r)+s(r')
 =s(r+r')+\iota(c(r,r')).
 \label{p12-83-c1-s9:eq:u006d-cociclo-extension}
\end{equation}

\begin{theorem}[Exactitude, non-scindement et classe cohomologique]
\label{p12-83-c1-s9:thm:u006d-extension-no-escindida}
La suite \eqref{p12-83-c1-s9:eq:u006d-extension} est exacte et ne se scinde pas. Pour
l’action triviale,
\begin{equation}
 H^2(C_9,C_{12})
 \simeq C_{12}/9C_{12}
 \simeq C_3,
 \label{p12-83-c1-s9:eq:u006d-cohomologia-extension}
\end{equation}
et le cocycle \eqref{p12-83-c1-s9:eq:u006d-cociclo-extension} représente sa classe
génératrice.
\end{theorem}

\begin{proof}
L’image de \(\iota\) est le sous-groupe des multiples de neuf, qui coïncide
avec \(\ker p\) ; \(p\) est surjective. Cela prouve l’exactitude. Si
l’extension se scindait, \(C_{108}\) serait isomorphe à
\(C_{12}\times C_9\), mais
\[
 \exp(C_{108})=108,
 \qquad
 \exp(C_{12}\times C_9)=\operatorname{mcm}(12,9)=36.
\]
L’identité de cocycle est la division euclidienne de
\(\widetilde r+\widetilde r'\) par neuf. Pour un groupe cyclique à action
triviale,
\(H^2(C_9,C_{12})\simeq C_{12}/9C_{12}\) ; la retenue unitaire se projette sur
\([1]\), qui engendre le quotient d’ordre trois.
\end{proof}

Le non-scindement type le retour :
\[
 g_0(\theta+9)=g_0(\theta),
 \qquad
 \beta(\theta+9)=\beta(\theta)+1\pmod{12}.
\]
Neuf itérations ferment la phase visible et translatent d’un secteur ; cent
huit ferment les deux coordonnées de l’horloge. L’état enrichi conserve
en outre la mémoire verticale, le cylindre, la frontière et la généalogie.

\subsection{Registre temporel enrichi : encodage conjoint de la phase, du feuillet, de la frontière et de la mémoire}

\begin{definition}[Registre dodécaphasique orienté]
\label{p12-83-c1-s9:def:u006d-registro-dodecafasico}
Pour une histoire enrichie, soit \(\tau_m\) la sous-route ordonnée sur
\(E_m\), \(\sigma_m\) son orientation, \(\kappa_m\) la retenue qui franchit la
frontière du bloc et \(\Lambda_m\) son livre local d’incidences. Le
registre complet est
\begin{equation}
 \mathcal R_{12}:=
 \bigl(
 (E_m,\tau_m,\sigma_m,\kappa_m,\Lambda_m)
 \bigr)_{m=1}^{12}.
 \label{p12-83-c1-s9:eq:u006d-registro-r12}
\end{equation}
\end{definition}

On distingue ainsi quatre niveaux :
\begin{equation}
 \boxed{
 \mathcal R_{12}
 \neq C_{12}
 \neq\Theta_{108}
 \neq\Gamma_{108}.}
 \label{p12-83-c1-s9:eq:u006d-cuatro-objetos}
\end{equation}
Le premier est un registre enrichi ; le deuxième, un indice de secteur ; le
troisième, un groupe de translations ; le quatrième, un support ordonné. Il existe
des projections entre eux, mais pas d’identifications.

Soit \(q_{\rm mem}\in\mathbb Z_{\geq0}\) le nombre non borné de tours
nonadiques. La coordonnée dodécaphasique n’est que son résidu
\[
 \beta=[q_{\rm mem}]_{12}.
\]
Après cent huit pas,
\[
 \beta\longmapsto\beta,
 \qquad
 q_{\rm mem}\longmapsto q_{\rm mem}+12.
\]
Le calendrier revient ; la mémoire verticale n’est pas réinitialisée.

\subsection{Extracteur signé d’événements}

Soit \((d_t)_{t=1}^{108}\) le registre de chiffres sur
\(\Gamma_{108}\). Nous définissons
\begin{equation}
 \mathcal D_{\rm evt}:=\{2,4,6,8\},
 \qquad
 \Sigma_{12}:=(+,+,+,-,-,-,+,+,+,-,-,-).
 \label{p12-83-c1-s9:eq:u006d-datos-eventos}
\end{equation}
Pour \(m=0,\ldots,11\), soit
\begin{equation}
 e_m:=
 \Sigma_{12}(m+1)\,
 \#\{t\in E_{m+1}:d_t\in\mathcal D_{\rm evt}\}
 \in\mathbb Z.
 \label{p12-83-c1-s9:eq:u006d-eventos-firmados}
\end{equation}
L’application conserve le recensement local et l’orientation sectorielle comme
facteurs séparés.

Nous apparions les deux demi-couronnes par
\begin{equation}
 p_i:=e_i+e_{i+6},
 \qquad
 a_i:=e_i-e_{i+6},
 \qquad 0\leq i\leq5.
 \label{p12-83-c1-s9:eq:u006d-pares-antipodales}
\end{equation}
Nous définissons le recensement total, les cinq différences par rapport à la sixième paire et
les six antisymétries :
\begin{equation}
 s_C:=\sum_{i=0}^{5}p_i,
 \qquad
 M_i:=p_i-p_5\quad(0\leq i\leq4),
 \qquad
 A_6:=(a_0,\ldots,a_5).
 \label{p12-83-c1-s9:eq:u006d-coordenadas-uno-cinco-seis}
\end{equation}

\begin{definition}[Lecteur entier \(1+5+6\)]
\label{p12-83-c1-s9:def:u006d-lector-integral}
Le lecteur linéaire est
\begin{equation}
 \mathcal L_{12}:\mathbb Z^{12}\longrightarrow\mathbb Z^{12},
 \qquad
 \mathcal L_{12}(e)
 :=
 (s_C,M_0,\ldots,M_4,a_0,\ldots,a_5).
 \label{p12-83-c1-s9:eq:u006d-lector-l12}
\end{equation}
La distribution des coordonnées est
\begin{equation}
 \underbrace{1}_{s_C}
 +\underbrace{5}_{M_0,\ldots,M_4}
 +\underbrace{6}_{a_0,\ldots,a_5}
 =12.
 \label{p12-83-c1-s9:eq:u006d-uno-cinco-seis}
\end{equation}
\end{definition}

\begin{theorem}[Reconstruction exacte et image entière]
\label{p12-83-c1-s9:thm:u006d-inversa-l12}
Le lecteur \(\mathcal L_{12}\) est injectif. Sur son image,
\begin{align}
 p_5&=
 \frac{s_C-\sum_{i=0}^{4}M_i}{6},
 &
 p_i&=p_5+M_i\quad(0\leq i\leq4),
 \label{p12-83-c1-s9:eq:u006d-inversa-p}\\
 e_i&=\frac{p_i+a_i}{2},
 &
 e_{i+6}&=\frac{p_i-a_i}{2}
 \quad(0\leq i\leq5).
 \label{p12-83-c1-s9:eq:u006d-inversa-e}
\end{align}
Un tuple de sortie appartient à \(\operatorname{im}\mathcal L_{12}\) si et
seulement si
\begin{equation}
 s_C-\sum_{i=0}^{4}M_i\equiv0\pmod6
 \label{p12-83-c1-s9:eq:u006d-condicion-divisibilidad}
\end{equation}
et, pour les \(p_i\) reconstruits,
\begin{equation}
 p_i\equiv a_i\pmod2
 \qquad(0\leq i\leq5).
 \label{p12-83-c1-s9:eq:u006d-condicion-paridad}
\end{equation}
\end{theorem}

\begin{proof}
De \(M_i=p_i-p_5\), on obtient
\[
 s_C=\sum_{i=0}^{4}(p_5+M_i)+p_5
 =6p_5+\sum_{i=0}^{4}M_i,
\]
ce qui donne la première équation de \eqref{p12-83-c1-s9:eq:u006d-inversa-p} ; les autres sont
immédiates. Additionner et soustraire
\(p_i=e_i+e_{i+6}\) et \(a_i=e_i-e_{i+6}\) produit
\eqref{p12-83-c1-s9:eq:u006d-inversa-e}. La divisibilité par six et les six
congruences de parité sont précisément les conditions nécessaires et
suffisantes pour que toutes les coordonnées reconstruites soient entières.
\end{proof}

La division par six apparaît après le recensement orienté, non pendant celui-ci. Les
conditions d’intégralité de cette unité sont exactement
\eqref{p12-83-c1-s9:eq:u006d-condicion-divisibilidad} et
\eqref{p12-83-c1-s9:eq:u006d-condicion-paridad} ; on n’attribue pas ici de forme normale de
Smith à une matrice qui n’aurait pas été préalablement définie.

\subsection{\(2\) mots dodécaphasiques de provenances différentes}

La distinction doit être conservée
\begin{equation}
 U_{12}^{\rm cat}:=U_6\Vert U_6,
 \qquad
 U_{12}^{\rm sgn}:=\mathcal H_{12}(K).
 \label{p12-83-c1-s9:eq:u006d-dos-palabras-doce}
\end{equation}
La première concatène deux émissions du catalogue de longueur six. La seconde
est une coordonnée signée reliée réversiblement à \(K\) par l’opérateur
\(\mathcal H_{12}\). Leur longueur commune n’implique pas une même provenance,
sémantique ou loi de transport.

\begin{criterion}[Critères de réfutation]
\label{p12-83-c1-s9:crit:u006d-refutacion}
La construction est réfutée si la suite
\eqref{p12-83-c1-s9:eq:u006d-extension} n’est pas exacte, si sa classe cohomologique est
triviale, si un tuple de l’image viole
\eqref{p12-83-c1-s9:eq:u006d-condicion-divisibilidad} ou
\eqref{p12-83-c1-s9:eq:u006d-condicion-paridad}, si les formules inverses ne retrouvent pas
les douze entiers initiaux, ou si le retour de \(\beta\) est interprété comme
une réinitialisation de \(q_{\rm mem}\).
\end{criterion}


% Unidad U006B. Funtor ciclotomico y elevacion modular de caminos.

\section{Foncteur cyclotomique de phase--longueur et relèvement modulaire des chemins}
\label{p12-83-c1-s10:chap:funtor-ciclotomico-elevacion-caminos}

Le quotient odométrique \(\mathcal O_{9,12}\) conserve simultanément la
position nonadique, le registre dodécaphasique et la longueur entière de chaque
flèche. Dans cette unité, sa phase est réalisée sur des fibres \(A_2\) et l’on prouve
le relèvement cardinal qui remplace chaque arête par quatre arêtes
orientées. La représentation cyclotomique et la longueur absolue
resteront des coordonnées distinctes.

\subsection{Représentation de Coxeter d’ordre \(3\) sur les fibres dodécaphasiques de type \(A_2\)}

Soit
\begin{equation}
 G_{A_2}:=
 \begin{pmatrix}2&-1\\-1&2\end{pmatrix},
 \qquad
 C:=
 \begin{pmatrix}0&-1\\1&-1\end{pmatrix}.
 \label{p12-83-c1-s10:eq:u006b-dato-coxeter}
\end{equation}
Par le \cref{p12-83-c1-s5:prop:coxeter-a2-autonomo},
\begin{equation}
 C^3=I_2,
 \qquad
 C^{\mathsf T}G_{A_2}C=G_{A_2},
 \qquad
 \det(XI_2-C)=X^2+X+1.
 \label{p12-83-c1-s10:eq:u006b-identidades-coxeter}
\end{equation}
Ainsi, \(C\) réalise le caractère cyclotomique de \(C_3\) sur \(A_2\).

Pour chaque \(b\in\mathbb Z/12\mathbb Z\), soit \(A_{2,b}\) une copie de
\(A_2\), soit
\[
 \iota_b:A_{2,b}\xrightarrow{\;\sim\;}A_2
\]
une identification entière et soit
\(P_b\in O(G_{A_2})\) un repère qui dépend de \(b\), mais non du résidu
\(r\). Nous définissons
\begin{equation}
 \rho_b:=
 \iota_b^{-1}P_bCP_b^{-1}\iota_b,
 \qquad
 \rho_b^3=\operatorname{id}_{A_{2,b}}.
 \label{p12-83-c1-s10:eq:u006b-rho-fibra}
\end{equation}
L’indépendance par rapport à \(r\) lie le caractère à la fibre
dodécaphasique ; \(r\) conserve la position interne du tour nonadique.

\subsection{Catégorie odométrique du TPK : cylindres, troncatures et morphismes de succession}

Rappelons que \(\mathcal O_{9,12}\) a pour objets
\[
 (b,r)\in\mathbb Z/12\mathbb Z\times\{0,\ldots,8\}.
\]
Pour \(n\in\mathbb N\), l’unique flèche de longueur \(n\) est
\begin{equation}
 \gamma_{(b,r),n}:
 (b,r)\longrightarrow
 \left(
 b+\kappa_r(n),[r+n]_9
 \right),
 \qquad
 \kappa_r(n):=
 \left\lfloor\frac{r+n}{9}\right\rfloor.
 \label{p12-83-c1-s10:eq:u006b-flecha-odometro}
\end{equation}
La composition s’appuie sur
\begin{equation}
 \kappa_r(n+m)
 =\kappa_r(n)+\kappa_{[r+n]_9}(m).
 \label{p12-83-c1-s10:eq:u006b-cociclo-odometro}
\end{equation}
La longueur \(\ell(\gamma_{(b,r),n})=n\) est additive.

\begin{definition}[Foncteur cyclotomique de phase--longueur]
\label{p12-83-c1-s10:def:u006b-funtor-ciclotomico}
Soit \(\mathbf{Rep}_{\mathbb C}(C_3)\) la catégorie des représentations
complexes de \(C_3\), et soit \(\mathbf B\mathbb N\) la catégorie monoïdale
à un seul objet avec endomorphismes \((\mathbb N,+)\). Nous définissons
\begin{equation}
 \mathfrak F_{\rm cyc}:
 \mathcal O_{9,12}
 \longrightarrow
 \mathbf{Rep}_{\mathbb C}(C_3)\times\mathbf B\mathbb N
 \label{p12-83-c1-s10:eq:u006b-signatura-funtor}
\end{equation}
par
\[
 \mathfrak F_{\rm cyc}(b,r)
 =(A_{2,b}\otimes_{\mathbb Z}\mathbb C,\rho_b)
\]
et
\begin{equation}
 \mathfrak F_{\rm cyc}(\gamma)
 :=
 \left(
 \iota_{b'}^{-1}P_{b'}C^{\ell(\gamma)}
 P_b^{-1}\iota_b,
 \ell(\gamma)
 \right)
 \label{p12-83-c1-s10:eq:u006b-flecha-funtor}
\end{equation}
pour \(\gamma:(b,r)\to(b',r')\).
\end{definition}

\begin{theorem}[Fonctorialité et séparation des registres]
\label{p12-83-c1-s10:thm:u006b-functorialidad}
L’affectation de la \cref{p12-83-c1-s10:def:u006b-funtor-ciclotomico} est un foncteur. Sa
première composante conserve uniquement la classe de longueur modulo trois ; la
seconde conserve l’entier complet et, par réduction, détermine cette classe.
La composante cyclotomique ne permet pas de reconstruire la longueur absolue.
\end{theorem}

\begin{proof}
Si \(\gamma_1:(b_0,r_0)\to(b_1,r_1)\) et
\(\gamma_2:(b_1,r_1)\to(b_2,r_2)\) ont pour longueurs \(n_1,n_2\), la
trivialisation intermédiaire s’annule :
\begin{align}
 &\bigl(
 \iota_{b_2}^{-1}P_{b_2}C^{n_2}P_{b_1}^{-1}\iota_{b_1}
 \bigr)
 \bigl(
 \iota_{b_1}^{-1}P_{b_1}C^{n_1}P_{b_0}^{-1}\iota_{b_0}
 \bigr)\notag\\
 &\qquad=
 \iota_{b_2}^{-1}P_{b_2}C^{n_1+n_2}
 P_{b_0}^{-1}\iota_{b_0}.
 \label{p12-83-c1-s10:eq:u006b-composicion-funtor}
\end{align}
La seconde composante se compose par addition. Comme \(C^n\) ne dépend que de
\([n]_3\), mais que \(\ell=n\) conserve l’entier, les deux coordonnées ne sont pas
identifiées.
\end{proof}

Si
\[
 Q:\mathsf{TPK}\longrightarrow\mathcal O_{9,12}
\]
est le foncteur odométrique de U006, la réalisation typée est
\begin{equation}
 \mathfrak F_{\rm cyc}\circ Q:
 \mathsf{TPK}\longrightarrow
 \mathbf{Rep}_{\mathbb C}(C_3)\times\mathbf B\mathbb N.
 \label{p12-83-c1-s10:eq:u006b-compuesto-tpk-ciclotomico}
\end{equation}

\subsection{Relèvement modulaire des chemins cardinaux}

Soit \(\mathcal A_{N,d}\) la catégorie libre des chemins cardinaux de
\((\mathbb Z/N\mathbb Z)^d\). Soient
\[
 N\ge2,\quad u\ge1,\quad\gcd(u,N)=1,
 \quad k:=\operatorname{ord}_N(u),
 \quad q:=\frac{u^k-1}{N}.
\]

\begin{definition}[Substitution cardinale par \(u\)]
\label{p12-83-c1-s10:def:u006b-sustitucion-cardinal}
Le foncteur
\[
 \mathcal S_u:\mathcal A_{N,d}\longrightarrow\mathcal A_{N,d}
\]
envoie un objet \(p\) sur \(up\) et chaque arête orientée sur la concaténation
ordonnée de \(u\) copies de même direction.
\end{definition}

\begin{theorem}[Relèvement modulaire en endomorphisme de chemins]
\label{p12-83-c1-s10:thm:u006b-elevacion-modular}
Pour tout objet \(p\), chemin \(\gamma\) et
\[
 \kappa_{N,r}(n):=
 \left\lfloor\frac{r+n}{N}\right\rfloor,
\]
on a
\begin{align}
 \mathcal S_u^k(p)&=p,
 \label{p12-83-c1-s10:eq:u006b-retorno-objetos}\\
 \ell(\mathcal S_u^k\gamma)&=u^k\ell(\gamma),
 \label{p12-83-c1-s10:eq:u006b-longitud-sustituida}\\
 \kappa_{N,r}(u^kn)-\kappa_{N,r}(n)&=qn.
 \label{p12-83-c1-s10:eq:u006b-diferencia-vueltas}
\end{align}
\end{theorem}

\begin{proof}
La catégorie est libre, de sorte que les images des objets et des générateurs
déterminent un unique foncteur. Chaque application multiplie la longueur par
\(u\). Comme \(u^k=1+qN\), les objets reviennent modulo \(N\), et
\[
 \left\lfloor
 \frac{r+n+qNn}{N}
 \right\rfloor
 =
 \left\lfloor\frac{r+n}{N}\right\rfloor+qn.
\]
\end{proof}

Dans le cas nonadique,
\begin{equation}
 (N,u,k,q)=(9,4,3,7),
 \qquad
 4^3=64=1+7\cdot9.
 \label{p12-83-c1-s10:eq:u006b-parametros-nonadicos}
\end{equation}
Par conséquent,
\begin{equation}
 \boxed{
 \mathcal S_4^3(e)
 \text{ conserve les extrémités et contient sept tours nonadiques
 orientés par arête initiale}.}
 \label{p12-83-c1-s10:eq:u006b-siete-vueltas}
\end{equation}
Les extrémités et la position résiduelle reviennent en trois applications ;
l’état enrichi ne revient pas nécessairement. Sans preuve
d’inversibilité, \(\mathcal S_4^3\) est un endomorphisme vertical des chemins,
non une monodromie.

\begin{remark}[Séparation d’avec la dilatation de l’horloge]
\label{p12-83-c1-s10:rem:u006b-s4-no-d4}
\(\mathcal S_4\) substitue les objets et les arêtes dans une catégorie de chemins.
L’application \(D_4\) de U006C multipliera par quatre la coordonnée
\(\Theta_{108}\). Elles partagent un facteur entier, mais ni le domaine, ni le codomaine, ni la
loi de composition.
\end{remark}

\begin{criterion}[Critères de réfutation]
\label{p12-83-c1-s10:crit:u006b-refutacion}
Le bloc est réfuté si l’une quelconque des identités
\eqref{p12-83-c1-s10:eq:u006b-identidades-coxeter},
\eqref{p12-83-c1-s10:eq:u006b-cociclo-odometro},
\eqref{p12-83-c1-s10:eq:u006b-composicion-funtor} ou
\eqref{p12-83-c1-s10:eq:u006b-parametros-nonadicos} échoue ; il est également mal typé si l’on
identifie la phase \(C^n\), la longueur \(n\), la substitution
\(\mathcal S_4\) et la dilatation \(D_4\).
\end{criterion}
}%


% Unidad U004. Redacción autónoma desde los generadors TRIT de 1063/live.

\section{TRIT comme unité minimale d’information topologique : orientation ternaire, mémoire de retenue et ordre temporel}
\label{p12-83-c1-s11:chap:trit-levantado-regimenes}

La courbure mixte de l’APP a produit trois valeurs orientées : zéro dans les feuillets
purs et des signes opposés lorsque l’ordre somme--produit est interverti. Le TRIT
type cette triade et conserve, outre le résidu visible, le quotient nécessaire
pour reconstruire l’entier relevé. Cette unité sépare expressément quatre
niveaux : classe ternaire, représentant équilibré, retenue et régime
quadratique.

\begin{definition}[Unité d’information topologique TRIT]
Le \emph{TRIT} est l’unité élémentaire orientée dont l’espace des états est
\(\mathsf{TRIT}=\{-1,0,+1\}\). Sa capacité logarithmique native est
\(\log 3\) et son expression en unités binaires est \(\log_2 3\). À la différence
d’un bit, ses trois états ne se réduisent pas à un encodage cardinal : ils typent,
respectivement, l’ouverture, le seuil et le retour, et conservent le quotient de
retenue requis pour reconstruire la transition qui a produit le résidu
visible.
\end{definition}

\subsection{Système ternaire orienté et représentation équilibrée des états TRIT}

Nous définissons l’ensemble orienté
\[
 \mathsf{TRIT}:=\{-1,0,+1\}.
\]
La carte équilibrée est la bijection d’ensembles
\[
 \operatorname{bal}:\mathbb F_3\longrightarrow\mathsf{TRIT},
 \qquad
 [0]_3\mapsto0,\qquad[1]_3\mapsto+1,\qquad[2]_3\mapsto-1.
\]
Il ne s’agit pas d’une égalité de types. \(\mathbb F_3\) porte la somme et le produit
d’un corps ; \(\mathsf{TRIT}\) porte l’ordre et l’orientation. Toute opération
transportée entre les deux doit indiquer la carte utilisée.

L’orientation des neuf phases est
\[
 \varepsilon_9(r)=
 \begin{cases}
 +1,&r\in\{1,4,7\},\\
 -1,&r\in\{2,5,8\},\\
 0,&r\in\{3,6,9\}.
 \end{cases}
\]
Sa table complète est
\[
\begin{array}{c|ccccccccc}
r&1&2&3&4&5&6&7&8&9\\\hline
\varepsilon_9(r)&+1&-1&0&+1&-1&0&+1&-1&0
\end{array}.
\]
La périodicité de trois n’élimine pas la phase nonadique : les phases
\(1,4,7\), par exemple, ont la même orientation et des positions distinctes
dans le tour.

\subsection{Algorithme de division ternaire équilibrée et unicité de la paire quotient–résidu}

\begin{theorem}[Relèvement ternaire local]
\label{p12-83-c1-s11:thm:trit-levantamiento-local}
Pour chaque \(n\in\{-4,-3,\ldots,4\}\), il existe une unique paire
\((r,c)\in\mathsf{TRIT}^2\) telle que
\[
 n=r+3c.
\]
\end{theorem}

\begin{proof}
L’existence se vérifie dans la table
\[
\begin{array}{c|rrrrrrrrr}
n&-4&-3&-2&-1&0&1&2&3&4\\\hline
r&-1&0&+1&-1&0&+1&-1&0&+1\\
c&-1&-1&-1&0&0&0&+1&+1&+1
\end{array}.
\]
Si \(r+3c=r'+3c'\), alors \(r-r'=3(c'-c)\). Comme
\(r-r'\in\{-2,-1,0,1,2\}\), le seul multiple de trois possible est zéro ;
par conséquent, \(r=r'\) et \(c=c'\).
\end{proof}

La projection visible \((r,c)\mapsto r\) possède, au-dessus de la valeur zéro, trois
états locaux :
\[
 0^-:=(0,-1),
 \qquad
 0^0:=(0,0),
 \qquad
 0^+:=(0,+1).
\]
Ils publient le même résidu, mais reconstruisent \(-3,0,+3\). Cette fibre triple
est la forme minimale d’une distinction qui réapparaîtra dans l’état enrichi
du Topological Phase Kernel.

\begin{corollary}[Itération équilibrée]
\label{p12-83-c1-s11:cor:trit-iteracion-balanceada}
Tout entier \(N\) admet un développement ternaire équilibré fini
\[
 N=r_0+3r_1+\cdots+3^mr_m,
 \qquad r_j\in\mathsf{TRIT}.
\]
Il s’obtient en itérant la division \(N_k=r_k+3N_{k+1}\), où \(r_k\) est le
représentant équilibré de \(N_k\bmod3\).
\end{corollary}

\begin{proof}
La division euclidienne modulo trois donne un résidu dans \(\{0,1,2\}\) ; remplacer
\(2\) par \(-1\) accroît le quotient d’une unité et produit
\(r_k\in\mathsf{TRIT}\). La valeur absolue du quotient diminue en dehors
d’un ensemble fini, de sorte que le processus se termine. L’unicité s’obtient
en comparant le premier chiffre non coïncident et en réduisant modulo trois.
\end{proof}

\subsection{Composition et cocycle de retenue}

Pour \(r,s\in\mathsf{TRIT}\), nous définissons \(r\oplus s\in\mathsf{TRIT}\) et
\(\kappa_3(r,s)\in\mathbb Z\) par l’unique égalité
\[
 r+s=(r\oplus s)+3\kappa_3(r,s).
\]
La table est
\[
\begin{array}{c|ccc}
\oplus&-1&0&+1\\\hline
-1&+1&-1&0\\
0&-1&0&+1\\
+1&0&+1&-1
\end{array},
\qquad
\begin{array}{c|ccc}
\kappa_3&-1&0&+1\\\hline
-1&-1&0&0\\
0&0&0&0\\
+1&0&0&+1
\end{array}.
\]

\begin{proposition}[Identité cocyclique ternaire]
\label{p12-83-c1-s11:prop:trit-cociclo}
Pour \(r,s,t\in\mathsf{TRIT}\),
\[
 \kappa_3(r,s)+\kappa_3(r\oplus s,t)
 =\kappa_3(s,t)+\kappa_3(r,s\oplus t).
\]
\end{proposition}

\begin{proof}
Les deux membres sont le quotient total qui apparaît lorsque l’on écrit
\(r+s+t\) comme un représentant équilibré plus un multiple de trois.
L’unicité du relèvement prouve l’égalité.
\end{proof}

Si \((r,c)\) et \((s,d)\) sont des états relevés, leur somme est
\[
 (r,c)\boxplus(s,d)
 :=\bigl(r\oplus s,\ c+d+\kappa_3(r,s)\bigr).
\]
La \cref{p12-83-c1-s11:prop:trit-cociclo} prouve que \(\boxplus\) est associative et que la
reconstruction \(\operatorname{rec}_3(r,c)=r+3c\) satisfait
\[
 \operatorname{rec}_3(x\boxplus y)
 =\operatorname{rec}_3(x)+\operatorname{rec}_3(y).
\]

\subsection{Involution d’orientation et ordre temporel}

Pour un mot tritique \(\tau=(\tau_1,\ldots,\tau_n)\), nous définissons
\[
 \mathcal C_{\rm rev}(\tau_1,\ldots,\tau_n)
 :=(-\tau_n,\ldots,-\tau_1).
\]

\begin{lemma}[Renversement orienté]
\label{p12-83-c1-s11:lem:trit-reversion}
L’application \(\mathcal C_{\rm rev}\) est une involution. Elle échange les
secteurs \(+1\) et \(-1\) et fixe le secteur \(0\).
\end{lemma}

\begin{proof}
L’appliquer deux fois inverse deux fois l’ordre et change deux fois chaque
signe. On retrouve ainsi le mot initial. Les autres affirmations sont
immédiates dans chaque composante.
\end{proof}

Lorsqu’un paramètre temporel orienté est fixé, la convention du TPK est
\[
 +1:\text{retour avec mémoire},
 \qquad
 0:\text{seuil présent},
 \qquad
 -1:\text{ouverture de prolongement}.
\]
Cette affectation n’ajoute pas un quatrième état temporel. La direction est conservée
dans le signe et la position exacte est conservée dans la phase nonadique.

\subsection{Famille universelle d’algèbres quadratiques}

Pour \(\tau\in\mathsf{TRIT}\), nous définissons
\[
 \mathbb A_\tau:=\mathbb R[J_\tau]/(J_\tau^2+\tau I).
\]
Les trois cas sont :
\[
\begin{array}{c|c|c|c}
\tau&J_\tau^2&\text{régime}&\exp(tJ_\tau)\\\hline
+1&-I&\text{elliptique}&\cos t\,I+\sin t\,J_{+1}\\
0&0&\text{parabolique}&I+tJ_0\\
-1&+I&\text{hyperbolique}&\cosh t\,I+\sinh t\,J_{-1}
\end{array}.
\]

\begin{theorem}[Exponentielle quadratique uniforme]
\label{p12-83-c1-s11:thm:trit-exponencial-uniforme}
Si \(J_\tau^2=-\tau I\), alors
\[
 \exp(tJ_\tau)
 =C_\tau(t)I+S_\tau(t)J_\tau,
\]
où
\[
\begin{array}{c|ccc}
\tau&+1&0&-1\\\hline
C_\tau(t)&\cos t&1&\cosh t\\
S_\tau(t)&\sin t&t&\sinh t
\end{array}.
\]
\end{theorem}

\begin{proof}
Nous séparons la série exponentielle en puissances paires et impaires. Les relations
\(J_{+1}^{2m}=(-1)^mI\), \(J_0^2=0\) et
\(J_{-1}^{2m}=I\) produisent, respectivement, les séries du cosinus et du sinus,
la troncature linéaire et les séries hyperboliques.
\end{proof}

Sur \(\mathbb F_3\) apparaissent les analogues
\[
 \mathbb F_9=\mathbb F_3[\iota]/(\iota^2+1),
 \qquad
 \mathbb D_3=\mathbb F_3[\epsilon]/(\epsilon^2),
 \qquad
 \mathbb S_3=\mathbb F_3[j]/(j^2-1).
\]
Le premier est un corps à neuf éléments ; le deuxième contient un nilpotent
non nul ; le troisième se décompose au moyen des idempotents
\((1\pm j)/2\) et est isomorphe à \(\mathbb F_3\times\mathbb F_3\). La
coïncidence de cardinalités n’identifie pas ces trois algèbres.

La source entière du régime elliptique est
\[
 \mathbb Q_{\mathbb Z}:=\mathbb Z[u]/(u^2+1).
\]
La réduction modulo trois envoie \(u\) sur \(\iota\) et produit
\(\mathbb F_9\). L’extension des scalaires à \(\mathbb R\) et une
représentation \(u\mapsto J_{+1}\) produisent le réalisateur réel elliptique.
Il n’existe donc pas d’application non typée \(\mathbb F_9\to\mathbb R\) : les deux
réalisations procèdent d’une source entière commune.

\subsection{Représentations matricielles du TRIT sélectionnées par le signe de \(J_\tau^2\)}

Soit
\[
 C=\begin{pmatrix}0&-1\\1&-1\end{pmatrix},
 \qquad C^3=I,
\]
et définissons
\[
 J_e:=\frac{2C+I}{\sqrt3}
 =\frac1{\sqrt3}\begin{pmatrix}1&-2\\2&-1\end{pmatrix}.
\]
Une multiplication directe donne
\[
 (2C+I)^2=-3I,
 \qquad J_e^2=-I.
\]
Le réalisateur parabolique est
\[
 J_p:=N=\begin{pmatrix}0&1\\0&0\end{pmatrix},
 \qquad N^2=0.
\]
Pour le régime hyperbolique, nous prenons
\[
 F_{\rm av}=\begin{pmatrix}0&1\\1&1\end{pmatrix},
 \qquad
 A:=F_{\rm av}^2=\begin{pmatrix}1&1\\1&2\end{pmatrix},
\]
et
\[
 J_h:=\frac{2A-3I}{\sqrt5}
 =\frac1{\sqrt5}\begin{pmatrix}-1&2\\2&1\end{pmatrix}.
\]
Alors
\[
 (2A-3I)^2=5I,
 \qquad J_h^2=I.
\]

La forme bilinéaire
\[
 G=\begin{pmatrix}2&1\\1&-2\end{pmatrix}
\]
satisfait
\[
 A^TGA=G.
\]
De plus,
\[
 \cosh(2\log\varphi)=\frac32,
 \qquad
 \sinh(2\log\varphi)=\frac{\sqrt5}{2},
\]
de sorte que
\[
 A=\exp(2\log\varphi\,J_h)
  =\cosh(2\log\varphi)I+\sinh(2\log\varphi)J_h.
\]
Cette identité sera retrouvée causalement dans la biographie de
\(\varphi\) ; ici, seule la représentation matricielle est vérifiée.

Les trois contrôles exponentiels sont
\[
 \exp(\pi J_e)+I=0,
 \qquad
 \exp(J_p)=I+J_p,
 \qquad
 \exp(2\log\varphi\,J_h)=F_{\rm av}^2.
\]
La présence de \(\pi\) et de \(\varphi\) dans ces identités de vérification
n’autorise pas à les utiliser comme constructeurs de ces constantes. Leurs
constructeurs apparaîtront ensuite à partir de la généalogie nonadique.

\subsection{Couplage local APP--TRIT}

Avec l’échange
\[
 R=\begin{pmatrix}0&1\\1&0\end{pmatrix}
\]
nous définissons, pour \(\tau\in\mathsf{TRIT}\),
\[
 B_\tau=(4+3\tau)I+(5-3\tau)R.
\]
Puisque \(P_\pm=(I\pm R)/2\),
\[
 B_\tau=9P_++(6\tau-1)P_-.
\]

\begin{theorem}[Récupération spectrale de l’orientation]
\label{p12-83-c1-s11:thm:app-trit-recuperacion-espectral}
Le spectre de \(B_\tau\) est
\[
 \operatorname{Spec}(B_\tau)=\{9,6\tau-1\}.
\]
La valeur propre différentielle \(\lambda_-\) détermine univoquement le TRIT :
\[
 \tau=\frac{\lambda_-+1}{6}.
\]
De plus, le bloc central de l’APP est \(B_{+1}\).
\end{theorem}

\begin{proof}
Les projecteurs \(P_+\) et \(P_-\) diagonalisent \(R\) avec les valeurs propres
\(+1\) et \(-1\). La formule spectrale s’obtient par substitution. La deuxième
identité isole \(\tau\). Pour \(\tau=+1\),
\(B_{+1}=7I+2R=\begin{psmallmatrix}7&2\\2&7\end{psmallmatrix}\), qui est
l’unique bloc de la \cref{p12-83-c1-s1:thm:app-bloque-central}.
\end{proof}

Par conséquent, le TRIT n’est pas ajouté comme étiquette rétrospective : il peut être retrouvé
à partir du mode différentiel d’une famille construite sur l’échange interne
de l’APP.

\subsection{Obstructions du TRIT : séparation résidu–représentant–retenue, involutivité et récupération spectrale de \(\tau\)}

\begin{criterion}[Critères de réfutation du TRIT]
L’unité est réfutée si :
\begin{enumerate}
 \item on identifie la classe résiduelle, le représentant équilibré et la retenue ;
 \item on affirme que \(0^-\), \(0^0\) et \(0^+\) sont le même état ;
 \item la composition \(\boxplus\) cesse de reconstruire la somme entière ;
 \item \(\mathcal C_{\rm rev}\) n’est pas involutive ;
 \item on identifie les trois algèbres parce qu’elles ont le même cardinal ;
 \item un réalisateur ne satisfait pas \(J_\tau^2=-\tau I\) ;
 \item la valeur de \(\tau\) est choisie après observation du résultat et n’est pas
       retrouvée à partir du spectre de \(B_\tau\).
\end{enumerate}
\end{criterion}

L’APP fournit le support, la matrice multiplicative et la loi additive
interne ; TRIT fournit l’orientation,
le relèvement et le type quadratique. L’unité suivante construira le
Topological Phase Kernel comme catégorie de transport de ces données, sans
confondre le sélecteur interne avec l’état complet.

\HMTDeferredTPKBase
\HMTDeferredTPKDevelopment


% Unidad U006C. Dilatación aritmética, cociclo de acarreo y
% semiconjugación cuaternaria.

\section{Dilatation quaternaire du registre, obstruction observable et semi-conjugaison raffinée}
\label{p12-83-c1-s12:chap:dilatacion-cuaternaria-semiconjugacion}

La substitution cardinale \(\mathcal S_4\) construite dans l’unité précédente
agit sur les chemins. Dans cette unité, on définit une deuxième opération, de
nature arithmétique, sur le registre
\(\mathbb Z/12\mathbb Z\times\mathbb Z/9\mathbb Z\), et une troisième
opération, catégorielle, qui subdivise chaque transition observable en quatre
microtransitions. L’objectif n’est pas de les fusionner, mais de démontrer la relation
exacte entre leurs réalisations.

\subsection{Dilatation sur le registre produit}

Soit
\[
 \mathcal Q_{12,9}:=
 \mathbb Z/12\mathbb Z\times\mathbb Z/9\mathbb Z.
\]
Pour \(r\in\mathbb Z/9\mathbb Z\), nous notons
\(\widetilde r\in\{0,\ldots,8\}\) son représentant normalisé.

\begin{definition}[Dilatation quaternaire du registre]
\label{p12-83-c1-s12:def:u006c-dilatacion-registro}
Nous définissons
\begin{equation}
 D_4(b,r):=
 \left(
 \left[
  4b+\left\lfloor\frac{4\widetilde r}{9}\right\rfloor
 \right]_{12},
 [4r]_9
 \right).
 \label{p12-83-c1-s12:eq:u006c-d4-producto}
\end{equation}
La coordonnée totale est
\begin{equation}
 \Theta:\mathcal Q_{12,9}\xrightarrow{\;\sim\;}
 \mathbb Z/108\mathbb Z,
 \qquad
 \Theta(b,r):=[9b+\widetilde r]_{108}.
 \label{p12-83-c1-s12:eq:u006c-coordenada-total}
\end{equation}
\end{definition}

\begin{proposition}[Réalisation par multiplication modulaire]
\label{p12-83-c1-s12:prop:u006c-d4-total}
Dans la coordonnée totale,
\begin{equation}
 \boxed{
 \Theta\circ D_4\circ\Theta^{-1}(\theta)=[4\theta]_{108}.}
 \label{p12-83-c1-s12:eq:u006c-d4-total}
\end{equation}
En particulier,
\begin{equation}
 \operatorname{im}(D_4)
 =4\mathbb Z/108\mathbb Z
 \simeq\mathbb Z/27\mathbb Z,
 \label{p12-83-c1-s12:eq:u006c-imagen-d4}
\end{equation}
et la restriction induite sur l’image est d’ordre neuf.
\end{proposition}

\begin{proof}
Pour \(\theta=9b+\widetilde r\),
\[
 4\theta
 =9\left(
  4b+\left\lfloor\frac{4\widetilde r}{9}\right\rfloor
 \right)
 +(4\widetilde r\bmod9).
\]
C’est exactement la décomposition quotient--résidu de
\eqref{p12-83-c1-s12:eq:u006c-d4-producto}. L’image est formée des classes
multiples de quatre. Après division par quatre, on obtient
\(\mathbb Z/27\mathbb Z\), et
\[
\begin{aligned}
 4^1&\equiv4,& 4^2&\equiv16,& 4^3&\equiv10,\\
 4^4&\equiv13,& 4^5&\equiv25,& 4^6&\equiv19,\\
 4^7&\equiv22,& 4^8&\equiv7,& 4^9&\equiv1\pmod{27}.
\end{aligned}
\]
Aucune puissance précédente n’est égale à un ; par conséquent,
\(\operatorname{ord}_{27}(4)=9\).
\end{proof}

\subsection{Identité exacte de retenue}

Pour \(r\in\{0,\ldots,8\}\) et \(n\in\mathbb N\), nous définissons
\[
 \kappa_r(n):=\left\lfloor\frac{r+n}{9}\right\rfloor,
 \qquad
 r':=[r+n]_9.
\]

\begin{theorem}[Cocycle quaternaire de retenue]
\label{p12-83-c1-s12:thm:u006c-cociclo-d4}
On a
\begin{equation}
 \boxed{
 \kappa_{[4r]_9}(4n)
 =4\kappa_r(n)
 +\left\lfloor\frac{4\widetilde r'}{9}\right\rfloor
 -\left\lfloor\frac{4\widetilde r}{9}\right\rfloor.}
 \label{p12-83-c1-s12:eq:u006c-cociclo-d4}
\end{equation}
La dernière différence est le cobord produit par le choix des
représentants \(r\mapsto\widetilde r\).
\end{theorem}

\begin{proof}
Écrivons
\[
 r+n=9\kappa_r(n)+\widetilde r'.
\]
En multipliant par quatre,
\[
 4r+4n=36\kappa_r(n)+4\widetilde r'.
\]
Séparer le quotient et le résidu modulo neuf et soustraire le quotient déjà contenu dans
\(4\widetilde r\) produit
\eqref{p12-83-c1-s12:eq:u006c-cociclo-d4}.
\end{proof}

\begin{corollary}[Troisième itération]
\label{p12-83-c1-s12:cor:u006c-tercera-iteracion}
Pour tout \((b,r)\in\mathcal Q_{12,9}\),
\begin{equation}
 \boxed{D_4^3(b,r)=([4b+7\widetilde r]_{12},r).}
 \label{p12-83-c1-s12:eq:u006c-d4-cubo}
\end{equation}
\end{corollary}

\begin{proof}
Par la \cref{p12-83-c1-s12:prop:u006c-d4-total},
\[
 \Theta(D_4^3(b,r))
 =[4^3(9b+\widetilde r)]_{108}
 =[64(9b+\widetilde r)]_{108}.
\]
Comme \(64=1+7\cdot9\), le résidu modulo neuf redevient \(r\), tandis que
la coordonnée de bloc se transforme en
\([4b+7\widetilde r]_{12}\).
\end{proof}

Le retour de \(r\) n’est pas le retour du registre complet. La différence
\(7\widetilde r\) conserve la mémoire du représentant interne traversé.

\subsection{Obstruction dans le quotient observable}

Dans la configuration observable de
\cref{p12-83-c1-s3:def:tpk-configuracion-observable}, nous écrivons
\[
 q=(p^\Sigma,p^\Pi,d^\Sigma,d^\Pi,\theta)
 \in\mathcal Q_{\rm obs}.
\]
La projection positionnelle est
\begin{equation}
 \pi_{\rm pos}(q):=(p^\Sigma,p^\Pi)\in X_9^2.
 \label{p12-83-c1-s12:eq:u006c-proyeccion-posicional}
\end{equation}
La signature temporelle est
\[
 \tau(q):=\varepsilon_9(r(\theta))\in\{-1,0,+1\}.
\]
Nous définissons le déplacement observable orienté
\begin{equation}
 a(q):=
 \begin{cases}
  (\delta(d^\Sigma),0),&\tau(q)=+1,\\
  (0,\delta(d^\Pi)),&\tau(q)=-1,\\
  (0,0),&\tau(q)=0.
 \end{cases}
 \label{p12-83-c1-s12:eq:u006c-desplazamiento-observable}
\end{equation}
Alors
\[
 \pi_{\rm pos}(F_{\rm obs}q)=\pi_{\rm pos}(q)+a(q),
\]
avec l’actualisation des directions et de l’horloge définie dans U005.

\begin{theorem}[Obstruction à l’endofoncteur observable ordinaire]
\label{p12-83-c1-s12:thm:u006c-obstruccion}
Il n’existe pas d’endofoncteur de l’espace observable ordinaire qui satisfasse
simultanément les conditions suivantes :
\begin{enumerate}
 \item multiplier la position géométrique par quatre ;
 \item remplacer chaque transition par quatre microtransitions véritables ;
 \item conserver le feuillet et l’orientation à chaque microtransition ;
 \item envoyer une fibre dodécaphasique complète vers une seule fibre
 dodécaphasique.
\end{enumerate}
L’obstruction persiste à un pas de seuil avec \(a(q)=0\).
\end{theorem}

\begin{proof}
Soit
\[
 \mathcal B_b:=\{(b,r):0\leq r\leq8\}.
\]
Remplacer une itération élémentaire par quatre force
\[
 h(\theta+1)=h(\theta)+4
\]
et, par récurrence,
\[
 h(\theta)=4\theta+c\pmod{108}.
\]
Prenons \(0\leq\widetilde c<108\) et écrivons
\(\widetilde c=9c_b+c_0\), avec \(0\leq c_0<9\). Si toute la fibre
\(\mathcal B_b\) était envoyée vers une seule fibre, l’indice de bloc de
destination devrait être indépendant de \(r\) ; il contient cependant
\begin{equation}
 4b+c_b+
 \left\lfloor\frac{4r+c_0}{9}\right\rfloor
 \pmod{12},
 \label{p12-83-c1-s12:eq:u006c-bloque-obstruccion}
\end{equation}
qui n’est pas constant en \(r\). De plus, l’opérateur de transition neutre
satisfait \(P_0^{\rm TPK}=I\), différent des trivialisations
cyclotomiques \(P_b\) ; le calendrier ordinaire \((+1,-1,0)\) fait que
quatre itérations élémentaires visitent les deux modes arithmétiques et le mode
neutre.
\end{proof}

\subsection{Raffinement minimal par microinstants}

\begin{definition}[État observable subdivisé]
\label{p12-83-c1-s12:def:u006c-estado-subdividido}
Soit
\begin{equation}
 \widetilde{\mathcal Q}_4
 :=\mathcal Q_{\rm obs}\times\{0,1,2,3\},
 \qquad
 \widetilde\pi(q,j)
 :=4\pi_{\rm pos}(q)+j\,a(q).
 \label{p12-83-c1-s12:eq:u006c-estado-refinado}
\end{equation}
On définit
\begin{equation}
 \widetilde F(q,j):=
 \begin{cases}
  (q,j+1),&0\leq j<3,\\
  (F_{\rm obs}q,0),&j=3,
 \end{cases}
 \qquad
 \iota_4(q):=(q,0).
 \label{p12-83-c1-s12:eq:u006c-dinamica-refinada}
\end{equation}
\end{definition}

\begin{theorem}[Semi-conjugaison à quatre pas]
\label{p12-83-c1-s12:thm:u006c-semiconjugacion}
L’inclusion \(\iota_4\) satisfait
\begin{equation}
 \boxed{
 \widetilde F^4\circ\iota_4
 =\iota_4\circ F_{\rm obs}.}
 \label{p12-83-c1-s12:eq:u006c-semiconjugacion}
\end{equation}
Les positions intermédiaires sont exactement
\[
 4\pi_{\rm pos}(q)+j\,a(q),
 \qquad j=0,1,2,3.
\]
\end{theorem}

\begin{proof}
À partir de \((q,0)\), les trois premières applications n’accroissent que l’indice
de microinstant. La quatrième applique \(F_{\rm obs}\) et ramène l’indice à
zéro. Cela donne
\[
 (q,0)\mapsto(q,1)\mapsto(q,2)\mapsto(q,3)
 \mapsto(F_{\rm obs}q,0).
\]
La formule des positions est la définition de
\(\widetilde\pi\). Le raffinement rend visible une subdivision déterminée ;
il n’ajoute pas de décision dynamique.
\end{proof}

\subsection{Catégorie subdivisée et relèvement cyclotomique}

Soit \(\operatorname{sd}_4\mathsf P_{\rm obs}\) la catégorie obtenue en
remplaçant chaque générateur
\(\gamma:q\to F_{\rm obs}q\) par la chaîne
\begin{equation}
 (q,0)\to(q,1)\to(q,2)\to(q,3)\to(F_{\rm obs}q,0).
 \label{p12-83-c1-s12:eq:u006c-cadena-subdivision}
\end{equation}
Le foncteur
\[
 \operatorname{sd}_4:
 \mathsf P_{\rm obs}\longrightarrow
 \operatorname{sd}_4\mathsf P_{\rm obs}
\]
envoie chaque générateur sur cette chaîne. Le foncteur de contraction
\[
 \operatorname{col}_4:
 \operatorname{sd}_4\mathsf P_{\rm obs}
 \longrightarrow\mathsf P_{\rm obs}
\]
recompose les quatre microflèches et satisfait
\begin{equation}
 \operatorname{col}_4\circ\operatorname{sd}_4
 =\operatorname{id}_{\mathsf P_{\rm obs}}.
 \label{p12-83-c1-s12:eq:u006c-colapso}
\end{equation}

Dans une chaîne qui reste dans la fibre \(b\), les trois premières
microflèches agissent par \(\rho_b\) ; la quatrième incorpore le transport
interfibre s’il existe un franchissement dodécaphasique. Comme \(\rho_b^3=I\), le produit
ordonné des quatre actions coïncide avec le morphisme cyclotomique de la
flèche originale.

\begin{proposition}[Invariance cyclotomique par subdivision]
\label{p12-83-c1-s12:prop:u006c-invariancia-ciclotomica}
Il existe un relèvement typé
\[
 \widetilde{\mathfrak F}_{\rm cyc}:
 \operatorname{sd}_4\mathsf P_{\rm obs}
 \longrightarrow
 \mathbf{Rep}_{\mathbb C}(C_3)\times\mathbf B\mathbb N
\]
tel que, pour toute flèche \(\gamma\) de
\(\mathsf P_{\rm obs}\),
\begin{equation}
 \boxed{
 \widetilde{\mathfrak F}_{\rm cyc}
 \bigl(\operatorname{sd}_4\gamma\bigr)
 =
 \bigl(\mathfrak F_{\rm cyc}\circ Q\bigr)(\gamma).}
 \label{p12-83-c1-s12:eq:u006c-invariancia-ciclotomica}
\end{equation}
\end{proposition}

\begin{proof}
On attribue aux microflèches internes les puissances successives de
\(\rho_b\), et à la quatrième le transport entre les repères des fibres
d’origine et de destination. La composition des trois premiers facteurs est
l’identité parce que \(\rho_b^3=I\). Le facteur restant est exactement le
morphisme que \(\mathfrak F_{\rm cyc}\circ Q\) attribue à la flèche
contractée. La catégorie subdivisée est libre sur ces microflèches, de sorte que
l’affectation s’étend de manière unique.
\end{proof}

\subsection{\(3\) opérations quaternaires non interchangeables}

L’architecture contient trois opérations distinctes :
\begin{enumerate}
 \item \(\mathcal S_4\), substitution cardinale d’objets et de chemins ;
 \item \(D_4\), multiplication par quatre dans
 \(\Theta_{108}\) ;
 \item \(\operatorname{sd}_4\), subdivision fonctorielle d’une transition
 observable.
\end{enumerate}
Leurs compatibilités sont données, respectivement, par
\[
 4^3=1+7\cdot9,\qquad
 \Theta D_4=4\Theta,\qquad
 \widetilde F^4\iota_4=\iota_4F_{\rm obs}.
\]
Partager l’entier quatre ne transforme aucune d’elles en une autre.

\begin{criterion}[Critères de réfutation]
\label{p12-83-c1-s12:crit:u006c-refutacion}
La construction est réfutée si l’une quelconque des identités
\eqref{p12-83-c1-s12:eq:u006c-d4-total},
\eqref{p12-83-c1-s12:eq:u006c-cociclo-d4},
\eqref{p12-83-c1-s12:eq:u006c-d4-cubo},
\eqref{p12-83-c1-s12:eq:u006c-semiconjugacion} ou
\eqref{p12-83-c1-s12:eq:u006c-invariancia-ciclotomica} échoue ; également si une réalisation
observable satisfait les quatre conditions du
\cref{p12-83-c1-s12:thm:u006c-obstruccion} sans conserver un registre équivalent au
microinstant.
\end{criterion}


% Unidad U006. Lectores, odómetro y calendario; redacción autónoma.

\section{Lecteurs modulaires, odomètre et calendrier dodécaphasique}
\label{p12-83-c1-s13:chap:tpk-lectores-odometro-calendario}

L’état enrichi du TPK n’est identifié à aucune de ses lectures.
Cette unité définit les lecteurs modulo neuf, modulo trois et modulo mille,
l’odomètre phase--secteur, le quotient prédictif de trente-six états et
l’extension cyclique de cent huit positions. Chaque construction déclare ce qu’elle
publie et quelle mémoire elle conserve.

\subsection{Schéma typé des lecteurs modulaires du TPK : domaine, codomaine, publication et mémoire conservée}

\begin{definition}[Fiche opératoire]
\label{p12-83-c1-s13:def:tpk-ficha-operatoria}
Un opérateur canonique du TPK est enregistré au moyen d’un sextuplet
\[
 \mathcal O=(D_{\mathcal O},C_{\mathcal O},f_{\mathcal O},
 I_{\mathcal O},L_{\mathcal O},M_{\mathcal O}),
\]
où \(D_{\mathcal O}\) est le domaine, \(C_{\mathcal O}\) le codomaine,
\(f_{\mathcal O}\) la règle, \(I_{\mathcal O}\) les invariants préservés,
\(L_{\mathcal O}\) l’information publiée ou oubliée et
\(M_{\mathcal O}\) la mémoire qui demeure dans l’état enrichi.
\end{definition}

Deux symboles de même cardinalité ou de même période ne représentent pas le même
opérateur, à moins qu’il existe une équivalence entrelaçant leurs domaines,
règles et invariants. Cette convention empêchera de confondre le quotient
prédictif \(C_{36}^{\rm pred}\), la frontière \(R_{36}\) et un défaut
cardinal de valeur trente-six.

\subsection{Lecteur nonadique de l’état enrichi : phase, transition et retour avec mémoire}

Pour \(n\geq1\), nous définissons
\[
 \rho_9(n)=1+((n-1)\bmod9),
 \qquad
 q_9(n)=\left\lfloor\frac{n-1}{9}\right\rfloor.
\]
Alors
\[
 n=9q_9(n)+\rho_9(n).
\]
La fiche de \(\rho_9\) est
\[
 \bigl(
 \mathbb Z_{>0},\mathcal D_9,\rho_9,
 n\equiv\rho_9(n)\pmod9,
 \text{phase positive},
 q_9(n)
 \bigr).
\]
La racine numérique positive coïncide numériquement avec \(\rho_9\), mais sa
sémantique est distincte : l’une se présente comme section du quotient ; l’autre,
comme somme itérée de chiffres.

\subsection{Lecteur tritique de l’état orienté : ordre temporel et transport de retenue}

La division équilibrée s’écrit
\[
 n=r_3(n)+3c_3(n),
 \qquad r_3(n)\in\mathsf{TRIT},\quad c_3(n)\in\mathbb Z.
\]
Sa fiche est
\[
 \bigl(
 \mathbb Z,\mathsf{TRIT}\times\mathbb Z,
 n\mapsto(r_3(n),c_3(n)),
 \operatorname{rec}_3(r,c)=r+3c,
 r_3,
 c_3
 \bigr).
\]
La publication peut ne montrer que \(r_3\), mais le TPK conserve
\(c_3\) si une composition ultérieure exige une reconstruction.

\subsection{Lecteur positionnel par blocs en base \(10^3\)}

La division euclidienne définit
\[
 N=1000q_{1000}(N)+r_{1000}(N),
 \qquad 0\leq r_{1000}(N)<1000.
\]
Sa fiche est
\[
 \bigl(
 \mathbb Z,\mathbb Z\times\{0,\ldots,999\},
 N\mapsto(q_{1000},r_{1000}),
 \operatorname{rec}_{1000}(q,r)=1000q+r,
 r_{1000},
 q_{1000}
 \bigr).
\]

\begin{proposition}[Concaténation et lecture résiduelle]
\label{p12-83-c1-s13:prop:tpk-concatenacion-base-mil}
Si \(u_1,\ldots,u_m\in\{0,\ldots,999\}\) et
\[
 D_m=1000D_{m-1}+u_m,
 \qquad D_0=0,
\]
alors
\[
 D_m=\sum_{j=1}^m u_j1000^{m-j}.
\]
De plus,
\[
 D_m\equiv\sum_{j=1}^m u_j\pmod9
\]
parce que \(1000\equiv1\pmod9\).
\end{proposition}

\begin{proof}
La première égalité s’obtient par récurrence sur \(m\). La seconde procède
de sa réduction modulo neuf.
\end{proof}

La congruence ne reconstruit pas l’ordre des blocs. Le registre généalogique conserve la
séquence \((u_1,\ldots,u_m)\) ; le lecteur résiduel ne publie que leur somme
modulo neuf.

\subsection{Catégorie odométrique du TPK : objets cylindriques et composition par troncature}

\begin{definition}[Odomètre nonadique--dodécaphasique]
\label{p12-83-c1-s13:def:tpk-odometro-9-12}
La catégorie \(\mathcal O_{9,12}\) a pour objets
\[
 (b,r)\in C_{12}\times\{0,\ldots,8\}.
\]
Une flèche de longueur \(n\in\mathbb N\) est
\[
 (b,r)\xrightarrow{\ n\ }
 \left(
  b+\left\lfloor\frac{r+n}{9}\right\rfloor,
  r+n\bmod9
 \right).
\]
La composition additionne les longueurs.
\end{definition}

Soit
\[
 \kappa_r(n)=\left\lfloor\frac{r+n}{9}\right\rfloor.
\]

\begin{lemma}[Cocycle de l’odomètre]
\label{p12-83-c1-s13:lem:tpk-cociclo-odometro}
Si \(r_1=r+n_1\bmod9\), alors
\[
 \kappa_r(n_1+n_2)
 =\kappa_r(n_1)+\kappa_{r_1}(n_2).
\]
\end{lemma}

\begin{proof}
Nous écrivons \(r+n_1=9\kappa_r(n_1)+r_1\), avec
\(0\leq r_1<9\). En ajoutant \(n_2\) et en divisant de nouveau par neuf,
on obtient l’identité.
\end{proof}

Pour un état \(x\) dont la base contient \(\theta\), nous définissons
\[
 Q_0(x)=(\beta(\theta),\widetilde\theta\bmod9).
\]
Une flèche TPK a une longueur égale au nombre de générateurs observables qui
la composent.

\begin{theorem}[Foncteur phase--longueur]
\label{p12-83-c1-s13:thm:tpk-funtor-fase-longitud}
L’affectation qui envoie un état \(x\) sur \(Q_0(x)\) et une flèche de
longueur \(n\) sur la flèche odométrique de longueur \(n\) définit un foncteur
\[
 Q:\mathsf{TPK}\longrightarrow\mathcal O_{9,12}.
\]
\end{theorem}

\begin{proof}
Chaque générateur augmente \(\theta\) d’une unité. Après \(n\) pas, le
résidu de phase augmente de \(n\) modulo neuf et le secteur augmente de
\(\lfloor(r+n)/9\rfloor\). L’identité a une longueur nulle et le
\cref{p12-83-c1-s13:lem:tpk-cociclo-odometro} prouve la compatibilité avec la composition.
\end{proof}

Après neuf pas,
\[
 (b,r)\longmapsto(b+1,r).
\]
Après cent huit,
\[
 (b,r)\longmapsto(b,r).
\]
Ces identités appartiennent au quotient odométrique. Elles n’affirment pas que
le cylindre, la frontière, le registre généalogique ou la mémoire verticale soient revenus.

\subsection{Quotient minimal de mémoire prédictive}

Soit
\[
 X_{9,12}=\{(a,m):a\in\mathbb Z/9\mathbb Z, 0\leq m<12\}
\]
avec successeur
\[
 S(a,m)=
 \begin{cases}
  (a,m+1),&m<11,\\
  (a+1\bmod9,0),&m=11.
 \end{cases}
\]
Nous définissons les observables
\[
 \beta_{\rm pred}(a,m)=4a\pmod{12},
 \qquad
 d_{\rm pred}(a,m)=1+(m+\beta_{\rm pred}(a,m)\pmod{12}).
\]
Pour une observable \(q\), le mot futur inclusif d’horizon \(h\)
est
\[
 Q_h^q(x)=(q(x),q(Sx),\ldots,q(S^hx)).
\]

\begin{theorem}[Quotient prédictif minimal]
\label{p12-83-c1-s13:thm:tpk-cociente-predictivo-36}
Les observables \(\beta_{\rm pred}\), \(d_{\rm pred}\) et la paire
\((\beta_{\rm pred},d_{\rm pred})\) produisent le même quotient stable
\[
 \pi_{36}(a,m)=(a\bmod3,m).
\]
Il a trente-six états et des fibres de cardinal trois. Les horizons
minimaux sont \(11\), \(8\) et \(0\), respectivement.
\end{theorem}

\begin{proof}
Pour un système fini \((X,T)\), soit
\[
 E_h=\operatorname{Eq}(Q_h^q).
\]
Alors
\[
 E_{h+1}=E_h\cap(T\times T)^{-1}(E_h).
\]
La chaîne décroissante se stabilise et sa limite est la plus grande équivalence
\(T\)-compatible contenue dans \(\operatorname{Eq}(q)\). Dans le système
déclaré, l’énumération exacte produit le nombre de classes suivant :
\[
\begin{array}{c|rrrrrrrrrrrr}
h&0&1&2&3&4&5&6&7&8&9&10&11\\\hline
\#(X/E_h^{\beta})&3&6&9&12&15&18&21&24&27&30&33&36
\end{array},
\]
et
\[
\begin{array}{c|rrrrrrrrr}
h&0&1&2&3&4&5&6&7&8\\\hline
\#(X/E_h^{d})&12&15&18&21&24&27&30&33&36.
\end{array}
\]
La paire distingue déjà les trente-six classes à l’horizon zéro. Dans tous
les cas, les classes stables conservent \(m\) et \(a\bmod3\) ; les trois
possibilités de \(a\) ayant le même résidu modulo trois forment chaque fibre.
\end{proof}

Nous notons cet objet \(C_{36}^{\rm pred}\). Ce n’est pas la frontière
\(R_{36}\), qui contiendra trois canaux ordonnés de six trits, ni un recensement
de trente-six éléments d’une autre construction.

\subsection{Extension centrale du calendrier}

On définit
\[
 \iota:C_{12}\longrightarrow C_{108},
 \qquad \iota([b]_{12})=[9b]_{108},
\]
et
\[
 p:C_{108}\longrightarrow C_9,
 \qquad p([\theta]_{108})=[\theta]_9.
\]
On obtient la suite exacte
\begin{equation}
 0\longrightarrow C_{12}
 \xrightarrow{\ \iota\ }C_{108}
 \xrightarrow{\ p\ }C_9
 \longrightarrow0.
\label{p12-83-c1-s13:eq:tpk-extension-12-108-9}
\end{equation}
Une section ensembliste \(s([r]_9)=[r]_{108}\),
\(0\leq r<9\), induit le cocycle
\[
 c(r,r')=
 \left[\left\lfloor\frac{r+r'}{9}\right\rfloor\right]_{12}.
\]

\begin{proposition}[L’extension ne se scinde pas]
\label{p12-83-c1-s13:prop:tpk-extension-no-escindida}
La suite \cref{p12-83-c1-s13:eq:tpk-extension-12-108-9} n’est pas isomorphe à l’extension
produit \(C_{12}\times C_9\).
\end{proposition}

\begin{proof}
L’exposant de \(C_{108}\) est \(108\). L’exposant de
\(C_{12}\times C_9\) est
\(\operatorname{mcm}(12,9)=36\). Deux groupes finis isomorphes ont le
même exposant ; ils ne sont donc pas isomorphes.
\end{proof}

\begin{theorem}[Nécessité et période exacte de cent huit positions]
\label{p12-83-c1-s13:thm:tpk-periodo-exacto-108}
La translation \(\theta\mapsto\theta+1\) sur \(C_{108}\) est d’ordre
exactement \(108\). En conséquence, une itération \(n\) fait revenir le
calendrier complet si et seulement si \(108\mid n\).
\end{theorem}

\begin{proof}
La \(n\)-ième puissance envoie \([\theta]\) sur \([\theta+n]_{108}\). C’est
l’identité pour tout \(\theta\) exactement lorsque \([n]_{108}=0\), c’est-à-
dire lorsque \(108\mid n\).
\end{proof}

En particulier, \(1008\) n’est pas un autre retour du calendrier :
\[
 1008\equiv36\pmod{108}.
\]
Après \(1008\) pas, la phase modulo neuf revient, mais la coordonnée de
\(C_{108}\) s’est déplacée de trente-six positions. Naturellement,
\(216,324,\ldots\) font également revenir le calendrier, mais ce sont des répétitions du
tour fondamental de longueur \(108\).

\subsection{Mémoire verticale non bornée}

Soit \(q_{\rm mem}\in\mathbb Z_{\geq0}\) le nombre non borné de tours
nonadiques et soit
\[
 \beta=[q_{\rm mem}]_{12}.
\]
Alors un tour de neuf pas satisfait
\[
 q_{\rm mem}\longmapsto q_{\rm mem}+1,
 \qquad
 \beta\longmapsto\beta+1\pmod{12}.
\]
Après cent huit pas,
\[
 \beta\longmapsto\beta,
 \qquad
 q_{\rm mem}\longmapsto q_{\rm mem}+12.
\]
Le calendrier revient ; la mémoire verticale, non.

\begin{corollary}[Hiérarchie des retours]
\label{p12-83-c1-s13:cor:tpk-jerarquia-retornos}
Dans le quotient de phase,
\(9\) pas constituent un tour. Les morphismes composés
\[
 \mathcal K_{27}=F_{\rm obs}^{27},
 \qquad
 F_{\rm obs}^{54},
 \qquad
 F_{\rm obs}^{108}
\]
contiennent, respectivement, trois, six et douze tours nonadiques. Aucun
n’efface par définition le registre généalogique ou la frontière.
\end{corollary}

\subsection{Système minimal de types de l’état enrichi et compatibilités entre ses lecteurs}

La table suivante répond explicitement à la question des objets internes qui
existent déjà dans le TPK avant l’introduction des germes :
\[
\begin{array}{p{.23\textwidth}p{.27\textwidth}p{.38\textwidth}}
\toprule
Objeto u operador&Dominio y codominio&Función\\\midrule
\(X_9\)&\((\mathbb Z/9)^2\)&soporte APP\\
\(T_N,T_E,T_S,T_O\)&\(X_9\to X_9\)&traslaciones cardinales\\
\(M_{\rm ph}\)&\(\mathcal D_9\to\mathsf{TRIT}\)&selección de hoja\\
\(F_{\rm obs}\)&\(\mathcal Q_{\rm obs}\to\mathcal Q_{\rm obs}\)&sucesor observable\\
\((\rho_9,q_9)\)&\(\mathbb Z_{>0}\to\mathcal D_9\times\mathbb Z_{\ge0}\)&fase y altura\\
\((r_3,c_3)\)&\(\mathbb Z\to\mathsf{TRIT}\times\mathbb Z\)&orientación y carry\\
\((q_{1000},r_{1000})\)&\(\mathbb Z\to\mathbb Z\times\{0,\ldots,999\}\)&bloque decimal y cociente\\
\(Q\)&\(\mathsf{TPK}\to\mathcal O_{9,12}\)&fase y longitud\\
\(C_{36}^{\rm pred}\)&cociente de \(X_{9,12}\)&memoria predictiva mínima\\
\(C_{108}\)&calendario cíclico&extensión no escindida\\
\(q_{\rm mem}\)&\(\mathbb Z_{\ge0}\)&memoria vertical no acotada\\
\(\operatorname{Ext}_r\)&
\(\operatorname{Ext}_r\subseteq\mathcal F_q\times\mathcal F_{q'}\)&
prolongaciones compatibles\\
\bottomrule
\end{array}
\]

\subsection{Obstructions des lecteurs et du calendrier du TPK : quotients euclidiens, \(C_{36}^{\mathrm{pred}}\), extension \(C_{108}\) et mémoire verticale}

\begin{criterion}[Critères de réfutation des lecteurs et du calendrier]
L’unité est réfutée si :
\begin{enumerate}
 \item un lecteur résiduel est présenté comme une transition complète ;
 \item on omet le quotient d’une division euclidienne utilisée ensuite ;
 \item on identifie \(C_{36}^{\rm pred}\) à \(R_{36}\) ;
 \item on affirme que l’extension \(C_{108}\) est le produit
       \(C_{12}\times C_9\) ;
 \item on déclare un retour complet pour un nombre non divisible par \(108\) ;
 \item le retour de \(\beta\) force \(q_{\rm mem}=0\) ;
 \item \(\mathcal K_{27}\) est présenté comme une horloge distincte et non comme
       une puissance du transport.
\end{enumerate}
\end{criterion}

\section{Structure opératorielle complète de l’APP et hiérarchie arithmétique--géométrique}

Le développement opératoriel de l’APP fixe la hiérarchie
\(4\to2\to6\to54\), les projecteurs arithmétiques, le bloc central et la
géométrie locale des chemins. Ces conséquences restent subordonnées au
noyau formel déjà construit et n’introduisent pas une seconde origine de l’APP.

L’unité suivante construira la subdivision quadruple, le foncteur
cyclotomique et les canaux internes de quatre-vingt-dix et cent vingt positions.
Ce n’est qu’ensuite que l’espace exhaustif des germes sera introduit.
